\documentclass[
ngerman,
paper=a4,
asymmetric,
fontsize=9pt,
% headings=small,
draft=false,
parskip=half,
% listof=totoc,         % verschiebt Minitoc/Hyperref
% bibliography=totoc,   % verschiebt Minitoc/Hyperref
bibtotocnumbered,
% index=totoc,
% noonelinecaption,
footnotes=multiple,
toc=flat,
]{scrbook}
% \setlength{\parindent}{1.5em}
% \usepackage{scrhack}

% UTF-8 Schriftkodierung
\usepackage[utf8]{inputenc}\usepackage[T1]{fontenc}
% get rid of the "no room for new \ldots". Standard TeX registers are limitied to 256, etex increases the max.
\usepackage{etex}

\pdfcompresslevel=9
\pdfimageresolution=72
% \pdfpkresolution=72
% \usepackage[a-1b]{pdfx}
% \def\Keywords{AASS} \def\Title{AASS} \def\Author{ARGE AASS} \def\Org{ARGE AASS}
% % \input glyphtounicode.tex \input glyphtounicode-cmr.tex \pdfgentounicode=1
% \usepackage[xcolor]{rvdtx}


\usepackage{blindtext}
\usepackage{microtype,fixltx2e,ellipsis}

% Optionen für Bilder:
% [draft] .. Es werden nur Rahmen angezeigt
% [final] .. Es werden Bilder normal angezeigt
% [demo} .. keine Bilder
\usepackage[draft]{graphicx} % Bilder-Support

% \usepackage{showidx} % shows index entries on page margin
% \usepackage{showkeys}

\input{./inputs/input-preamble-ExperimentalLayout}

% Trennungen
%%%%%%%%%%%%%%%%%%%%%%%%%%%%%%%%%%%%%%%%%%%%%%%%%%%%%%%%%%%%%%%%   Hyphenation

\hyphenation{
AASS
Ab-teil-ung
Ab-teil-ung-en
Alarm-zeichen
aller-dings
Anam-nese
An-eurys-ma
An-eurys-mas
An-eurys-men
An-forder-ung-en
Arbeits-ge-mein-schaft
ARGE
Auf-klär-ung
Aus-bild-ungs-stand-ards
Aus-bildungs-ver-ordnung
Aus-bildungs-zentrum
Aus-wurf-leist-ung
Bauch-aorten-an-eurys-ma
Bauch-aorten-an-eurys-mas
Bauch-aorten-an-eurys-men
Be-at-mung
Be-hand-lungs-mö-glich-keiten
Be-weg-lich-keit
be-wusst-seins-klar
Be-zeich-nung
be-droh-ten
be-reits
Carotis-sono-gra-phie
Des-in-fekt-ion
Dienst-nehmer-haft-pflicht-ge-setz
Donau-stadt
ein-rast-en
ent-weder
Ent-wick-lung
er-ken-nen
Fett-em-bolie
Fing-er-end-glied
Fing-er-end-glied-er
Florids-dorf
Fremd-an-am-nese
Füssig-keits-de-fi-zit
Füssig-keits-de-fi-zits
Füssig-keits-de-fi-zit-es
Ge-lenks-ver-letz-ung
Ge-lenks-ver-letz-ung-en
gleich-ge-stellt
gleich-ge-stellt-en
gleich-zeit-ig
gleich-zeit-ige
gleich-zeit-ig-es
Halb-seiten-läh-mung
Halb-seiten-läh-mung-en
Halb-seiten-schwäche
Haupt-be-schwer-de
Haupt-be-schwer-den
Herz-in-farkt
Herz-in-farkt-es
Herz-in-farkts
hier
hypo-vol-ämi-schen
im-mun-sup-res-sion
im-mun-sup-res-siv
im-mun-sup-res-sive
In-ten-si-tät
Keim-über-träger
Keim-über-träger-in
Keim-über-trägers
Koron-ar-ge-fäße
Körper-kapil-la-ren
Körper-kreis-lauf
Körper-ober-fläche
Kranken-an-stalt
Kranken-an-stalt-en
Kranken-an-stalt-en-ver-bund
Kranken-an-stalt-en-ver-bund-es
Kranken-an-stalt-en-ver-bunds
Kranken-haus
Kranken-haus-es
Kranken-häuser
Kranken-trans-port
Kranken-trans-porte
Kranken-trans-portes
Kranken-trans-ports
Kranken-trans-port-dienst
Kranken-trans-port-dienste
Kranken-trans-port-dienst-es
Kranken-trans-port-wagen
Kranken-trans-port-wag-ens
Kur-an-stalt-en-ge-setz
leicht
Lungen-er-krank-ung
Lungen-er-krank-ung-en
Lungen-kapil-la-re
Lungen-kapil-la-ren
Lungen-kreis-lauf
man-geln-des
Mega-code
Mega-code-train-ing
Mega-code-train-ings
Mehr-höhl-en-ver-letz-ung
Mehr-höhl-en-ver-letz-ung-en
nach-ge-wies-en
nach-ge-wies-en-en
nach-ge-wies-en-es
neuro-lo-gisch
neuro-lo-gische
neuro-lo-gisch-en
neuro-lo-gisch-es
Nie-der-halte-sieb
Nie-der-halte-siebe
Nie-der-halte-siebes
Nie-der-halte-siebs
Not-arzt
Not-arzt-ein-satz-fahr-zeug
Not-arzt-ein-satz-fahr-zeugs
Not-ärzt-in
Not-arzt-hub-schrau-ber
Not-arzt-hub-schrau-bers
Not-arzt-wa-gen
Not-arzt-wa-g-ens
Patient-en-ver-füg-ungs-ge-setz
Per-for-at-ion
Pneu-mo-kok-ken
prä-klin-isch
Puls-oxy-me-ter
Puls-oxy-me-ters
Rea-nim-at-ion
Rea-nim-at-ions-be-reit-schaft
rea-nim-at-ions-pflicht-ig
Rechts-pfle-ge
Rettungs-dienst
Rettungs-trans-port-wagen
Rettungs-trans-port-wag-ens
Sani-tä-ter
Sani-täts-dienst
Schluck-test
schmerz-los-en
Schutz-wirk-ung
sit-uat-ions-ge-recht
sit-uat-ions-ge-rechte
Span-nungs-pneu-mo-thor-ax
Staph-ylo-kok-ken
Symp-tom
symp-tom-los
Tachy-kar-die
Tage
Taschen-atlas
Ter-min-ver-ein-bar-ung
thera-peut-isch
thera-peut-ische
thera-peut-ischen
Tot-raum
un-be-friedi-gend
un-be-friedi-gend-en
Ver-dau-ungs-trakt
Ver-dau-ungs-trakts
Ver-dau-ungs-trakt-es
Ver-nichtungs-schmerz
Vor-hof-flim-mern
Wien
Wien-er
zu-reich-en
}

% Titel
%%%%%%%%%%%%%%%%%%%%%%%%%%%%%%%%%%%%%%%%%%%%%%%%%%%%%%%%%%%
%%% Title

\titlehead{{\DocVersionTyp} {\DocVersion} (Var. ASB\,921 --
{\DocVersionInfo}) \par\DocGitCommitTiny}

\author{{\normalsize Herausgeber} \vspace{0.5em}\\Gabriel, Koch, Emhofer, Motal, Steuer
\vspace{1em}\\{\itshape\normalsize{ Mit Beiträgen von
Christof Koller, Michael Withofner u.a.}}}

\title{Arbeits- und Ausbildungsstandards für den Sanitätsdienst}
\subtitle{Angepasste Variante für das Ausbildungszentrum des ASB
Floridsdorf-Donaustadt}
\subject{\Huge{\AASS}}
\date{\today}
% \date{November 2009}
\publishers{{\fontfamily{lmss}\selectfont ARGE \AASS}}


% \extratitle{{\scalefont{10}\AASS}}

%%% uppertitleback
\uppertitleback{
	\emph{Covergestaltung: Sebastian Gabriel.}

    \emph{Titelbild: Nicht immer ist die Richtung für den qualifizierten
    Krankentransport und den Rettungsdienst so klar vorgegeben wie hier. Meist
    liegt es am Einzelnen, die richtigen Entscheidungen zu treffen, basierend
    auf Wissen und Erfahrung. Die {\AASS} versuchen, ihm in diesem
    Dschungel aus Möglichkeiten und Verantwortung Halt zu geben und ihn zu
    leiten.}
    
    \emph{Bild: GABS.}}
    
    \lowertitleback{\begin{gWideParEnv}\centering\begin{minipage}{8cm}
    \begin{center} 
        Dieses Projekt enstand in Zusammenarbeit mit dem
        
        \emph{\bfseries Ausbildungszentrum des ASB Floridsdorf-Donaustadt} 
        
        und wurde maßgeblich von der
        
        \emph{\bfseries Cyberservice Internetdienstleistungsgesellschaft~m.b.H.}
        
        und
        
        \emph{\bfseries Grissu.org} 
        
        unterstützt.\vspace{2em}
    
        \includegraphics[width=3cm]{./icons/asb.png}
        
        {\sffamily\fontfamily{phv}Gruppe Floridsdorf-Donaustadt}\par\vspace{2em}
        
        \includegraphics[width=5cm]{./images/logos/cyberservice.png}
    
        {\sffamily\fontfamily{phv}www.cyberservice.net}\par\vspace{2em}
    
        {\Huge Grissu.org}\par\vspace{1em}
        \textit{serving the community} \par \textit{since 1999} 
    \end{center}
    \end{minipage}
    \vspace{2em}
    \begin{center}
        {\bfseries\itshape Nur für den internen Gebrauch gem §\,42 Abs.\,6
        UrheberG bestimmt. Weitergabe nicht gestattet.}
    \end{center}
    
    \end{gWideParEnv}}

\dedication{

\bigskip

saluti et solatio aegrorum

\bigskip

\bigskip

\bigskip

\bigskip

\bigskip

\bigskip

\bigskip

\bigskip

\bigskip

\bigskip

\dictum[Andrew S. Tanenbaum (übersetzt)]{Das schöne an Standards
ist, es gibt so viele, aus denen man auswählen kann.}}





% Dokumenteinfos (Status, Version, ...)
% \RequirePackage[iso]{isodateo}
\RequirePackage{datetime}
\renewcommand{\dateseparator}{-}
\newdateformat{iso}{\THEYEAR-\twodigit{\THEMONTH}-\twodigit{\THEDAY}}
\newtimeformat{nummeric}{\twodigit{\THEHOUR}\twodigit{\THEMINUTE}}
\settimeformat{nummeric}


\newcommand{\DocVersion}{0.17.x-{\pdfdate}}
% \renewcommand{\DocVersion}{0.17.20110802-02}
\newcommand{\DocVersionTyp}{Entwicklungsversion}
% \newcommand{\DocVersionTyp}{Vorabversion}
\newcommand{\DocVersionInfo}{Entwicklungsversion -- Für den internen Gebrauch.}
% \newcommand{\DocVersionInfo}{Vorabversion -- Für den internen Gebrauch.} 
\hypersetup{
draft,
colorlinks=false,
linkcolor=darkBlue,
citecolor=darkBlue,
filecolor=darkBlue,
pagecolor=darkBlue,
urlcolor=darkBlue,
% Rahmen um Links:
pdfborderstyle={/S/U/W 0},  % Unterstreichung
% pdfborderstyle={/S/U/W 0.5},  % Unterstreichung
pdfborder={0 0 0},          % Dicke
% linkbordercolor=DefaultB,     % Farbe
% urlbordercolor=DefaultB,     % Farbe
pdftitle={Arbeits- und Ausbildungsstandards für den Sanitätsdienst \DocVersion},
% Sets the document information Title field pdfauthor=ARGE AASS, % Sets the document information Author field
pdfsubject={AASS \DocVersion}, % Sets the document information Subject field
pdfcreator=ARGE AASS, % Sets the document information Creator field
pdfproducer=ARGE AASS, % Sets the document information Producer field
pdfkeywords=ARGE AASS, % Sets the document information Keywords field
% pdfpagelayout=TwoPageRight,
}
% Staus Möglichkeiten:
% Final 				Endversion
% DraftLight			+ Bilder, - Kommentare
% DraftMedium		- Bilder, - Kommentare
% DraftFull			- Bilder, + Kommenatre
% DraftSuper			- Bilder, + Kommenatare, + draft=yes
\newcommand{\DocStatus}{DraftFull}
\newcommand{\DocGitCommit}{Unofficial build, no SHA1 included!}
% \renewcommand{\DocGitCommit}{\input{../tmp/git-lastcommitsha1.tmp}}
\newcommand{\DocGitCommitFormatted}{\ttfamily\DocGitCommit}
\newcommand{\DocGitCommitTiny}{\tiny\DocGitCommitFormatted}


% Kopf- und Fusszeilen
\newcommand{\Fancyheadreset}{
  \fancyheadoffset[RE,RO]{\marginparsep+\marginparwidth}
  \fancyfootoffset[RE,RO]{\marginparsep+\marginparwidth}
  }


\pagestyle{fancy}
\fancyhf{} % clear all header and footer fields
\fancyheadoffset[RE,RO]{\marginparsep+\marginparwidth}
\fancyfootoffset[RE,RO]{\marginparsep+\marginparwidth}

\fancyhead[LE,RO]{\thepage}

\fancyhead[RE]{\EE{\AuthNotes{ | {\today} | ENTWURF |}}}
\fancyhead[LO]{\EE{\AuthNotes{ | ENTWURF | {\today} |}}}

\fancyhead[LE]{{\sffamily{\color{ubuntured}\thepage} | \leftmark}}
\fancyhead[RO]{{\sffamily{\rightmark} | {\color{ubuntured}\thepage}}}

%%% IMMER %%%
\fancyfoot[RO]{{\noindent\AASS{}~ \itshape v.\,\DocVersion}}
\fancyfoot[LE]{{\noindent\AASS{}~ \itshape v.\,\DocVersion}}
%%% DRAFT %%%
% \fancyfoot[CE,CO]{{\noindent\color{red}\ttfamily\small
% \upshape{\DocVersionTyp} | Nur für den internen
% Gebrauch!}}

% 
% % {\tiny
% % Zweig: }\input{./tmp/git-current-branch.tmp} -- {\tiny
% % Letztes Commit: }\input{./tmp/git-lastcommit.tmp}
% }}
%
% %%% FINAL %%%
\fancyfoot[CE,CO]{{{\tiny {\DocVersionTyp} | Nur für den internen
Gebrauch!}}}

\fancyfoot[RE]{{\noindent\DocGitCommitTiny}}
\fancyfoot[LO]{{\noindent\DocGitCommitTiny}}
% % %



% Redefining pagestyle plain
\fancypagestyle{plain}{%
\fancyhf{} % clear all header and footer fields
%%% IMMER %%%
\fancyfoot[RO]{{\noindent\AASS{}~ \itshape v.\DocVersion}}
\fancyfoot[LE]{{\noindent\AASS{}~ \itshape v.\,\DocVersion}}
%%% DRAFT %%%
\fancyfoot[CE,CO]{{\noindent\color{red}\ttfamily\small
\upshape{\DocVersionTyp} -- \today{}  -- Nur für den internen
Gebrauch!}}
% 
% % {\tiny
% % Zweig: }\input{./tmp/git-current-branch.tmp} -- {\tiny
% % Letztes Commit: }\input{./tmp/git-lastcommit.tmp}
% }}
%
% %%% FINAL %%%
% \fancyfoot[CE,CO]{\EE{{\itshape {\DocVersionTyp} -- Nur für den internen Gebrauch!}}}
\fancyfoot[RE]{{\noindent\DocGitCommitTiny}}
\fancyfoot[LO]{{\noindent\DocGitCommitTiny}}
\renewcommand{\headrulewidth}{0pt}
\renewcommand{\footrulewidth}{0pt}
}


% Ausgeschlossene Kapitel
\usepackage{excludeonly}
% %\excludeonly{include-betriebsintern-chapter,include-regional-chapter,include-taetigkeitsbild-chapter,include-erstehilfe-chapter,include-recht-chapter,include-reserve-ende}
\excludeonly{
./chapters/chapter-NEW,
./chapters/chapter-ARZ,
./chapters/chapter-WIS,
% ./chapters/chapter-INT,
}

% Schliesse nur bestimmte Kapitel ein ...
% \includeonly{%
% ./chapters/chapter-ABD,
% ./chapters/chapter-AED,
% ./chapters/chapter-ALS,
% ./chapters/chapter-ANA,
% ./chapters/chapter-ARZ,
% ./chapters/chapter-ASB,
% ./chapters/chapter-ASS,
% ./chapters/chapter-AUP,
% ./chapters/chapter-BAU,
% ./chapters/chapter-BET,
% ./chapters/chapter-BLS,
% ./chapters/chapter-DOK,
% ./chapters/chapter-EHI,
% ./chapters/chapter-EKG,
% ./chapters/chapter-EIN,
% ./chapters/chapter-GEF,
% ./chapters/chapter-GLO,
% ./chapters/chapter-GYN,
% ./chapters/chapter-HKL,
% ./chapters/chapter-HYG,
% ./chapters/chapter-INK,
% ./chapters/chapter-INT,
% ./chapters/chapter-JUF,
% ./chapters/chapter-JUS,
% ./chapters/chapter-KHD,
% ./chapters/chapter-KIN,
% ./chapters/chapter-KRI,
% ./chapters/chapter-LER,
% ./chapters/chapter-MP1,
% ./chapters/chapter-NEU,
% ./chapters/chapter-NEW,
% ./chapters/chapter-NKA,
% ./chapters/chapter-NWG,
% ./chapters/chapter-ADVANCED-Atemweg,
% ./chapters/chapter-PAM,
% ./chapters/chapter-PKU,
% ./chapters/chapter-PSI,
% ./chapters/chapter-PSY,
% ./chapters/chapter-REG,
% ./chapters/chapter-RSP,
% ./chapters/chapter-RWT,
% ./chapters/chapter-SAN,
% ./chapters/chapter-SHF,
% ./chapters/chapter-SPE,
% ./chapters/chapter-STW,
% ./chapters/chapter-SYM,
% ./chapters/chapter-TAE,
% ./chapters/chapter-THE,
% ./chapters/chapter-TOD,
% ./chapters/chapter-TOX,
% ./chapters/chapter-TRA,
% ./chapters/chapter-URO,
% ./chapters/chapter-VIT,
% ./chapters/chapter-WEI,
% ./chapters/chapter-WIS,
% ./chapters/chapter-WUN,
% inc-APP,
% inc-RES-ende,
% ./chapters/part-Lehrmaterialien,
% }

% optional --- Optionaler Text. 
%%%%%%%%%%%%%%%%%%%%%%%%%%%%%%%%%%%%%%%%%%%%%%%%%%%%%%%%%%%%%%%%%%%%%%%%%%%%%%%
%%%%%%%%% OPTIONEN
%%%%%%%%%%%%%%%%%%%%%%%%%%%%%%%%%%%%%%%%%%%%%%%%%%%%%%%%%%%%%%%%%%%%%%%%%%%%%%%
% optional --- Optionaler Text. 
\usepackage
[
%%% Status
% STATUS-DRAFT,
% STATUS-DRAFT-SHOWINDEX,
% STATUS-DRAFT-SHOWKEYS,
% STATUS-DRAFT-WITHPICTURES,
STATUS-WITH-BIBLIO,
% STATUS-FINAL,
% STATUS-EXPERIMENTAL,
%%% Organisation
ORGASBFLO,
ORG-ASBFLO,
%%% Varianten
VAR-STANDARD,
VAR-ERNAEHRUNGSPAED,
%%% Niveau
NIVRSA, % Rettungssanitäter .at
NIVNFS, % Notfallsanitäter .at
% NIVDOC, % Arzt, Mediziner, Med.-Student
% NIVEHT, % EH-Trainer
%%% Region
REGAUT,
% REGGER;
REGVIE,
% REGNOE,
%%% Medizinprodukte
%%% MPG-Hersteller-Baureihe-Typ
%% Tragem
MPG-FERNO-TRAGE, % 
MPG-STOLLENWERK, % 
%% Monitore / Defis 
MP-MEDTRONIC-LP, % 
MP-MEDTRONIC-LP-12, % 
MP-MEDTRONIC-LP-15, % 
MP-MEDTRONIC-LP-500, %
MP-MEDTRONIC-LP-1000, %  
MP-SCHILLER-DEFIGARD, % 
MP-SCHILLER-DEFIGARD-1002, % 
MP-SCHILLER-DEFIGARD-2002, %  
MP-SCHILLER-FRED, %  
%% Beatmungsgeräte
MP-DRAEGER-OXYLOG-1000, % 
MP-DRAEGER-OXYLOG-2000, % 
MP-DRAEGER-OXYLOG-3000, % 
MP-WEINMANN-MEDUMAT-COMPACT, %
MP-WEINMANN-MEDUMAT-STANDARD, %
MP-WEINMANN-MEDUMAT-STANDARDA, %
MP-WEINMANN-MEDUMAT-TRANSPORT %     
]{optional}

\setcounter{errorcontextlines}{10}

% \usepackage{fontspec}
% \usepackage{setspace}


\usepackage{ragged2e}
% \usepackage[xcolor]{changebar}
% % \driver{PDFTeX} % changebar

% Randnummernzähler einrichten 
\newcounter{randnr} 
\newcommand{\newrandnummer}{\refstepcounter{randnr}} 
\newcommand{\rdnr}{\newrandnummer\marginnote{{\sffamily\bfseries\scriptsize\therandnr}}} 

\makepagenote

% \newcommand{\MultiColMassnahmenStart}[1]{\addtocontents{lomassnahmen}{\protect\begin{multicols}{#1}}}
% \newcommand{\MultiColMassnahmenEnd}{\addtocontents{lomassnahmen}{\protect\end{multicols}}}
% \newcommand{\MultiColFazitStart}[1]{\addtocontents{lofazit}{\protect\begin{multicols}{#1}}}
% \newcommand{\MultiColFazitEnd}{\addtocontents{lofazit}{\protect\end{multicols}}}
  
\newcommand{\MultiColListStart}[2]{\addtocontents{#1}{\protect\begin{multicols}{#2}}}
\newcommand{\MultiColListEnd}[1]{\addtocontents{#1}{\protect\end{multicols}}}
\newcommand{\Sign}[1]{{\sffamily#1}}


\newcommand{\KAMEL}[1]{{\color{DefaultB}\small\textcircled{\ttfamily\bfseries\scriptsize#1}}}
\newcommand{\TxKAMEL}{(S-)KAMEL}
\newcommand{\SpeechMargin}[1]{\marginpar{{\small\Speech{#1}}}}
\newcommand{\SpeechTextMargin}[1]{\marginpar{{\small\Speech{#1}}}\Speech{#1}}
\newcommand{\Study}[1]{\textit{#1}}
\newcommand{\FontTypewriter}{\fontfamily{lmtt}\fontseries{lc}\selectfont}


\newenvironment{PictureSeries}[2]
{%
\begin{figure}[H]
% \ifthenelse{\equal{#1}{Wide}}{\begin{WidePar}}{}
\Captionof{figure}{#2}\par}
{%
% \ifthenelse{\equal{#1}{Wide}}{\end{WidePar}}{}
\end{figure}
}

\newenvironment{PictureSeriesWide}[2]
{%
\begin{figure}[H]
\begin{WidePar}
\Captionof{figure}{#2}\par}
{%
\end{WidePar}
\end{figure}
}

\newcommand{\PictureSeriesRow}[9]{
\ifthenelse{\equal{#1}{2}}
{\subfloat
[#3]
{\ifthenelse{\equal{#2}{empty}}{\parbox{0.32\linewidth}{~}}{\includegraphics[width=0.45\linewidth]{#2}}}
\hfill
\subfloat[#5]
{\ifthenelse{\equal{#4}{empty}}{\parbox{0.32\linewidth}{~}}{\includegraphics[width=0.45\linewidth]{#4}}}
}
{}
\ifthenelse{\equal{#1}{3}}
{\subfloat
[#3]
{\ifthenelse{\equal{#2}{empty}}{\parbox{0.32\linewidth}{~}}{\includegraphics[width=0.32\linewidth]{#2}}}
\hfill
\subfloat[#5]
{\ifthenelse{\equal{#4}{empty}}{\parbox{0.32\linewidth}{~}}{\includegraphics[width=0.32\linewidth]{#4}}}
\hfill
\subfloat[#7]
{\ifthenelse{\equal{#6}{empty}}{\parbox{0.32\linewidth}{~}}{\includegraphics[width=0.32\linewidth]{#6}}}
}
{}
\ifthenelse{\equal{#1}{4}}
{\subfloat
[#3]
{\ifthenelse{\equal{#2}{empty}}{\parbox{0.24\linewidth}{~}}{\includegraphics[width=0.24\linewidth]{#2}}}
\hfill
\subfloat[#5]
{\ifthenelse{\equal{#4}{empty}}{\parbox{0.24\linewidth}{~}}{\includegraphics[width=0.24\linewidth]{#4}}}
\hfill
\subfloat[#7]
{\ifthenelse{\equal{#6}{empty}}{\parbox{0.24\linewidth}{~}}{\includegraphics[width=0.24\linewidth]{#6}}}
\hfill
\subfloat[#9]
{\ifthenelse{\equal{#8}{empty}}{\parbox{0.24\linewidth}{~}}{\includegraphics[width=0.24\linewidth]{#8}}}
}
{}
}

% \usepackage{colortbl}
% \arrayrulecolor{Beige}

\renewcommand{\Abk}[1]{#1}


% #######################################################################################################
% #######################################################################################################
% #######################################################################################################
% #######################################################################################################
% #######################################################################################################
% #######################################################################################################
% #######################################################################################################
% BEGIN DOCUMENT ########################################################################################
% #######################################################################################################
\begin{document}
% \pagestyle{scrheadings}

% \setmainfont[SmallCapsFont={* Caps},
% Mapping=tex-text]{Latin Modern Roman}
% \setsansfont[Mapping=tex-text]{Latin Modern Sans}
% \setmonofont[SmallCapsFont={* Caps}]{Latin Modern Mono}

% \defaultfontfeatures{Numbers=OldStyle}
\MultiColListStart{toc}{3}
\MultiColListStart{lofazit}{2}
\MultiColListStart{lomassnahmen}{2}



\layout
% \pagediagram\pagedesign\paragraphdiagram\paragraphdesign\floatpagediagram\floatpagedesign\listdiagram\listdesign\footnotediagram\footnotedesign\tocdiagram\tocdesign

\clearpage
\section{Vorwort}

\begin{WidePar}
\begin{multicols}{3}
  Das vorliegende Werk ist das (Zwischen-)Ergebnis einer seit mehr als zwei Jahren
andauernden Arbeit. Oft werden wir gefragt \Speech{"`Wann wird Euer Buch fertig
sein?"'} Nun, das wird es wohl nie sein. Die {\AASS} sind darauf ausgelegt,
ständig weiterentwickelt, gewartet und verbessert zu werden. Das medizinische
und organisatorische Umfeld wandelt sich ständig, sodass man es sich gar nicht
leisten könnte, still zu stehen.

Ziel dieses Werkes ist es, das Wissen und die Welt des Sanitätsdienstes abseits
bereits breitgetretener Pfade darzustellen, und sowohl eine Unterlage für die
Ausbildung, als auch ein Hilfsmittel für die Praxis zu schaffen.  Besonderen
Wert haben wir dabei nicht nur auf die Praxis, sondern auch -- heutzutage nicht
mehr ganz modern -- auf das Hintergrundwissen und relevante Randbereiche gelegt.
Denn in einer Zeit, in der das praktische Wissen innerhalb weniger Jahre bereits
wieder veraltet ist, bildet dieses Fundament die Chance für eine
langjährige Tätigkeit und erlaubt selbstständiges, regel\E{geleitetes} Arbeiten
(anstatt strikten, regel\E{gebundenen} Handelns). 

In diesem Sinne entstanden in Zusammenarbeit mit dem Ausbildungszentrum des ASB
Florisdorf-Donaustadt ausgehend von der Ausbildungsverordnung 
österreichischer Rettungssanitäter und den Erfordernissen der entsprechenden
Kurse die {\AASS}. Sie bilden die Grundlage der Curricula und konnten sich
schon, etappenweise und unter kontrollierten Bedingungen eingeführt, seit über
einem Jahr als Lern- und Ausbildungsunterlage
bewähren.
\end{multicols}
\end{WidePar}



\begin{multicols}{2}
  Eines muss vorweg klargestellt werden: Auch wenn es sehr in
Mode gekommen ist, Abläufe zu standardisieren und
festzuschreiben, so ist dies keine taugliche
Herangehensweise, um komplexe Probleme in der Medizin oder
im Sanitätsdienst zu bewältigen. Zu komplex sind die
Variablen, und zu gering bzw. zu unangepasst an konkrete
Situationen ist das gesicherte Wissen, auf dem
festgeschriebene Prozesse basieren. Der menschliche Faktor,
insbesonders die jeweilige Erfahrungen, ist und bleibt der
entscheidende Faktor jeglichen Handelns. Kein Lehrbuch oder
standardisierter Prozess kann das ändern, sofern das Ziel
ein "`gutes"' Ergebnis, und keine statistische Kennziffer
ist.
\end{multicols}

\paragraph{Dieses Werk ist keine "`Dienstanweisung"'!}~~



In diesem Sinne ist der Titel dieses Werkes nicht im Sinne eines
standardisierten Prozessmanagements zu verstehen, sondern soll einen
\E{Basisstandard} bezüglich medizinisch-sanitätsdienstlichem Wissen und
Kompetenz bilden, der individuell -- je nach Situation bzw. durch vom Helfer
andernorts erworbenen Wissens -- grundsätzlich verbesserbar ist. Die {\AASS}
geben Empfehlungen, es obliegt dem Helfer diese gemäß seiner 
fachlichen Erfahrung und Kenntnisse zu interpretieren, einzusetzen oder
anzupassen -- 

\Speech{salus aegroti suprema lex}\footnote{salus aegroti suprema lex \Lang{lat.}: \Speech{Das
Wohl des Kranken ist oberstes Gebot.}}.

\bigskip

\bigskip

\emph{Wien, im Sommer 2010}



\clearpage
\section{Offenlegung, Impressum, Nutzungsbedingungen und Hinweise}

\begin{WidePar}

\begin{multicols}{3}
\minisec{Medieninhaber}

{\noindent\sffamily Arbeitsgemeinschaft Arbeits- und Ausbildungsstandards für
den Sanitätsdienst -- }{\ARGEAASS}
\\
\noindent{}Engerthstraße 146\,/\,8\,/\,13 \\
A-1200 Wien -- AUSTRIA

\begin{compactdesc}
  \item[E-Mail:] office@aass.at
  \item[Homepage:] \url{http://www.aass.at}
  \item[ZVR:] 846678982
\end{compactdesc}


\minisec{Weitere Informationen}

\begin{compactdesc}
  \item[Verlags- und Herstellungsort:] Wien
  \item[Publikationsart:] Elektronisch. Die etwaige Herstellung von
  Druckwerken erfolgt durch den autorisierten Nutzer.
  \item[ISBN:] ---
  \item[ISSN:] ---
  \item[Status:] \DocVersionTyp{} (\DocVersionInfo{})
  \item[Versionsgeschichte:]~
  
  {\scriptsize %
  \begin{description}
    \item[0.2.0] 2009-07 (VV\footnote{Vorabversion}, 349 Änderungen\footnote{Seit
    Übernahme in das Versionsverwaltungssystem.})
    \item[0.4.0] 2009-09 (VV, 284 Änderungen)
    \item[0.6.0] 2009-10 (VV, 322 Änderungen)
    \item[0.8.0] 2010-01 (VV, 653 Änderungen)
    \item[0.10.0] 2010-06 (VV, 935 Änderungen)
    \item[0.12.0] 2010-08 (VV, 489 Änderungen)
    \item[0.14.0] 2011-01 (VV, 706 Änderungen)
    \item[0.16.0] 2011-03 (VV, 141 Änderungen)
    \item[0.18.0] 2011-08 (VV, 449 Änderungen)
    \item[0.20.0] 2011-?? (VV, ??? Änderungen\footnote{Lt.
    vorläufiger Statistik, Änderungen nach Redaktionsschluss möglich.})
  \end{description}
  }
\end{compactdesc}

\noindent{}Die {\ARGEAASS} ist ein gemeinnütziger Verein mit Sitz in Wien,
welcher die Entwicklung, Wartung, Qualitätssicherung, Anpassung, Verbreitung und Lehre von
wissenschaftlich und sachlich fundierten Arbeits- und Ausbildungsstandards für
den Sanitätsdienst (AASS), sowie die wissenschaftliche Forschungstätigkeit in
diesem Gebiet bezweckt.

\minisec{Mitarbeit und Meldung von Fehlern}

\noindent{}Wünsche, Anregungen, Beschwerden und Fehlermeldungen nehmen wir
gerne unter 
% \url{http://www.aass.at/bugs}
\href{mailto:anregungen@aass.at}{\nolinkurl{anregungen@aass.at}} entgegen.

\minisec{Nutzungsbestimmungen}

\noindent{}Dieses Werk ist nur für den internen Gebrauch gem. §\,42
Abs.\,6 UrheberG bestimmt. Die Weitergabe ist nicht
gestattet. 

Sollten Sie Interesse an den Arbeits- und
Ausbildungsstandards für den Sanitätsdienst haben, setzen
Sie sich bitte mit uns in Verbindung.

\minisec{Hinweise}

\noindent{}Wie jede Wissenschaft ist die Medizin ständigen Entwicklungen
unterworfen. Soweit fachliche Angaben gemacht werden, darf der Leser
zwar darauf vertrauen, dass die Autoren große Sorgfalt
darauf verwandt haben, dass diese Angaben dem Wissensstand bei
Fertigstellung des Werkes entspricht. Für Angaben, speziell
über Dosierungsanweisungen und Applikationsformen, kann jedoch
keine Gewähr übernommen werden. 

Jeder Leser ist angehalten, durch
sorgfältige Prüfung der Literatur, der Lehrmeinung so\-wie der
Beipackzettel der verwendeten Präparate und gegebenenfalls nach
Konsultation eines Spezialisten festzustellen, ob die gegebene
Aussage oder Empfehlung für Dosierungen oder die Beachtung von
Kontraindikationen gegenüber der Angabe in diesem
Werk abweicht. Jede Dosierung oder Applikation erfolgt
auf eigene Gefahr des Benutzers. 

Die {\ARGEAASS} bittet jeden Leser, ihm aufgefallene Ungenauigkeiten und
Fehler, sowie Anregungen und konstruktive Kritik, über die Email-Adresse
\href{mailto:anregungen@aass.at}{\nolinkurl{anregungen@aass.at}} mitzuteilen.

Geschützte Warennamen und Warenzeichen werden nicht durchgehend besonders
kenntlich gemacht. Aus dem Fehlen eines solchen Hinweises kann also
nicht geschlossen werden, dass es sich um einen freien Warennamen
handelt.

\minisec{Die Technik}

\begin{compactdesc}
  \item[Satz und Druckvorstufe] \LaTeX{} mit {\sffamily KOMA-Script}
  
  Distribution: \TeX live 2009, u.\,a. mit den Paketen: AASS
  Style Collection (asc),
  \Code{microtype},
  \Code{babel},
  \Code{babelbib} (mehrsprachiges Literaturverzeichnis),
  \Code{shorttoc} (Inhaltsübersicht),
  \Code{titeltoc},
  \Code{makeidx} (Index),
  \Code{prettyref},
  \Code{units},
  \Code{nicefrac},
  \Code{soul},
  \Code{soulutf8},
  \Code{booktabs},
  \Code{marvosym} (Symbole),
  \Code{txfonts} (Symbole),
  \Code{skul} (Symbole),
  \Code{ifsym} (Symbole),
  \Code{dingbat} (Symbole),
  \Code{paralist},
  \Code{geometry} (Seitenlayout),
  \Code{fancybox},
  \Code{lmodern} (Latin Modern Fonts),
  \Code{titlesec} (Sectioning titles),
%   \Code{float},
  \Code{hyperref} (Hypertext links),
  \Code{url} (URL's),
  \Code{quotchap},
  \Code{minitoc} (Inhaltsverzeichnisse für Kapitel),
  \Code{chappg} (Seitennummerierung pro Kapitel),
  \Code{caption} (Customizing captions),
  \Code{caption3},
  \Code{mhchem} (chemische Formeln),
  \Code{optional} (optionaler Text),
  \Code{ulem}
  \item[Entwicklungsumgebung] Eclipse Ganymed mit Texlipse- und JGit/EGit-Plugin,
  Texmaker
  \item[Versionsverwaltung] Git
  \item[Literaturverzeichnis] JabRef 2.5.0, BibTeX, Babelbib
  \item[Textverarbeitung und Rohentwurf] OpenOffice
  \item[Betriebssysteme] \E{Rohentwurf}: Linux (Debian, Gentoo, Ubuntu), MacOS X,
  MS Windows XP, Vista
  
  \E{Endproduktion}: Ubuntu Linux 10.10 \textit{"`Maverick
  Meerkat"'}, \url{http://www.ubuntu.com}
\end{compactdesc}


\end{multicols}
\end{WidePar}



\clearpage






% \begin{WidePar}
%   \setcounter{tocdepth}{2}\shorttableofcontents{Übersicht}{0}
% \end{WidePar}
% {\loadgeometry{FullPage}\twocolumn\setcounter{tocdepth}{2}\shorttableofcontents{Übersicht}{0}\onecolumn\restoregeometry}

%\renewcommand\contentsname{Inhalt - Schnell\"ubersicht}
\dominitoc[n]

\begin{WidePar}
  {\setcounter{tocdepth}{3}\scriptsize\tableofcontents}
\end{WidePar}




\begin{savequote}
\begin{verbatim}
 _______________________________
( Don't look back, the lemmings )
( are gaining.                  )
 -------------------------------
       o   ,__,
        o  (oo)____
           (__)    )\
              ||--|| *

\end{verbatim}
\end{savequote}

\chapter
[\CO{[TAE]} Tätigkeitsbild und Orientierung]
{\CC{[TAE]}\\Tätigkeitsbild und Orientierung}
\label{chp:TAE}

\minitoc

\section{Präklinik}

\subsection{Sanitäter}

\subsubsection{Sanitäter in der Theorie}

Der Beruf bzw. die Tätigkeit des Sanitäters wurde 2002 durch das neue 
\E{Sanitätergesetz} (SanG) neu geregelt. Demnach unterteilt man in
\E{Rettungssanitäter} oder \E{Notfallsanitäter}. Der Beruf bzw. die Tätigkeiten des Sanitäters dürfen nur nach Maßgabe des Sanitätergesetzes
 ausgeübt werden. 

Sanitäter haben ihre Tätigkeit \E{ohne Ansehen der Person} gewissenhaft
auszuüben. Sie haben das Wohl der Patienten und der betreuten Personen nach
\E{Maßgabe der fachlichen und wissenschaftlichen Erkenntnisse und Erfahrungen}
zu wahren. 
% Nötigenfalls ist ein Notarzt oder, wenn ein solcher nicht zur Verfügung steht,
% ein sonstiger zur selbständigen Berufsausübung berechtigter Arzt anzufordern.

\Es{Fortbildung} ist ein wichtiger Teil der Tätigkeit, Sanitäter haben sich
laufend tätigkeitsrelevant fortzubilden. Darüber hinaus muss er \E{Berufs- bzw.
Tätigkeitspflichten} erfüllen (Dokumentationspflicht, Verschwiegenheitspflicht,
Auskunftspflicht).


\Layout{60}{Allgemeiner Tätigkeitsbereich}{%
Der Tätigkeitsbereich des Sanitäters beinhaltet eine Reihe von Tätigkeiten,
welche \Es{eigenverantwortlich} durchgeführt werden. Dazu gehört die
eigenverantwortliche Anwendung von Maßnahmen der qualifizierten Ersten Hilfe,
Sanitätshilfe und Rettungstechnik, einschließlich diagnostischer und
therapeutischer Verrichtungen.
}
{}
{./images/TAE/dif-w2-zelt-1.jpg}
{Fürsorge für erkrankte oder
hilfsbedürftige Personen: Die Kernaufgabe eines jeden Sanitäters.}
{}


\Layout{60}{Tätigkeitsbereich des Rettungssanitäters}{%
Der Tätigkeitsbereich des \Es{Rettungssanitäters} umfasst die selbständige und
\E{eigenverantwortliche Versorgung und Betreuung kranker, verletzter und
sonstiger hilfsbedürftiger Personen}, die medizinisch indizierter Betreuung
bedürfen, vor und während des Transports, die Übernahme sowie die Übergabe des
Patienten oder der betreuten Person im Zusammenhang mit einem Transport, 
eine qualifizierte Durchführung von lebensrettenden Sofortmaßnahmen sowie  
die sanitätsdienstliche Durchführung von Sondertransporten.
}
{}
{./images/TAE/tae-teamwork-transfer-lbk.jpg}
{Auch die gute Zusammenarbeit mit dem
Spitalspersonal ist wichtig, wie zum Beispiel bei dieser Überstellung eines
Intensivpatienten. Wenn alle Beteiligten danach so fröhlich sind ist das ein
gutes Zeichen.}
{}



\Layout{2}{Tätigkeitsbereich des Notfallsanitäters}{%
Der Tätigkeitsbereich des \Es{Notfallsanitäters} umfasst zusätzlich die
\E{Unterstützung des Arztes} bei allen notfall- und katastrophenmedizinischen
Maßnahmen einschließlich der Betreuung und des sanitätsdienstlichen Transports von
\E{Notfallpatienten}, 
die Verabreichung von dafür freigegebenen
\E{Arzneimitteln}, 
die eigenverantwortliche Betreuung der berufsspezifischen Geräte, Materialien und
Arzneimittel und die Mitarbeit in der \E{Forschung}. Notfallsanitäter können
die Berechtigung zur Durchführung allgemeiner und besonderer
\E{Notfallkompetenzen} erwerben (Arzneimittellehre, Venenzugang und Infusion,
Beatmung und Intubation). 
}{}{}{}{}


\Layout{60}{Dienstverhältnis}{%
Tätigkeiten des Sanitäters dürfen \E{ehrenamtlich}, \E{berufsmäßig}, als
\E{Zivildienstleistender} oder unter anderen bestimmten Umständen\ZFootnote{Die
Tätigkeiten des Sanitäters dürfen ferner als Soldat im Bundesheer, als Organ des
öffentlichen Sicherheitsdienstes, Zollorgan, Strafvollzugsbediensteter,
Angehöriger eines sonstigen Wachkörpers} ausgeübt werden. Die \Es{berufsmäßige}
Ausübung von Tätigkeiten des Sanitäters setzt neben der jeweiligen
Fachausbildung die Absolvierung eines \E{Berufsmoduls} voraus. 
}
{}
{./images/TAE/tae-okt26-alle.jpg}
{Im Notfall ist die
Organisationszugehörigkeit nebensächlich.}
{}


\Layout{2}{Befristung}{%
Die Berufs- und Tätigkeitsberechtigung ist mit jeweils zwei Jahren
\E{befristet}. Zur Verlängerung der Berufs- und Tätigkeitsberechtigung bedarf
es der Absolvierung von Fortbildungen sowie einer Rezertifizierung.
}{}{}{}{}


% (2) Notfallpatienten gemäß Abs. 1 Z 2 sind Patienten, bei denen im Rahmen einer
% akuten Erkrankung, einer Vergiftung oder eines Traumas eine lebensbedrohliche
% Störung einer vitalen Funktion eingetreten ist, einzutreten droht oder nicht
% sicher auszuschließen ist.

\subsubsection{Sanitäter in der Praxis}

Es ist derzeit noch relativ unklar wie der Sanitäter in das etablierte
Berufsfeld der Gesundheitsberufe "`hineinpasst"'. Die Berufsgruppen der Ärzte,
Hebammen, Pflegepersonal usw. gibt es schon seit Jahrzehnten oder gar
Jahrhunderten und es blieb genug Zeit für Streitigkeiten bezüglich der
jeweiligen Kompetenz und der Abgrenzung der Tätigkeitsfelder. Des öfteren hat
der Gesetzgeber Regelungen aufgestellt und die Gerichte Orientierungspunkte
gegeben; darüber hinaus wurden durch die ständige Arbeit miteinander Realitäten
"`geschaffen"'. In manchen Punkten aber auch heutzutage einiges Unklar oder
einem stetigen Wandel unterworfen. 

\Layout{2}{Sanitätergesetz 2002}{%
Mit der Schaffung des Sanitätergesetzes im Jahr 2002 wollte man die Tätigkeit
der Sanitäter endlich klar regeln und Rechtssicherheit schaffen. Gleich vorweg,
wirkliche Klarheit über die Tätigkeiten des Sanitäters wurde nicht geschaffen,
wie unzählige Diskussionen zu diversen Themen beweisen. Bezüglich der genauen
rechtlichen Regelungen sei auf das Kapitel \prettyref{chp:JUS} verwiesen, an
dieser Stelle wollen wir uns aus praktischer Sicht dem Tätigkeitsfeld und den
Aufgaben der Sanitäter widmen. 
}{}{}{}{}

\Layout{2}{Das Miteinander}{%
Es scheint ein Naturgesetz zu sein dass, sobald zwei oder mehr Berufsgruppen
aufeinander treffen, es zu Streitigkeiten bezüglich der jeweiligen Kompetenzen,
Rechte und Befugnisse kommt. Auf der Strecke bleibt oft das Potential welches
entwickelt werden könnte, würden die Berufsgruppen konstruktiv und
zielgerichtet zusammenarbeiten. In diesem Sinne sind die \AASS{} auf ein
interdisziplinäres Miteinander der Berufsgruppen.}{}{}{}{}

\opt{STATUS-DRAFT}{
\Layout{2}{Der Sanitäter -- früher}{%
\TakeNotesLarge
}{}{}{}{}
}

\Layout{2}{Der Sanitäter -- heute}{%
Das Aufgabengebiet des Sanitäter von heute wird immer mehr um Tätigkeiten
erweitert, die medizinisch-wissenschaftliche Kenntnisse voraussetzen
und noch vor einiger Zeit nur den Ärzten oder anderen Berufsgruppen
vorbehalten waren. 

Dies geschah anfangs durch das SanG, später aber auch durch die Anforderungen,
welche sich zwangsläufig aus der Arbeit in der Praxis ergeben und z\,T. gar
nicht gesetzlich geregelt sind. 

In letzter Zeit kommt es auch zu einem Einfluss privater, meist kommerzieller,
internationaler Ausbildungsprogramme wie z.\,B.
\Abk{AMLS}{\texttrademark}\footnote{Advanced Medical Life
Support{\texttrademark}}, \Abk{ITLS}{\texttrademark}\ZFootnote{International
Trauma Life Support{\texttrademark}}, \Abk{PHTLS}{\texttrademark}\ZFootnote{Prehospital Trauma Life Support{\texttrademark}} etc. Deren Lehrinhalte sind meist auf den Kompetenzbereich von
Paramedics\ZFootnote{Paramedics sind gutausgebildete Sanitäter im
englischsprachigen Raum, welche in nicht-Notarzt-basierten Systemen eingesetzt
werden, d.\,h. sie haben in diesen Ländern das Personal mit der höchsten
medizinischen Ausbildung, welches im präklinischen Regelbetrieb eingesetzt
wird.} und reizen Grenzen dessen, was allgemein als Aufgaben der Sanitäter
angesehen wird, weitgehend aus, oder überschreiten diese. Diesbezüglich kommt
es immer wieder zu Diskussionen, was Sanitäter tun bzw.
unterlassen\footnote{Schädliche Dinge zu unterlassen ist genauso wichtig wie
nützliche Dinge zu tun.} \sout{dürfen} \E{müssen}\footnote{Meist kann man es
sich nicht aussuchen, ob man eine Maßnahme ergreift oder nicht. Handlungen,
die man machen oder unterlassen darf und die das Patientenwohl fördern,
\E{müssen} i.\,d.\,R. entsprechend vorgenommen oder unterlassen werden.}, und was sie
\E{nicht} tun oder unterlassen dürfen. }{}{}{}{}

% \clearpage


\subsubsection{Anforderungen -- Beanspruchung -- Belastung: Was erwartet die
Welt von mir?}

\label{topic:belastung}

\Layout{2}{{\TxQuerverweise}}{%
\begin{inparaitem}
  \item Gesprächsführung: \CO{[AUP]} (\ref{chp:AUP});
  \item Belastungsreaktion, Burn-out:
  \prettyref{topic:ptss}, \prettyref{topic:burn-out}
\end{inparaitem}
}{}{}{}{}

\begin{description}
  \item[Anforderungen:] \H{fachliches Wissen --
  praktische Kompetenzen -- dieses in schwierigen Situationen umsetzen zu
  können professionelle Haltung -- social/soft skills zwischenmenschliche
  Fähigkeiten, Frustrationstoleranz -- Umgang mit mehrdeutigen
  Situationen bzw. "`schwierigem Klientel"' -- Teamdifferenzen -- \ldots}
\end{description}


\begin{center}
\ttfamily\small
\begin{minipage}{45em}
\begin{verbatim}
 _________________________________________
( Your object is to save the world, while )
( still leading a pleasant life.          )
 -----------------------------------------
  o
   o   \_\_    _/_/
    o      \__/
           (oo)\_______
           (__)\       )\/\
               ||----w |
               ||     ||

\end{verbatim}
\end{minipage}
\end{center}

\Layout{2}{Theoretische Anforderungen}{%
Die Tätigkeit im Rettungs- und Krankentransportwesen geht aufgrund der
Komplexität der Aufgaben mit einer Fülle von körperlichen, wissensmäßigen und
psychischen Anforderungen einher. Einen nur kargen Hinweis darüber, welche Qualitäten von
einem (Zivildienstleistenden) Sanitäter erwartet werden, liefert das
Anforderungsprofil des AMS, nämlich:



\begin{itemize}
\item Körperliche und psychische Belastbarkeit
\item rasches Auffassungs- und Reaktionsvermögen
\item Einfühlungsvermögen
\item Beobachtungsgabe
\item Kommunikations- und Teamfähigkeit
\item Zuverlässigkeit
\item Verantwortungsbewusstsein
\end{itemize}
}{%

}{}{}{}



Im Folgenden soll der Versuch unternommen
werden diese nüchternen Auflistung von Begrifflichkeiten
plastischer darzustellen und zu erweitern.

\Layout{0}{Anforderungen und Belastungen in der Praxis}{%
Es versteht sich von selbst, dass die professionelle Versorgung erkrankter bzw.
verunfallter Menschen ein medizinisches Wissen voraussetzt und damit auch
einhergehend die praktischen Kompetenzen, dieses \Es{zuverlässig und
eigenverantwortlich} in schwierigen Situationen umzusetzen zu können. 

Dringliche
Versorgungssituationen bzw. Notfälle, in denen jede zeitliche Verzögerung
zulasten des Patienten fallen könnten, machen ein \Es{schnelles Einschätzen der
Lage} so wie die Bereitschaft \Es{Verantwortung zu übernehmen} notwendig. Da es
sich stets um eine Interaktion zwischen zwei Personen -- nämlich Versorgendem
und zu Versorgenden -- handelt, gelten zwischenmenschliche Fähigkeiten wie
höfliche Umgangsformen, \Es{Kommunikationsfähigkeit} sowie \Es{Einfühlungsvermögen} als
Grundvoraussetzung. 

Erkrankung bzw. Unfall stellen für die Betroffenen in
häufigen Fällen \E{kein alltägliches Geschehen} dar, sondern haben oft einen
kritisch-krisenhaften, jedenfalls aber frustrierend-belastenden Charakter.
Deshalb soll an dieser Stelle die besondere Wichtigkeit einer \Es{einfühlenden
und verständnisvollen Grundhaltung} betont werden. Dies kann für den
professionellen Helfer auch damit einhergehen, \E{seine eigenen Bedürfnisse}
bzw. Empfindungen für den Moment \E{hintanzustellen}, also eine hohe
\Es{Frustrationstoleranz}
erforderlich machen. 

\Fazit{Der (richtige) Umgang mit Frustration ist wesentlich, um im Rettungs-
Krankentransportdienst auf Dauer "`zu überleben"'.}

Im Rettungs- und Krankentransportdienst beschäftigte Personen arbeiten
 die meiste Zeit in Teams, sodass Qualitäten wie Toleranz,
Kritikfähigkeit und Selbstdisziplin ("`\Es{Teamfähigkeit}"') ebenso von hoher
Bedeutung sind (Die Praxis zeigt, dass hier häufig viel Spielraum für
Verbesserungen besteht \Code{;-)}).

Weiters sei auf die Wichtigkeit der \Es{\E{körperlichen} Belastungsfähigkeit}
als erforderliches Kriterium hingewiesen: das Tätigkeitsfeld des Rettungssanitäters
umfasst diverse Aufgaben, die mit körperlichem Einsatz und dementsprechender
Anstrengung verbunden sind.

Zu guter Letzt und im besonderen sei auf die Bedeutung \Es{psychischer
Belastbarkeit} hingewiesen, \AuthNotes{die auch in den folgenden Kapiteln zum
Thema gemacht werden soll,} v.\,a. in Hinblick auf den Umgang mit psychischem
Stress: 

Sanitätsdienst zu leisten, bedeutet tägliche \E{Konfrontation} mit dem
persönlichen Leid von Mitmenschen, des weiteren die \E{Verantwortlichkeit} –
ggfs. auch noch unter Zeitdruck – darauf professionell zu reagieren und zusätzlich die
besonderen und generellen \E{Stressauslöser des Arbeitsbereichs} eines
Sanitäters zu meistern. Dies benötigt eine hinreichend funktionierende
psychische Verarbeitungsfähigkeit, um ein gesundes seelisches Gleichgewicht bzw. sein
persönliches Wohlbefinden wahren und pathologischen Entwicklungen (Burnout,
Belastungsreaktion, etc.) präventiv entgegenwirken zu können.
}{%
}{}{}{}
\begin{multicols}{2}

% \Layout{27}{}{
% 
% }
% {}
% {}
% {}
% {}
\begin{figure}[H]
\includegraphics[width=\linewidth]{./images/TAE/dif-nacht-atmosphaere.jpg}
\Captionof{figure}{\AuthNotesPicLegalOk{}{}Hitze, Dunkelheit und
kontrolliertes Chaos: Ambulanzdienst auf dem Donauinselfest.}
\end{figure}

% \Layout{27}{}{
% 
% }
% {}
% {}
% {}
% {}
\begin{figure}[H]
\includegraphics[width=\linewidth]{./images/TAE/erste-hilfe-dif-schild-stange-2.jpg}
\Captionof{figure}{\AuthNotesPicLegalOk{}{}Licht, Menschenmassen und noch mehr
Hitze: Auch das ist das Donauinselfest.}
\end{figure}




% \Layout{27}{}{
% 
% }
% {}
% {}
% {}
% {}
% \begin{figure}[H]
% \includegraphics[width=\linewidth]{./images/TAE/jacken-alt-grau.jpg}
% \Captionof{figure}{\AuthNotesPicLegalOk{}{}%
% % In den letzten zehn Jahren hat sich viel getan: 
% Die alten, grauen Jacken erscheinen heute wie ein Relikt aus grauer
% Vorzeit.}
% \end{figure}



% \Layout{27}{}{
% 
% }
% {}
% {}
% {}
% {}
\begin{figure}[H]
\includegraphics[width=\linewidth]{./images/TAE/tae-okt26-bundesheer1.jpg}
\Captionof{figure}{\AuthNotesPicLegalOk{}{}Integraler Bestandteil der zivilen
und militärischen Landesverteidigung: Der Bundesheersanitäter.}
\end{figure}


\begin{figure}[H]
\includegraphics[width=\linewidth]{./images/gabriel-sebastian-aass/IMG_1853_1.jpg}
\Captionof{figure}{Die Zusammenarbeit mit der Exekutive -- auch an
sozialen Brennpunkten -- ist Alltag.}
\end{figure}

% \Layout{27}{}{
% 
% }
% {}
% {}
% {}
% {}
\begin{figure}[H]
\includegraphics[width=\linewidth]{./images/TAE/tae-teamwork-arzt-sani.jpg}
\Captionof{figure}{\AuthNotesPicLegalOk{}{}Arzt und Sanitäter bilden eine
Einheit.}
\end{figure}
\end{multicols} 
% \begin{savequote}
\includegraphics[width=4cm]{./images/wikipedia-pd/Vesalius_164frc.png}

{\tiny\itshape Antonius Vesalius}
\end{savequote}
\chapter
[\CO{[ANA]} Anatomie und Physiologie]
{\CC{[ANA]}\\Anatomie und Physiologie}
\label{chp:ANA}

\emph{Maintainer: Lena Hirtler}

% \AuthNotes{
% Changelog: 
% 
% \begin{description}
%     \item[2009-02-25] Start Changelog.
%     \item[2009-05-06] GABS: Korrekturen von KOCH
%     eingearbeitet.
% \end{description}
% }

\minitoc

\section{Die Zelle}\index{Zelle}
\Layout{0}{\TxBeschreibung}
{
Die Zelle ist die kleinste lebende Einheit in unserem
Körper. Sie besteht u.\,a. aus einem
\E{Zellkern}\index{Zellkern}, der die gesamten Erbinformationen in Form des
Riesenmoleküls \E{DNS}\ZFootnote{Desoxyribonukleinsäure, \Lang{engl.} \E{DNA}.}
enthält, sowie verschiedenen Zellorganellen\index{Zellorganellen}. Die Zellorganellen sind
wichtig für den Zellstoffwechsel, man kann sie mit der Funktion der Organe im
menschlichen Körper vergleichen. Der überwiegende Bestandteil der Zelle ist
\E{Flüßigkeit}, das Zellplasma. Die Zelle wird von der
\E{Zellmembran}\index{Zellmembran} umhüllt.
\cite{AlbertsMolZellbiologie2,LoefflerBiochemie7}
}{
\begin{itemize}
    \item Kleinste lebende Einheit.
    \item Zellkern
    \begin{itemize}
        \item Erbinformationen (DNS)
    \end{itemize}
    \item Zellorganellen
\end{itemize}
}
{}{}{}



\begin{figure}[H]
    \includegraphics[width=\linewidth]{./images/anatomie/Animal_cell_structure_de-edited.png}
    \Captionof{figure}{\AuthNotesPicLegalOk{}{PD}Eine
    Zelle, Schema. \SourcePicture{Mariana Ruiz Villarreal,
    Public domain}}
\end{figure}

\Layout{0}{Gewebe}{
Artverwandte Zellen können einen
Zellverband bilden, welcher \T{Gewebe}\index{Gewebe} genannt wird. Diese
Gewebe können dann größere Strukturen, wie z.\,B. Herz, Hirn, Lunge, Knochen,
Muskel, "`erbauen"'.
\A{
Diskussion: Was ist ein Organ?

\ACh{GABS}{2010-03-09}{Ist Muskel nicht ein Organ?}
\ACh{HIRT}{2010-03-09}{Wie du sicher weißt gibt es verschiedene Definitionen
von Organen.Bei den gängigsten fallen die Muskeln nicht hinein.}
\ACh{GABS}{2010-03-12}{Bei welcher? 

Pschyrembel: "`aus Zellen u. Geweben
zusammengesetzte Teile des Körpers, die eine Einheit mit best. Funktionen
bilden."' 

Wollen wir's nicht allgemein halten? Zumal die Definition"'Organ"' noch
ausständig ist \ldots{}
\ACh{HIRT}{2010-03-13}{Damit ist die Diskussion beendet.}
\ACh{GABS}{2010-03-14}{Neue Formulierung passt, weil unmissverständlich.}
}
}


}{
\begin{itemize}
    \item Zellverband.
    \item → größere Strukturen.
\end{itemize}
}{}{}{}

\nocite{UlfigHistologie1,LippertAnatomie7}


\section{Das Skelett und der Bewegungsapparat}
\index{Skelett}

\begin{figure}[H]
    \includegraphics[width=10cm]{./images/anatomie/Human_skeleton_front_de.png} \Captionof{figure}{\AuthNotesPicLegalOk{}{ PD }Das Menschliche Skelett.
    \SourcePicture{Mariana Ruiz Villarreal, Public domain}}

\end{figure}

\Layout{0}
{Funktionen des Skeletts}
{
Das Skelettsystem besteht aus \E{Knochen} und \E{Knorpeln} und dient 
hauptsächlich als

\begin{itemize}
    \item Stützgerüst  des Körpers, wichtig für die aufrechte Haltung
    \ACh{GABS}{2010-03-09}{Auch für nicht-aufrechte Haltung wichtig?}
    \ACh{HIRT}{2010-03-09}{Hier geht es um die grundsätzliche Vorstellung,
    dass wir ohne Knochen nur ein Haufen Gatsch sind.}
    \item Ursprungs- und Ansatzfläche für die Muskulatur, wichtig für
    Wirkungsentfaltung der Muskel
\end{itemize}

und weiters als

\begin{itemize}
    \item Kalziumspeicher (\ce{Ca++})
    \item Bildung von roten und weißen Blutkörperchen sowie den Blutplättchen
    (im roten Knochenmark!) \ACh{GABS}{2010-03-09}{Wenn man die Körperchen aufzählt müssen die
    Thrombos aber auch hinein \ldots}
    \ACh{HIRT}{2010-03-09}{Vielen vielen Dank für den Hinweis\ldots}
\end{itemize}
}
{
\begin{itemize}
  \item Haltung
  \item Muskelbewegung
  \item Kalziumspeicher
  \item Blutbildung
\end{itemize}
}
{}{}{}

\subsection{Der Knochen}
\index{Knochen}

\Layout{0}
{Knocheneinteilung}
{%
Man kann die Knochen anhand ihrer Form einteilen in: 
\index{Knochen!platte}\index{Knochen!Röhren-}\index{Knochen!Würfel-}
\begin{description}
    \item[\T{Platte Knochen}] (z.\,B. Beckenknochen,
    Schulterblatt, Brustbein, Schädel),
    \item[\T{Röhrenknochen}] (Oberarmknochen, Schienbein)
    und
    \item[\T{Würfelknochen}] (z.\,B. Hand-, Fußwurzelwurzelknochen,
    Wirbelknochen)
\end{description}
}
{
\begin{itemize}
  \item Platte Knochen
  \item Röhrenknochen
  \item Würfelknochen
\end{itemize}
}
{}{}{}

\AuthNotes{GABS \& KOCH: evtl. umstellen auf deutsch -- latein, aber was heißt
Kortikalis auf deutsch?}\AuthNotes{HIRT: Man könnte sie als Knochenrinde
bezeichnen.}

\Layout{20}
{Aufbau eines Knochens}
{%
\subparagraph{\T{Periost}} \index{Periost}\E{Beinhaut}, überzieht den Knochen und
ist sehr gut mit Nerven versorgt. Es ist daher die einzige Struktur des Knochens die
schmerzempfindlich ist. Verantwortlich für Dickenwachstum und Frakturheilung
\subparagraph{\T{Kompakta}} (Kortikalis, "`Knochenrinde"'),
\index{Kompakta}\index{Kortikalis}feste äußere Schicht,
äußere Form, Verantwortlich für Stabilität des Knochens
\subparagraph{\T{Spongiosa}} \index{Spongiosa}\index{Knochenbälkchen}Knochenbälkchen
${\rightarrow}$ Zug- und Druckfestigkeit
\subparagraph{\T{Markhöhle}} \index{Markhöhle}s.u.
\subparagraph{\T{Epiphysenfuge}} \index{Epiphysenfuge}Wachstumsfuge
\subparagraph{\T{Gelenkfläche}} \index{Gelenkfläche!Aufbau des
Knochens}mit knorpeligen Überzug als Teil eines Gelenkes.
Manche Knochen haben keine Gelenksfläche.


Lange Röhrenknochen kann man außerdem in 
\begin{inparaitem}
    \item \T{Kopf}\index{Kopf!Knochenaufbau},
    \item \T{Hals}\index{Hals!Knochenaufbau} und
    \item \T{Schaft}\index{Schaft!Knochenaufbau} 
\end{inparaitem}
unterteilen.
}{
\begin{itemize}
  \item Periost
  \item Kompakta
  \item Spongiosa
  \item Markhöhle
  \item Epiphysenfuge
  \item Gelenkfläche
\end{itemize}
}
{./images/anatomie/roehrenknochen-querschnitt.png}
{Ein langer Röhrenknochen im Querschnitt.}
{Hirtler}

\subsubsection{Das Knochenmark}
\index{Knochenmark}


% \fxnote[author=HIRT]{Review: ANA: Knochenmark, rotes: Entweder lange oder kurze
% Überschrift.} 
% GABS 2010-03-12: ok.

\Layout{0}
{Rotes Knochenmark}
{%
% 
% Bitte Index-Einträge in den Fließtext hineingeben, nicht in Überschrift
% (es gibt noch alte Stellen wo die Einträge in der Überschrift stehen,  aber
% das wird nach und nach geändert!) GABS 2010-03-12
% 
\index{Knochenmark!rotes}
Es ist das wichtigste blutbildende Organ
des Körpers. Hier werden die \Es{zellulären Bestandteile des Blutes
gebildet}. Dazu gehören die \E{Roten Blutkörperchen} (\E{Erythrozyten}), die
\E{weißen Blutkörperchen} (\E{Leukozyten}) sowie die \E{Blutplättchen}
(\E{Thrombozyten}) -- \prettyref{topic:blut}. 

Bei Neugeborenen findet man das rote Knochenmark in
sämtlichen Knochen, im Laufe des Lebens bildet es sich zurück und verwandelt
sich in {\quotedblbase}Gelbes Fettmark{\textquotedblleft}. Rotes Knochenmark
findet man dann lediglich in platten Knochen wie Becken, Brustbein und
Schädel, wo es für die Blutbildung verantwortlich ist und normalerweise
zeitlebens erhalten bleibt.

\TODO{GABS}{yyyy-mm-dd}{Review: 2010-03-12: Alte Überschrift als Fazit
eingefügt:}{}
\Fazit{Das \E{rote Knochenmark} ist für die Blutbildung zuständig}

}
{%
\begin{itemize}
  \item Bildung der zellulären Blutbestandteile
  	\subitem rote Blutkörperchen
  	\subitem weiße Blutkörperchen
  	\subitem Blutplättchen
  \item im Erwachsenenalter nur noch in platten Knochen  
\end{itemize}
}
{}{}{}

\AuthNotes{KOCH: wichtig? ev. weglassen?

GABS: Ja, wichtig wegen KT, onko-Patienten, etc. Soll
bleiben. 

KOCH: ok.} 

\Fazit{Einige medizinische Behandlungsmethoden
verursachen als Nebenwirkung eine massive \Es{Schädigung
des Knochenmarks}\index{Knochenmark!Schädigung} (z.\,B.
Bestrahlung, Chemotherapien bei Krebspatienten ("`onkologische
Patienten"'), etc.). Diese Patienten leiden neben einer
Blutarmut auch an einem schwachen Immunsystem, da nicht genügend weiße
Blutkörperchen gebildet werden. Sie können daher sehr leicht schwerwiegende
Infektionen bekommen.

Diese Patienten
sind im Krankentransport häufig anzutreffen.} 

% \fxnote[author=HIRT]{Review: ANA: Knochenmark, gelbes: Entweder lange oder kurze
% Überschrift.} 
% GABS 2010-03-12: ok.

\Layout{0}
{Gelbes Knochenmark\index{Knochenmark!gelbes}}
{%
Im Laufe des Lebens bildet sich in den
Röhrenknochen das rote Knochenmark  zu fettreichem, gelben
Knochenmark um.

Das gelbe Fettmark kann bei schweren Knochenverletzungen von Bedeutung
sein, da verletzungs\-bedingt kleine Fetttröpfchen in die
Blutbahn gelangen können, die dann als Fettgerinnsel zu ernsten
Komplikationen führen können
(Fettembolie\AuthNotes{+ KOCH}).

Dieses Knochenmark ist nicht unbedingt für den Körper notwendig - deshalb kann
man nach Frakturen auch Nägel in die Markhöhle einsetzen ohne größere
Auswirkungen auf den Knochen zu verursachen.

\TODO{GABS}{yyyy-mm-dd}{Review: 2010-03-12: Alte Überschrift als Fazit
eingefügt:}{}
\Fazit{Das \E{gelbe Knochenmark} besteht aus Fett}
}
{%
\begin{itemize}
  \item fettreich
  \item Grund für Komplikationen bei Frakturen
\end{itemize}
}
{}{}{}

\subsubsection{Knochenwachstum}~

\Layout{0}
{Knochenentwicklung}
{%
Das Skelett eines Neugeborenen besteht bei der Geburt zum größten
Teil aus Knorpel. Erst mit zunehmendem Lebensalter beginnt dieser von
Knochenkernen ausgehend, zu verknöchern. Eine
\Es{Wachstumszone}
(\index{Epiphysenfuge}\T{Epiphysenfuge}) bleibt dabei allerdings erhalten.

Die Epiphysenfuge enthält beim Kind/Jugendlichen bis
etwa zum Ende der Pubertät \E{teilungsfähiges
Knorpelgewebe}, daher ist ein \Es{Längenwachstum}
möglich. Beim Erwachsenen sind diese knorpeligen Fugen
hingegen vollständig \E{verknöchert}, sodass
keinerlei Wachstum mehr stattfinden kann.
}
{%
\begin{itemize}
  \item Geburt: Knorpel
  \item Knochenkerne
  \item Wachstumszone
  \item Verknöcherung der Wachstumszone
\end{itemize}
}
{}{}{}

\subsubsection{Gelenke}
\index{Gelenke}

\Definition{(Meistens bewegliche) Verbindung zwischen
mindestens zwei Knochen.}

\Layout{0}
{Allgemeine Gelenklehre}
{Bei den Gelenken kann man zwischen Gelenken \E{mit Gelenksspalt} ("`echtes
Gelenk"', \Z{Diarthrose}) und Gelenken \E{ohne Gelenksspalt} ("`unechtes
Gelenk"', \Z{Synarthrose}) unterscheiden.

Zu den Gelenken ohne Gelenksspalt gehören alle Knochenverbindungen die \E{nicht}
frei beweglich sind, z.\,B. die Schädelnähte, die Schambeinfuge und die
Rippen-Brustbein-Verbindungen des Thorax (\Z{Syndesmose, Synchondrose,
Synostose}).
\ACh{HIRT}{2010-03-11}{Um allen Diskussionen zuvorzukommen: ich finde es
wichtig, dass zumindest diese Tatsachen erwähnt werden.}
\ACh{GABS}{2010-06-07}{Fachbegriffe wie besprochen grau.}

% orientierende Zusammenfassung (GABS, 2010-06-07)
\Fazit{Das, was man umgangssprachlich unter dem Begriff "`Gelenk"' versteht,
sind "`echte"' Gelenke mit Gelenksspalt.} 
}
{%
\begin{itemize}
  \item "`echtes Gelenk"'
  \item "`unechtes Gelenk"'
\end{itemize}
}{}{}{}

\Layout{20}
{Prinzipieller Aufbau eines "`echten Gelenks"'}
{%
Ein Gelenk ist von einer bindegewebigen Kapsel umgeben und durch
Sehnen und Bänder stabil verbunden.

Die Gelenksflächen sind zur Reibungs\-ver\-minderung mit 
Knorpel überzogen und der Gelenksspalt mit einer dünnen Schicht 
Schmierflüssigkeit \Z{(Synovialflüssigkeit)} gefüllt. 
}{%
\begin{itemize}
  \item Kapsel
  \item Knorpelüberzug der Gelenkflächen
  \item Gelenksflüssigkeit
  \item Gelenkarten
  \begin{itemize}
    \item Kugelgelenk
    \item Sattelgelenk
    \item Scharniergelenk
  \end{itemize}
\end{itemize}
}
{./images/anatomie/anatomie-gelenk-edited.png}
{\AuthNotesPicLegalOk{}{Lena }Ein Gelenk im Querschnitt.}
{Hirtler}
\AuthNotes{KOCH. löschen: Membrana synovialis und membrana fibrosa}


\paragraph{Es gibt verschiedene Arten von Gelenken}
\index{Gelenke!Arten}

\index{Kugelgelenk}
\index{Eigelenk}
\index{Scharniergelenk}
\index{Sattelgelenk}
\index{Zapfengelenk|Z}
Klassische Beispiele sind das 
\T{Kugelgelenk} (Hüftgelenk),
 das \T{Eigelenk} (z.\,B. Handwurzelgelenk),
das das \T{Sattelgelenk} (Daumengrundgelenk), 
das \T{Scharniergelenk} (Fingergliedgelenke, Ellbogengelenk) \Z{sowie das
Zapfengelenk (z.\,B. Radioulnargelenk)}.


\noindent\begin{minipage}{\linewidth}
\includegraphics[width=\linewidth]{./images/hirtler-aass/Gelenkarten.pdf}
% \TODO{GABS}{2010-04-11}{Bild: von HIRT anpassen und einfuegen.}{}
% DONE 2010-06-07 GABS

\CaptionofDiff{figure}{Arten von Gelenken. 
\begin{inparadesc}
  \item[A:] Kugelgelenk (z.\,B. Hüftgelenk),
  \item[B:] Eigelenk (z.\,B. Handwurzelgelenk),
  \item[C:] Sattelgelenk (z.\,B. Daumengrundgelenk),
  \item[D:] Scharniergelenk (z.\,B. Fingergliedgelenke, Ellbogengelenk),
  \item[E:] \Z{Zapfengelenk (z.\,B. Radioulnargelenk)}.
\end{inparadesc}
\SourcePicture{Hirtler.}
}{Arten von Gelenken.}
\end{minipage}


\subsection{Der Schädel} 
%%% Index und Label
\index{Schädel}
\index{Caput}
\index{Cranium}
\label{topic:schaedel}


Kopf (\T{Caput}) / Schädel (\T{Cranium})

 


\begin{gWideParEnv}
\Parallel{
\includegraphics[width=\linewidth]{./images/anatomie/Human_skull_side_simplified_hirngesicht_de-edited.png}
\Captionof{figure}{\AuthNotesPicLegalOk{}{PD}
\index{Hirnschädel}\index{Gesichtsschädel}Der Schädel
gliedert sich in den \T{Hirnschädel} und den \T{Gesichtsschädel}}. 
\SourcePicture{Mariana Ruiz Villarreal, Public Domain}
}
{%
\includegraphics[width=\linewidth]{./images/anatomie/Human_skull_front_simplified_bones_de-edited.png}
\Captionof{figure}{\AuthNotesPicLegalOk{}{PD}Schädel Vorderansicht.
\SourcePicture{Mariana Ruiz Villarreal, Public Domain}}
% \includegraphics[width=\linewidth]{./images/anatomie/Human_skull_side_simplified_bones_de-edited.png}
% \Captionof{figure}{\AuthNotesPicLegalOk{}{PD}Schädel seitlich.
% \SourcePicture{Mariana Ruiz Villarreal, Public Domain}}
}\end{gWideParEnv}

\A{STEU 2009-06-28: Genaue Bezeichnungen der einzelnen
Schädelknochen wirklich notwendig? Unterkiefer, Oberkiefer,
Kalotte.

EMHO: Schrift grau. }

\Layout{0}{Aufbau des Hirnschädels}
{%
Der Hirnschädel besteht aus der \Es{Schädelbasis} und dem \Es{Schädeldach}
(Kalotte). Beide Strukturen werden durch platte Knochen gebildet, welche
beim Erwachsenen durch Schädelnähte  \index{Schädelnähte}  miteinander verzahnt
und verknöchert sind.
}{%
\begin{itemize}
    \item Schädelbasis
    \item Schädeldach (Kalotte)
    \item Platte Knochen, durch Schädelnähte verzahnt \& verknöchert
\end{itemize}
}
{}{}{}



\Layout{20}{Fontanelle}{%
\index{Fontanelle}%
Die \T{Fontanellen} sind \Es{knorpelige 
Verbindungen} zwischen Schädelknochen beim Neu\-geborenen
(der Schädel muss sich für Geburt verformen
können!). \Z{Sie schließen  sich bis zum 18.
Lebensmonat,} verknöchern  aber erst vollständig mit
zunehmendem  Lebensalter \Z{(ca. 20 Lebensjahr
abgeschlossen)}.

\Z{Es gibt zwei Fontanellen: Die 
\begin{inparaitem}
\item \E{große} Fontanelle befindet sich zwischen dem Stirnbein und dem
Scheitelbeinen, die
\item \E{kleine} Fontanelle zwischen den Scheitelbeinen und dem
Hinterhauptsbein.
\end{inparaitem}
Die Unterscheidung ist für die Geburtshilfe interessant, da man dadurch
unterscheiden kann in welche Richtung der Kopf des Kindes schaut.} 

\vspace{1cm}

}
{
\begin{itemize}
  \item Verschluss bis 18. Lebensmonat
  \item große Fontanelle
  \item kleine Fontanelle
\end{itemize}
}
{./images/anatomie/fontanelle-dt.jpg}
{Fontanellen.}
{Gray's Anatomy, Copyright expired.}

\Z{Die \T{Sutur} ist die endgültig fest verknöcherte Naht
zwischen den Schädelknochen beim Erwachsenen ${\rightarrow}$ starrer  stabiler
Schädel, Schutz  für Gehirn.}


\Layout{20}{Schädelbasis}{%
% a ... Stirnbein (Os frontale)
% 
% b ... Siebbein (Os ethmoidale)
% 
% c ... Keilbein  (Os sphenoidale)
% 
% d ... Schläfenbein(e)  (Os temporale)
% 
% e ... Felsenbein  (Pars petrosa ossis temporalis)
% 
% f ... Hinterhauptsbein  (Os occipitale)
% 
% g ... großes Hinterhauptsloch (Foramen magnum)
Das Gehirn liegt auf der Schädelbasis \index{Schädelbasis}
auf \Z{(außer wenn man auf dem Kopf steht)}. Hirnnerven treten durch die diversen Löcher aus dem
Schädel aus. \Es{Das \T{große
Hinterhauptsloch}\index{Großes
Hinterhauptsloch}  \Z{(Foramen magnum)} ist die
Durchtrittsstelle des Rückenmarks.} 

\subparagraph{Exkurs: \T{Hirnstammeinklemmung}}
\index{Hirnstammeinklemmung} 
Das \Es{große Hinterhauptsloch} ist die Durchtrittsstelle des Rückenmarks. Bei
einer \Es{Hirndrucksteigerung} (\E{Hirnödem, Hirnblutung}) ist es die einzige
Möglichkeit zur Ausdehnung des Gehirns! \Es{Dabei drückt sich das Gehirn
allerdings selber ab}, es kommt zu einer tödlichen \Ca{Hirnstammeinklemmung!}
${\rightarrow}$ \Ca{Lebensgefahr!}
}
{
\begin{itemize}
  \item Auflagefläche des Gehirns
  \item Hirnnerven
  \item großes Hinterhauptsloch (Foramen magnum)
  \item Hirnstammeinklemmung
\end{itemize}
}
{./images/anatomie/cavernous_sinus.png}
{Aufsicht auf die Schädelbasis.}
{Gray's Anatomy, Copyright expired.}



\subsection{Die Wirbelsäule}
\index{Wirbelsäule}
\Layout{20}{\TxBeschreibung}
{%
Die Wirbelsäule bildet das knöcherne Rückgrat des Körpers und ist damit eine
der wichtigsten Stützen des Körpers, andererseits ist sie auch Ansatzfläche
vieler wichtiger Muskeln. Sie hat eine doppelt \E{S-förmig gekrümmte Form},
welche den \E{aufrechten Gang} ermöglicht.

\subparagraph{Bestandteile}
Die Wirbelsäule besteht aus mehreren Komponenten. Die knöchernen, übereinander
"`gestapelten"' \T{Wirbel} \Z{(Vertebrae)} geben der Wirbelsäule ihre Form
und Stabilität. Die Wirbel sind i.d.R. zueinander im gewissen Umfang
\E{beweglich}.

Zwischen den
einzelnen Wirbeln finden sich \index{Bandscheiben}\T{Bandscheiben}
(\T{Discus}), welche als \Es{Stoßdämpfer} dienen. Mit zunehmenden Alter können
diese \E{verschleißen}. 
    
Von den Wirbeln wird ein Kanal, der sog. \T{Wirbelkanal}\index{Wirbelkanal}
gebildet. In diesem ist das \Es{Rückenmark} eingebettet.

\subparagraph{Gliederung} Die Wirbelsäule gliedert sich in verschiedene
Abschnitte: 
\index{Halswirbel}
\index{Brustwirbel}
\index{Lendenwirbeln}
\index{Kreuzbeinwirbel}
\index{Steißbeinwirbel}
Es gibt 
\begin{inparaitem}
    \item 7 \Es{Hals}wirbel \Z{(Cervical-Wirbel)},
    \item 12 \Es{Brust}wirbel  \Z{(Thorakal-Wirbel)},
    \item 5 \Es{Lenden}wirbel  \Z{(Lumbal-Wirbel)},
    \item 5 verschmolzene \Es{Kreuzbein}wirbel \Z{(Os sacrum)}, sowie
    \item 3-4 \Es{Steißbein}"`wirbel"' \Z{(Os coccygis)}.
\end{inparaitem}

\subparagraph{Klinischer Bezug} Wenn die Wirbelsäule beschädigt wird kann es
gleichzeitig zu einer Beschädigung des Rückenmarks kommen
(\E{Wirbelsäulentrauma}, \prettyref{topic:wirbelsaeulentrauma}). Bandscheiben
können aufplatzen und der weichere Kern \E{in den Wirbelkanal rutschen}. Durch
den Druck auf das Rückenmark kann das starke Schmerzen und
Nervenausfälle verursachen (\E{Bandscheibenvorfall},
\prettyref{topic:lumbago}).


}{%
Doppelt S-förmig gekrümmtes
Achsenskelett, ermöglicht  aufrechten Gang 

\TabSub{Bestandteile}

\begin{itemize}
    \item Wirbel \Z{(Vertebrae)}
    \item \index{Bandscheiben}\T{Bandscheiben}
    (\T{Discus}): dienen der "`Stoßdämpfung"'
    \item \Es{Im Wirbelkanal liegt das Rückenmark
    eingebettet.}
\end{itemize}


\TabSub{Gliederung}

\begin{itemize}
    \item 7 \Es{Hals}wirbel \Z{(Cervical-Wirbel)}
    \index{Halswirbel}
    \item 12 \Es{Brust}wirbel  \Z{(Thorakal-Wirbel)}
    \index{Brustwirbel}
    \item 5 \Es{Lenden}wirbel  \Z{(Lumbal-Wirbel)}
    \index{Lendenwirbeln}
    \item 5 verschmolzene \Es{Kreuzbein}wirbel \Z{(Os
    sacrum)} \index{Kreuzbeinwirbel}
    \item 3-4 \Es{Steißbein}"`wirbel"' 
    \Z{(Os coccygis)} \index{Steißbeinwirbel}
\end{itemize}

Bandscheibenvorfall, Wirbelsäulentrauma
}
{./images/anatomie/wirbelsaeule-lena-edited.png}
{\AuthNotesPicLegalOk{}{}Die Wirbelsäule.}
{Hirtler. }





\subsubsection{Die Wirbel: Aufbau und Arten}~
\index{Wirbel}

\Layout{20}{\TxBeschreibung}
{%
Ein Wirbel besteht aus einem 
\begin{inparaitem}
  \item \T{Wirbelkörper}\index{Wirbelkörper} und einem
  \item \T{Wirbelbogen}\index{Wirbelbogen}, welcher das Rückenmark umschließt.
  \end{inparaitem}
Vom Wirbelbogen entspringen
  \begin{inparaitem}
    \item 2 \T{Querfortsätze}\index{Querfortsätze} und
    \item 1 \T{Dornfortsatz}\index{Dornfortsatz}.
  \end{inparaitem}


Der Wirbekörper trägt die Last des Rumpfes, zwischen ihnen liegen die
\E{Bandscheiben}. Der Wirbelbogen \Es{schützt das Rückenmark} vor Verletzungen,
die Querfortsätze dienen z.\,B. zur Befestigung der Rippen. Die Dornfortsätze sind
jene Teile der Wirbelsäule, die am Rücken tastbar sind.


In den jeweiligen Abschnitten sehen die Wirbel
unterschiedlich aus.
}
{%
Aufbau:
\begin{itemize}
  \item Wirbelkörper
  \item Wirbelbogen
  \begin{itemize}
    \item 2 Querfortsätze
    \item 1 Dornfortsatz
  \end{itemize}
\end{itemize}
}{./images/anatomie/Gray94.png}
{Typischer Wirbelkörper.}
{Gray's Anatomy, Copyright expired.}






\subsection{Der Brustkorb -- Thorax}
\index{Brustkorb}
\index{Thorax}
\index{Brustbein}
\index{Sternum}
\label{topic:brustkorb-skelett}



\Layout{20}{\TxBeschreibung}
{%
Der  Brustkorb besteht \E{vorne} aus dem \T{Brustbein}
 (\T{Sternum}) und \Es{12
Rippenpaaren}, welche \E{hinten} mit der \Es{Brustwirbelsäule} (12
Brust\-wirbel) in gelenkiger Verbindung stehen.

\index{Rippen}
\index{Rippen!echte}
\index{Rippen!unechte}
\index{Rippen!fliegende}
Die obersten 7 Rippen-Paare werden als 
{\quotedblbase}\T{echte Rippen}{\textquotedblleft}
bezeichnet. Sie erreichen mit ihren Spitzen das Brustbein und bilden praktisch
{\quotedblbase}Ringe{\textquotedblleft} um die Brustorgane. Die 5
unteren Rippenpaare bilden den bogenförmigen unteren Rand des
Brustkorbes. Die Rippenpaare 8-10 sind dabei nur indirekt
über einen Knorpel mit dem Brustbein verbunden
({\quotedblbase}\T{unechte Rippen}{\textquotedblleft},
\emph{{\quotedblbase}echt{\textquotedblleft}} bedeutet, die Rippe hat eine
direkte Verbindung zum Sternum), das 11. und das 12. Paar endet mit knorpeligen Spitzen frei zwischen den
Bauchmuskeln({\quotedblbase}\T{fliegende
Rippen}{\textquotedblleft}).

Der Brustkorb enthält und schützt die Brust\-organe (Lunge, Herz)
sowie die Oberbauchorgane (Leber, Magen). Zwischen den
Rippen spannt sich die
\T{Zwischenrippenmuskulatur}\index{Zwischenrippenmuskulatur} (wichtig für \Es{Atemmechanik}), in diese sind segmentale
Nerven- und Gefäßbündel eingelagert. Das \Es{Zwerchfell}
\index{Zwerchfell} schließt den Brustraum nach unten hin ab
und trennt somit den Brustraum vom Bauchraum.

Die Intaktheit des Brustkorbes ist für die Atmung von großer Bedeutung. Nach
Rippenfrakturen kann es durch spitze Fragmente zu einer
Lungenfell- oder Lungenverletzung kommen, die lebensgefährlich werden kann.
}
{
\begin{itemize}
    \item \Es{Brustbein} (\T{Sternum})
    \item \Es{12  Rippenpaare}
    \begin{itemize}
        \item \Z{7 "`echte"'}
        \item \Z{3 "`unechte"'}
        \item \Z{2 "`fliegende"'}
    \end{itemize}
    \item \Es{Brustwirbelsäule} (12 Brust\-wirbel), mit Rippen in gelenkiger Verbindung
\end{itemize}
}
{./images/anatomie/thorax.png}
{Der knöcherne Thorax.}
{Mariana Ruiz Villarreal, Public Domain}



\subsection{Der Schultergürtel und die obere Extremität}
\index{Schultergürtel}
\label{topic:schulterguertel}

\begin{figure}[H]
    \begin{center}
    \includegraphics[width=\linewidth]{./images/anatomie/Human_arm_bones_diagram_de-edited.png}
\Captionof{figure}{\AuthNotesPicLegalOk{}{PD}Der
Schultergürtel und die obere Extremität.
\SourcePicture{Mariana Ruiz Villarreal, Public Domain}}
    \end{center}
\end{figure}

\subsubsection{Der Schultergürtel} 

\Layout{0}{\TxBeschreibung}
{
Der Schultergürtel besteht (von vorne nach
hinten) aus :

\begin{itemize}
    \item \index{Brustbein}\T{Brustbein}
    (\index{Sternum}\T{Sternum})
    \item \index{Schlüsselbein}\T{Schlüsselbein}
    \Z{(Clavicula)}
    \item \index{Schulterblatt}\T{Schulterblatt} \Z{(Scapula)}
\end{itemize}

Der Schultergürtel ermöglicht die Bewegung der Arme. Über eine Bandverbindung
des Schlüsselbeines am Brustbein und wiederum eine Bandverbindung zwischen
Schlüsselbein und Schulterblatt sind alle Knochen miteinander verbunden. Die
Gelenkspfanne für den Oberarmknochen befindet sich am Schulterblatt.

Der Oberarm ist hauptsächlich durch Muskelzug \Z{(M. deltoideus,
Rotatorenmanschette)} am Schultergürtel befestigt.

\Fazit{Bänder und Muskel sind für die Stabilität des Schultergelenkes wichtig.}
}
{
\begin{itemize}
    \item Brustbein (Sternum)
    \item Schlüsselbein (Clavicula)
    \item Schulterblatt(Scapula)
    \item Befestigung des Oberarmes durch Muskelzug
\end{itemize}
}{}{}{}

\subsubsection{Die obere Extremität}
\index{Extremität!obere}
\label{topic:obere-extremitaet}

\Layout{0}{\TxBeschreibung}
{
Der Arm beginnt mit dem \T{Schultergelenk}. Dieses ist ein klassisches
\E{Kugelgelenk}, es ermöglicht die Bewegung in alle drei Raumebenen. Die
beteiligten Knochen im Schulterglenk sind das \E{Schulterblatt} (Scapula) und
\T{Oberarmknochen} (Humerus). Letzterer bildet als einzelner Knochen den
Oberarm.

Das \T{Ellenbogengelenk} liegt zwischen Oberarm und Unterarm. Es wird
aus \Es{Oberarmknochen} und \Es{Elle} (Ulna) sowie \Es{Speiche} (Radius)
gebildet. Das Ellenbogengelenk ist ein zusammengesetztes Gelenk. Die Verbindung zwischen
Elle und Oberarmknochen ist ein Scharniergelenk und ermöglicht das Abwinkeln des
Armes. Die Verbindung zwischen Speiche und Oberarmknochen hingegen ist ein
Kugelgelenk. In Kombination mit dem Zapfengelenk zwischen Elle und Speiche
ermöglicht dieses Gelenk das Drehen des Unterarmes und der Hand.

Die beiden Knochen des \Es{Unterarmes} sind die \T{Elle} (Ulna) und die
\T{Speiche} (\T{Radius}) sind, \Z{sie liegen bei Drehung der Handfläche nach
oben parallel, bei Drehung der Handfläche zum Boden hingegen überkreuz.} Die \Es{Speiche} (Radius) liegt
\Es{daumenseitig}\footnote{Daher ist die Radialarterie daumenseitig zu tasten!},
die Elle an der Seite des kleinen Fingers.

Es folgt das \Es{Handwurzelgelenk}, welches hauptsächlich aus
der Speiche und den \E{Handwurzelknochen} gebildet wird. Das Handwurzelgelenk
ist ein \Es{Eigelenk}, \Z{es hat fast so viele Freiheitsgrade wie ein Kugelgelenk, in
einer Ebene ist jedoch die Bewegung durch die Ei-Form eingeschränkt
(Kippbewegungen in Richtung von Daumen und Kleinfinger).} 
Die \Es{Hand} beginnt auf Höhe des Gelenkes.

Es gibt insgesamt 8
\T{Handwurzelknochen}, Würfelknochen, \Z{welche in zwei Reihen zu jeweils 4
Knochen angeordnet sind.}

Hierauf folgen \Es{5 \T{Mittelhandknochen}} und dann die 5 Finger mit jeweils 2
(Daumen) oder 3 (Finger \VonBis{II}{V}) \T{Fingergliedknochen} (\E{körpernahe},
\E{mittlere} und \E{körperferne} Fingergliedknochen). Die Grundgelenke
(Mittelhandknochen $\leftrightarrow$ körpernahe Fingergliedknochn) sind
Sattelgelenke, die \Es{Fingergliedgelenke} sind \Es{Scharniergelenke}.
}{
\begin{itemize}
    \item \Es{Schultergelenk} \index{Schultergelenk}
    \item \Es{Oberarmknochen} \index{Oberarmknochen} 
    (\T{Humerus}\index{Humerus})
    \item \Es{Ellenbogengelenk} \index{Ellenbogengelenk}
    \item (2) \Es{Unterarmknochen} \index{Unterarmknochen} 
    \begin{itemize}
        \item \Es{Elle} \index{Elle}
        (\T{Ulna}\index{Ulna})
        \item \Es{Speiche} \index{Speiche}
        (\T{Radius}\index{Radius}), daumenseitig!
    \end{itemize}
    \item Handwurzelgelenk \index{Handwurzelgelenk}
    \item (8) \Es{Handwurzelknochen}
    \index{Handwurzelknochen}
    \item (5) \Es{Mittelhandknochen}
    \index{Mittelhandknochen}
    \item \Es{Fingergliedknochen}
    \index{Fingergliedknochen} mit \VonBis{2}{3} Gliedern
\end{itemize}
}
{}{}{}

\AuthNotes{GABS: Bild evtl. nur mit einer Seite, dafür aber
größer und neben Text?

KOCH: passt so.}

\Fazit{Achtung: Unterscheide: Hand --
Arm!\index{Hand}\index{Arm}}

\subsection{Der Beckengürtel und die untere Extremität}
\index{Beckengürtel}
\index{Extremität, untere}
\label{topic:beckenguertel}

\Layout{0}{Der Beckengürtel}
{
Der Beckengürtel besteht (von hinten nach vorne) aus dem

\begin{inparaitem}
    \item \T{Kreuzbein} (Os sacrum) und dem 
    \item \T{Beckenknochen}\Z{, welcher wiederum aus dem 
\begin{inparaitem}
        \item \Z{\index{Darmbein|Z}Darmbein, dem}
        \item \Z{\index{Sitzbein|Z}Sitzbein und dem }
        \item \Z{\index{Schambein|Z}Schambein.}
    \end{inparaitem}
    besteht.}
\end{inparaitem}

\Z{Die \index{Symphyse|Z}Symphyse stellt eine straffe
Knorpelverbindung vorne zwischen den beiden Beckenknochen
dar (Dehnungsfuge für Geburt).}

Die Stabilität des Beckengürtels ist für Stehen, Gehen und alle sonstigen
Bewegungen der Beine unabdingbar. Die Gelenkspfanne des Hüftgelenkes befindet
sich im Beckenknochen. Der Oberschenkelknochen wird hauptsächlich durch Bänder
im Gelenk gehalten.

\noindent Nach einem Unfall kann es zu einer Fraktur des Beckengürtels kommen
(Beckentrauma, \prettyref{topic:beckentrauma}. Stehen und Gehen ist dann nicht
mehr möglich, der Patient verspürt starke Schmerzen bei jeder passiven
Bewegung \E{beider} Beine. 

}
{
\begin{itemize}
    \item Kreuzbein (Os sacrum)
    \item Beckenknochen
    \begin{itemize}
        \item \Z{Darmbein}
        \item \Z{Sitzbein}
        \item \Z{Schambein}
    \end{itemize}
    \item Symphyse
\end{itemize}
}
{}{}{}

\Layout{20}{Die untere Extremität}
{%
\label{topic:untere-extremitaet}%
\index{Schenkelhals}
Das Bein beginnt mit dem \T{Hüftgelenk}, einem \Es{Kugelgelenk}, gebildet aus
dem \E{Beckenknochen} und dem \E{Oberschenkelknochen}. Im Gegensatz zum
Schultergelenk ist die Pfanne im Hüftgelenk besser ausgebildet und die
Befestigung des Gelenkskopfes in der Pfanne wird durch starke Bänder
bewerkstelligt. 

Der \T{Oberschenkelknochen} (Femur) bildet als einziger Knochen den Oberschenkel.
Er ist gekenzeichnet durch einen ausgeprägten \T{Schenkelhals} mit "`Knick"'
zwischen Kopf-Hals-Bereich und Körper. Körperfern endet der Oberschenkelknochen
im Kniegelenk. 

 in den Unterschenkel - bestehend  über.

Das \T{Kniegelenk} verbindet den Ober- und Unterschenkel Und wird aus dem
Oberschnkelknochen und dem Schienbein (Unterschenkel) gebildet. Es ist \E{kein}
reines Scharniergelenk, Drehungen sind bei gebeugtem Knie möglich. 
% Im Kniegelenk sind
% nur Oberschenkelknochen und Schienbein beteiligt, nicht das Wadenbein. 
Die
Besonderheit des Kniegelenkes besteht einerseits in den Bandbefestigungen:
Kreuzbänder (2) und Seitenbänder (2) werden gemeinsam mit den Menisken (2) häufig
bei diversen Sportarten verletzt. Andererseits gibt es die \Es{Kniescheibe}, sie
dient der besseren Kraftübertragung des Kniestreckers.

Der \T{Unterschenkel} wird vom \T{Schienbein} (Tibia; dick, vorne) und
\T{Wadenbein} (Fibula; dünn, seitlich hinten) gebildet. Sie gehen im Sprunggelenk
in den Fuß über und bilden dort den Innen- (Schienbein) und Außenknöchel
(Wadenbein).

Im \T{Sprunggelenk} gibt es, ähnlich wie bei der Hand, \T{Fußwurzelknochen}
(7, \Z{der bekannteste ist das Fersenbein}), \T{Mittelfußknochen} (5) und 5
Zehen mit jeweils 2 (Großzehe) bis 3 Gliedern. 
}{
\begin{itemize}
    \item \Es{Hüftgelenk} \index{Hüftgelenk} 
    \item \Es{Oberschenkel} \index{Oberschenkel}  
    (\T{Femur})
    \index{Femur} 
    \item \Es{Kniegelenk} \index{Kniegelenk} 
    \item \Es{Kniescheibe} \index{Kniescheibe} 
    \item \Es{Unterschenkelknochen} (2)
    \index{Unterschenkelknochen}

    \begin{itemize}
        \item \Es{Schienbein}  (Tibia) , dick, vorne
        \index{Schienbein}  \index{Tibia} 
        \item \Es{Wadenbein}  (Fibula), dünn, seitlich
        hinten \index{Wadenbein}  \index{Fibula} 
    \end{itemize}
    \item \Es{Sprunggelenk} \index{Sprunggelenk} 
    \item 7 \Es{Fußwurzelknochen}, \Z{u.a. Sprungbein,
    Fersenbein} \index{Fußwurzelknochen} 
    \item 5 \Es{Mittelfußknochen}  
    \index{Mittelfußknochen}
    \item \Es{Zehengliederknochen} mit je 2-3 Gliedern
    \index{Zehengliederknochen} 
\end{itemize}
}
{./images/anatomie/untere-extremitaet.png}
{\AuthNotesPicLegalOk{}{}Der Beckengürtel und die untere Extremität.}
{Mariana Ruiz Villarreal, Public Domain}


% GABS, 2010/05/04: Die Beschreibung bezgl. Trauma und Massnahmen ist glaub ich
% im Trauma-Kapitel besser aufgehoben, der Text in TRA ist eh noch nciht berühmt
% -> dorthin verschoben.

% % GABS: Ad Trauma: 
% Eine Kompression des Beckens ist wichtig um den
% möglichen großen Blutverlust (einige Liter!) einzudämmen.

Als \Es{Knöchel} bezeichnet man die Knochengabel, welche aus Schien- und
Wadenbein auf Höhe des Sprunggelenkes gebildet wird. Es gibt einen Innen- und
Außenknöchel (jeweils körperfernes Ende von Schien bzw. Wadenbein). Bei einer
Verletzung durch Umknicken des Fußes nach außen (häufiger) bzw. innen kann
einerseits dieser Knochen-Fortsatz brechen, oder andererseits die Bandfixierung
dieser Strukturen zu den Mittelfußknochen reißen.

% \A{STEU: Knöchel erwähnen.}
% \A{HIRT: Review: ANA: Knöchel erwähnt, gut so?} 
% \A{GABS: Vl. könnte man Innen- und Aussenknöchel erwähnen.} 
% \A{HIRT: Review: ANA: gemacht!} 

 
\Fazit{Achtung: Unterscheide Fuß -- Bein!}

\subsection{Die Muskulatur}
 \index{Muskel}  
 \index{Muskulatur!glatte} 
 \index{Muskulatur!quergestreifte}
 \index{Muskulatur!Skelett-}
 \index{Muskulatur!Herz-}

\Layout{0}{\TxBeschreibung}
{
Prinzipiell gibt es 3  Arten von Muskelzellen:

\begin{itemize}
    \item \T{Quergestreifte (Skelett-)Muskulatur}:
    willkürlich steuerbar, trainierbar, er\-müd\-bar,  schnell und
    kräftig.     
    \item \T{Glatte  Muskulatur}: Nicht willentlich
    kontrollierbar, unermüdlich, ausdauernd
    \item \T{Herzmuskulatur}: Sonderstellung zwischen glatt und
    quergestreift, Eigenschaften von beiden: nicht willentlich
    steuerbar, trainierbar, nicht ermüdend, quergestreift.
    Verfügt über eigenen Impulsgeber (Reizleitungssystem).
\end{itemize}

\noindent Die \E{quergestreifte Muskulatur} heißt deswegen
{\quotedblbase}quergestreift{\textquotedblleft}, da man unter dem
Mikroskop tatsächlich eine Streifung sehen kann, allerdings nicht mit
dem freien Auge. 

Die glatte Muskulatur kommt u.a. in vielen
Verdauungsorgangen (z.\,B. Darm) und den Blutgefäßen vor.
}
{
\begin{itemize}
  \item Quergestreifte Muskulatur
  \item Glatte Muskulatur
  \item Herzmuskulatur
\end{itemize}
}
{}{}{}

\Layout{20}{Gegenspieler}
{%
Muskel können nur Zug und keinen
Druck ausüben, daher sind \Es{Gegenspieler} (Agonist u.
Antagonist) erforderlich (z.\,B. ein Muskel streckt den Arm,
ein \E{anderer} beugt ihn).

Jeder Muskel hat einen \Es{Ursprung} (\Speech{"`Wo kommt er her?"'})
und einen \Es{Ansatz} (\Speech{"`Wo geht er hin?"'}), diese beiden Punkte bestimmen
die Bewegung, die der Muskel machen kann. Der Ursprung eines Muskels ist meist
näher zum Körper, der Ansatz der Punkt, der weiter weg ist. Der Ursprung ist
definitionsgemäß außerdem immer der Punkt, der stillgehalten wird. Der Ansatz
dagegen wird durch den Muskelzug auf ihn zu bewegt. Für alle dieser Erklärungen
gibt es natürlich Ausnahmen (z.\,B. Unterstützung der Atemhilfsmuskulatur),
darauf soll aber hier nicht näher eingegangen werden.
}
{
\begin{itemize}
  \item Agonist - Antagonist
  \item Ursprung - Ansatz
\end{itemize}
}
{./images/anatomie/Gray361.png}{\AuthNotesPicLegalOk{}{}Muskelzug. }{Gray's Anatomy, Copyright expired.}






%%%%%%%%%%%%%%%%%%%%%%%%%%%%%%%%%%%%%%%%%%%%%%%%%%%%%%%%%%%
%%%%%%%%%%%%%%%%%%%%%%%%%%%%%%%%%%%%%%%%%%%%%%%%%%%%%%%%%%%
%%%%%%%%%%%%%%%%%%%%%%%%%%%%%%%%%%%%%%%%%%%%%%%%%%%%%%%%%%%
%%%%%%%%%%%%%%%%%%%%%%%%%%%%%%%%%%%%%%%%%%%%%%%%%%%%%%%%%%%
\clearpage
\section{Das Herz}
 \index{Herz} 
 \label{topic:herz}

\Layout{20}{\TxBeschreibung}
{%
Das Herz ist ein etwa faustgroßer\ZFootnote{Faustgröße des \E{Pateinten!}}
\Es{Hohlmuskel} und hat die Funktion einer
\Es{Pumpe}. Es liegt \Es{hinter dem Sternum\footnote{Brustbein}}
\A{\footnote{\A{EMHO: statt Fußzeile: "`Brustbein (Sternum)"'. damit
einheitlich.\par GABS: nein. Hier sollen sie bereits primär den Begriff
"`Sternum"' verwenden und daran gewöhnt werden.}}} und teilweise im linken
Brustkorb. Unten liegt es auf \Es{dem Zwerchfell auf}.

\index{Scheidewand}  
\index{Septum}  
\index{Vorhof}
\index{Atrium}  
\index{Herz!-kammer}
\index{Ventrikel}  
Das Herz ist durch eine \Es{Scheidewand} (Septum) in eine
\E{rechte} und eine \E{linke} Seite geteilt.
Umgangssprachlich spricht man dabei jeweils vom
\emph{{\quotedblbase}rechten{\textquotedblleft}} oder
\emph{{\quotedblbase}linken Herz{\textquotedblleft}}. Jede
Hälfte wird durch \Es{Klappen} in einen \T{Herzvorhof}
(\T{Atrium}) und eine \T{Herzkammer} (\T{Ventrikel})
geteilt.

\subparagraph{Blutversorgung}
\index{Herzkranzgefäße}
\index{Koronargefäße}

Die Blutversorgung des Herzens erfolgt über die
\T{Herzkranzgefäße} (\T{Koronargefäße}), welche direkt
aus der Aorta (Hauptschlagader) -- gleich oberhalb der Aortenklappe beim
Ausgang des linken Ventrikels -- abgehen.

\subparagraph{Schichten}

3 Schichten (von innen nach außen):
\AuthNotes{\footnote{GABS: Schichten relevant? KOCH: ja!}}
\begin{itemize}
\item Endokard (innen) \index{Endokard}     
\item \Es{Myokard} (\Es{Muskel}) \index{Myokard|Es} 
\item Epikard (aussen) \index{Epikard} 
\item Perikard (ganz aussen), bildet den Herzbeutel.
\end{itemize}

\noindent Im Herzbeutel befindet sich eine geringe Menge Flüssigkeit, die die
reibungslose Bewegung des Herzens ermöglicht. Wenn durch Eintritt von z.\,B. Blut nach einer Verletzung der
Herzbeutel zu viel gefüllt ist, kann es dadurch zu einer Einschränkung der
Herzbewegung und schlimmstenfalls zum Tod kommen.
}
{
\begin{itemize}
  \item Hohlmuskel
  \item Aufbau
  \begin{itemize}
    \item Scheidewand (Septum) - rechts / links
    \item Vorhof (Atrium)
    \item Kammer (Ventrikel)
    \item Herzklappen
  \end{itemize}
  \item Schichten
  \begin{itemize}
    \item Endokard
    \item Myokard
    \item Epikard
    \item Perikard
  \end{itemize}
\end{itemize}
}
{./images/hirtler-aass/Herz-edited-0800.png}{\AuthNotesPicLegalOk{}{}Das Herz.} {Hirtler}

~


\Layout{20}{Herzklappen}
{%
\index{Herzklappen}%
\index{Segelklappen}%
\index{Taschenklappen}%
%
Insgesamt \Es{4
Herzklappen} sorgen als Ventile dafür, dass das Blut nur in eine Richtung
strömt. Es gibt 2 Segelklappen und 2 Taschenklappen. Die Segelklappen befinden
sich zwischen den Vorhöfen und den Herzkammern, die Taschenklappen befinden
sich an den Ausflussöffnungen der Kammern am Übergang zur Lungenarterie (rechte
Seite) bzw. zur  Aorta (Hauptschlagader, linke Seite). 
}
{
\begin{itemize}
  \item 4 Stück
  \item 2 Segelklappen - Übergang zwischen Vorhof und Kammer
  \item 2 Taschenklappen - Übergang zwischen Kammer und Kreislauf
\end{itemize}
}
{./images/hirtler-aass/Herz-Schema_edited.pdf}
{Das Herz - schematisch.}
{Hirtler}







 
% \Captionof{figure}{ Schema der
% Herz-Anatomie.  }


\paragraph{Herzfunktion}\label{topic:herzfunktion}

Der Herzmuskel pumpt und entspannt sich beim Erwachsenen
in Ruhe ca. 60 bis 100~mal in der Minute. Bei jeder
Anspannung (Kontraktion) wirft das Herz ca. 70\,ml Blut in die Arterien
aus. Dabei pumpt es ca. 4\,-\,6\,l  Blut (bzw. 1\,x das
gesamte Blutvolumen) in der Minute.

Durch die Pumpfunktion des Herzens entsteht ein
lebenswichtiger Druck im (arteriellen) Gefäßsystem, der
(arterielle systolische) \E{Blutdruck} \index{Blutdruck}.
Vgl. \prettyref{topic:blutdruck}. 

%
%
\subsection{Reizleitungssystem}
 \index{Reizleitungssystem} 

\AuthNotes{GABS @ STEU: Reizleitungssystem -- wie
ausführlich? Prüfungsfrage?

KOCH: Beim Defi schon.

GABS: Defi-Teil noch nicht abgefasst!}

\Layout{22}{TxBeschreibung}
{%
Das Reizleitungssystem des Herzens ist ein eigenständiges  Erregungs- und 
Reizleitungssystem, das streng hierarchisch gegliedert ist. Jede dieser
Gliederungsebenen ist aber zur eigenständigen Erregungsbildung fähig, sofern
ihr nicht die übergeordnete Ebene zuvor kommt.

Es besteht aus:

\begin{itemize}
\item Sinusknoten  \index{Sinusknoten} 
({\sffamily 1}):  
{\quotedblbase}Schrittmacher{\textquotedblleft}, feuert mit einer
Eigenfrequenz von ca. 60/min. elektrische Impulse ab)
\item AV-(AtrioVentricular)-Knoten ({\sffamily 2}): Leitung
vom Vorhof in die Kammern, filtert zu hohe Frequenzen
\index{AV-Knoten} 
\item His-Bündel \index{His-Bündel} 
\item \Z{Tawara-} (Kammer-) Schenkel (linker/rechter)
\item \Z{Purkinje-}Fasern
({\quotedblbase}Endstrecke{\textquotedblleft})
\end{itemize}

\Fazit{Der normale Schrittmacher ist der Sinusknoten,
welcher die elektrischen Impulse für die Vorhöfe und die Kammern erzeugt. Man spricht dann
vom Sinusrhythmus.}

Beispiele von Rhythmen werden unter \prettyref{topic:herzrhythmusstoerungen}
vorgestellt. 

}
{
\begin{itemize}
  \item Sinusknoten
  \item AV-Knoten
  \item His-Bündel
  \item Tawara- (Kammer-) Schenkel
  \item Purkinje-Fasern
\end{itemize}
}{./images/anatomie/Reizleitungssystem_1.png}{\AuthNotesPicLegalOk{}{CC-BY}
Reizleitungssystem, Schema.}{J. Heuser,
basierend auf der Arbeit von Patrick J. Lynch; illustrator;
C. Carl Jaffe; MD; cardiologist Yale University Center for
Advanced Instructional Media. Lizenz: CC-BY}






\AuthNotes{KOCH: Tabelle mit Normalwerten einfügen

GABS: in Kap Vitalfunktionen}


\section{Der Kreislauf }
\index{Kreislauf}
\label{topic:blutkreislauf} 

\Fazit{\Es{Herz + Gefäßsystem + Blut = Kreislauf}}

\Layout{0}{TxBeschreibung}
{
Der Blutdruck und die Pumpfunktion des Herzens werden im
Kapitel \emph{Vitalfunktionen} unter \ref{topic:blutdruck}
(Blutdruck) sowie im Abschnitt \ref{topic:herzfunktion}
(Herzfunktion) in diesem Kapitel besprochen.

\index{Körperkreislauf}  \index{Kreislauf!großer} 
\index{Kreislauf!kleiner}
\index{Lungenkreislauf}
\index{Hochdrucksystem}  \index{Niederdrucksystem}
Es wird der \T{Körperkreislauf} (\E{"`Großer
Kreislauf"'}, "`Hochdrucksystem"') vom
\T{Lungenkreislauf} (\E{"`Kleiner Kreislauf"'},
"`Niederdrucksystem"') unterschieden. Das Herz ist für 
beide Systeme die gemeinsame Pumpe.

Der Kreislauf wird vom \Es{Stammhirn} gesteuert
(Herzfrequenz, Blutdruckregulation, etc.).
}
{
\begin{itemize}
  \item großer Kreislauf
  \item kleiner Kreislauf
  \item Herz = Pumpe
  \item Steuerung: Hirnstamm
\end{itemize}
}
{}{}{}

\subsection{Abschnitte des Kreislaufs}
\vspace{-1em} \index{Gasaustausch} 
\index{Kreislauf!Abschnitte}

\begin{gWideParEnv}
\Parallel
{\subsubsection{Kleiner (Lungen-)Kreislauf}
Die obere und untere Hohlvene münden in den \Es{rechten
Vorhof} des Herzens und liefern sauerstoffarmes Blut aus
dem Körper an. Von hier wird das Blut in die \Es{rechte
Herzkammer} gepumpt, um dann in die \Es{Lungenarterie} zu
gelangen. In der Lunge verästeln sich die Lungenarterien
zu haardünnen Gefäßen, den \Es{Kapillaren}. Diese
umspannen die Lungenbläschen es findet ein
\Es{Gasaustausch} statt: Sauerstoff wird aufgenommen und
Kohlendioxid abgegeben. Das jetzt mit Sauerstoff
angereicherte Blut wird über die \Es{Lungenvenen} in den
\Es{linken Vorhof} geleitet, der wiederum das Blut in die
\Es{linke Herzkammer} pumpt.

\vspace{1em}

\subsubsection{Großer Kreislauf}
Das Blut wird aus der linken Herzkammer ausgestoßen und
gelangt über die die Hauptschlagader
("`\Es{Aorta}"') in den
Körperkreislauf. Aus der Aorta entspringen \Es{Arterien}
({\quotedblbase}Schlagadern{\textquotedblleft}), die sich
immer weiter verzweigen und an Durchmesser abnehmen. Sie
versorgen den gesamten Körper mit Blut. Der durch die
Pumpleistung des Herzens erzeugte Puls ist an oberflächlich
gelegenen Arterien tastbar (Radialis-Puls, Carotis-Puls, ..).
Daher rührt auch der Name
"`Schlagader"'.

Sie verzweigen sich, weiter an Durchmesser abnehmend,
ähnlich dem Geäst eines Baumes, bis hin zu
\Es{Arteriolen} und haardünnen Gefäßen, den
Haargefäßen ("`\Es{Kapillaren}"').
Hier findet wieder ein \Es{Gasaustausch} statt, diesmal
allerdings spiegelbildlich zum Lungenkreislauf:
\Es{Sauerstoff wird an das Gewebe abgegeben, im Gegenzug
wird Kohlendioxid abtransportiert}.

Das nun sauerstoffarme Blut gelangt über kleine
\Es{Venolen} in das Venensystem und fließt in den
\Es{Venen} herzwärts. Am Weg zum Herzen nehmen die Venen
an Umfang immer mehr zu, bis sie letztendlich in die großen Venen
({\quotedblbase}\Es{obere} und \Es{untere
Hohlvene}{\textquotedblleft}, Vena cava \Z{superior und inferior}) münden. 
}
{
\includegraphics[width=\linewidth]{./images/hirtler-aass/Kreislauf-Schema.pdf}

\Captionof{figure}{\AuthNotesPicLegalOk{}{Blutkreislauf, Schema.}Der
Blutkreislauf.}

\SourcePicture{Hirtler}

\Fazit{Rechter Vorhof ${\rightarrow}$ rechte Herzkammer
${\rightarrow}$ Lungenarterie(n) ${\rightarrow}$
Lungenkapillaren (\ce{CO2}-Abgabe, \ce{O2}-Aufnahme) ${\rightarrow}$ Lungenvenen ${\rightarrow}$ linker Vorhof 

${\rightarrow}$ linke Herzkammer ${\rightarrow}$ Aorta ${\rightarrow}$
abzweigende Arterien ${\rightarrow}$ Kapillaren
(\ce{O2}-Abgabe, \ce{CO2}-Aufnahme)  ${\rightarrow}$ Venen
${\rightarrow}$ Hohlvene ${\rightarrow}$ Rechter Vorhof

 ${\rightarrow}$ alles beginnt wieder von vorne
 $\circlearrowright$}}
\end{gWideParEnv}

~




\subsection{Die Blutgefäße}
\index{Arterien}
\index{Venen}
\index{Kapillargefäße}

\Layout{0}{\TxBeschreibung}
{
Man unterscheidet zwischen:

\begin{labeling}[\dots{}]{Kapillargefäßexxx}
    \item[\T{Arterien}] führen \E{vom Herzen \Es{weg}} und haben eine dicke
    % GABS, 2010/05/04 Arterien und Muskelschicht: Erwähnung der dickeren
    % Muskelschicht ist wichtig, weil PRüfungsfrage.
    Wand und Muskelschicht. 
    \item[\T{Venen}] führen \E{zum Herzen \Es{hin}} und
    besitzen Venen\E{klappen}, um Rückfluss\footnote{Rückfluss:
    Fluss entgegen der Flussrichtung} zu vermeiden. 
    
    Venen
hingegen haben eine deutlich dünnere Wand als die
muskulösen Arterien.     
    \item[\T{Kapillargefäße}] sind dünne "`Haargefäße"',
    in denen der \E{\Es{Gasaustausch}} stattfindet; sie \E{verbinden Arterien
    und Venen}. 
\end{labeling}
}
{
\begin{itemize}
  \item Arterien
  \item Venen
  \item Kapillaren
\end{itemize}
}{}{}{}

\Layout{0}{Sonderfall Shunt}
{%
\label{topic:shunt}
\index{Dialyse-Patient!Shunt}
Ein \T{Shunt} \index{Shunt}  ist eine Verbindung zwischen Arterien und Venen. Er
kann künstlich (chirurgisch) angelegt, angeboren oder akzidentiell entstanden
sein. Bei \Es{Dialyse-Patienten}  wird oft künstlich -- mittels einer Operation
-- ein Shunt gelegt\ZFootnote{Dadurch dass durch die Vene umgeleitetes,
arterielles, Blut mit höherem Druck fließt, verdickt sich die Vene und man kann
sie mit dickeren Kanülen (wie sie für die Dialyse benötigt werden) punktieren.}.
Ein derartger Dialyse-Shunt ist \Ca{verletzlich}, daher darf an einem Arm mit
Shunt \Es{nicht gestaut} werden; das betrifft insbesonders die
\E{Blutdruckmessung} (vgl. \prettyref{topic:blutdruckmessung}). 
}{%
\begin{itemize}
  \item Künstliche Verbindung Arterie -- Vene
  \item Dialyse
  \item \Ca{Verletzlich}
  \item Nicht stauen!
\end{itemize}
}{}{}{}


\Layout{0}{Blutgefäße sind in Schichten aufgebaut}
{
Mit Außnahme der dünnsten Gefäße (Kapillaren) sind alle Gefäße annähernd gleich
aufgebaut. Von innen nach außen gibt es drei Schichten:

\Z{
\begin{itemize}
  \item Tunica intima: Grenze zwischen Blut und Blutgefäß, enthält elastische
  Fasern (bei Arterien dicker) und längsgerichtete Muskelfasern.
  \item Tunica media: besteht häuptsächlich aus ringförmigen Muskelfasern, sie
  ist bei Arterien dicker als bei Venen
  \item Tunica adventitia: lockeres Bindegewebe
\end{itemize}
}

Die größten Gefäße, wie zum Beispiel die Aorta, haben zur geeigneten
Blutversorgung ihrer Wand dort eigene Gefäße (Vasa vasorum). Arterien die
herznah sind, besitzen in ihrer Wand eher elastische Fasern. Arterien die der
Organversorgung dienen sind muskulärer aufgebaut.
}
{
\begin{itemize}
  \item 3 Schichten
  \item Unterschied Arterien - Venen
  \item elastische und muskuläre Arterien
\end{itemize}
}
{}{}{}


\subsubsection{Wichtige Blutgefäße}

\begin{gWideParEnv}

{\sffamily
\label{topic:blutgefaesse-wichtige}

% \begin{minipage}{\linewidth}
\noindent\Captionof{table}{Übersicht wichtiger Blutgefäße}

\noindent\begin{tabulary}{\linewidth}{lllL}
\\\addlinespace
\noindent\TabHeading{Gefäß}
&
\noindent\TabHeading{Fachbegriff}
&
\noindent\TabHeading{Typ}
&
\noindent\TabHeading{Bemerkungen}
\\\midrule\addlinespace
\Es{Hauptschlagader} \index{Hauptschlagader}
&
\T{Aorta} \index{Aorta}
&
Arterie
&
"`Hauptarterie"', geht von der linken Herzkammer ab, von dort zweigen alle
anderen Körperarterien ab.
% GABS: Aorta steht schon in der Fachbegriff-Spalte.
\\\addlinespace
\Es{Halsschlagader} \index{Halsschlagader}
&
\T{Karotis} \index{Karotis}
&
Arterie
&
1 Paar (2 Stk.), versorgen den Kopf und das Gehirn, Puls tastbar seitlich vom
Kehlkopf
\\\addlinespace
\Es{Herzkranzgefäße} \index{Herzkranzgefäße}
&
\T{Koronararterien} \index{Koronararterien}
&
Arterie
&
Versorgen das Herz mit Blut, gehen aus der Aorta direkt über ihrer Herzklappe ab
\\\addlinespace
\Es{Lungenarterie} \index{Lungenarterie}
&
\T{Pulmonalarterie} \index{Pulmonalarterie}
&
Arterie
&
\noindent\Es{\ce{O2}-arm.} Geht von der rechten Herzkammer weg in
die Lunge zu den Lungenkapillaren
\\\addlinespace
\Es{Oberarmarterie} \index{Oberarmarterie}
&
Brachialarterie
&
Arterie
&
Blutversorgung des Armes, Taststelle zum Puls-fühlen beim Baby, Blutdruckmessung
\\\addlinespace
\Es{Radialarterie} \index{Radialarterie}
&
\T{A. radialis} \index{A. radialis}
&
Arterie
&
Verläuft am Unterarm speichenseitig (dort wo der Daumen ist). Puls
fühlbar.
\\\addlinespace
\Es{Leistenarterie} \index{Leistenarterie}
&
\Z{A. femoralis} \index{A. femoralis}
&
Arterie
&
verläuft am Oberschenkel, Puls tastbar in der Leistenbeuge, wird für
Herzkatheter als Zugang verwendet. 
\\\addlinespace
\Es{Obere Hohlvene} \index{Obere Hohlvene}
\index{Hohlvene!obere}
&
\T{Vena cava} sup. \index{Vena cava}
&
Vene
&
Obere Vene die in den rechten Vorhof mündet
\\\addlinespace
\Es{Untere Hohlvene}
\index{Untere Hohlvene}
\index{Hohlvene!untere}
&
\T{Vena cava} inf. \index{Vena cava}
&
Vene
&
Untere Vene die in den rechten Vorhof mündet.
\\\addlinespace
\Es{Lungenvene} \index{Lungenvene}
&
\T{Pulmonalvene} \index{Pulmonalvene}
&
Vene
&
\Es{\noindent\ce{O2}-reich.} Kommt von den Lungenkapillaren und
mündet in den linken Vorhof
\\\addlinespace
\Es{Pfortader} \index{Pfortader}
&
\Z{Vena portae}
&
Vene
&
Führt \ce{O2}-armes, aber nährstoffreiches Blut von Teilen des Darms zur
Leber (zur Entgiftung)
\\\addlinespace
\Es{Haargefäße} \index{Haargefäße}
&
\T{Kapillargefäße} \index{Kapillargefäße}
&
Kapillare
&
Dünne, dicht verzweigte Gefäße, Verbindung zwischen Arterien
und Venen. Hier findet der Gas- und Nährstoffaustausch statt. Es gibt
Körperkapillaren (Körperkreislauf) und Lungenkapillaren
(Lungenkreislauf)
\\\addlinespace\bottomrule
\end{tabulary}
% \end{minipage}
}
% STEU: Leistenarterie! Puls
% GABS: Wozu?
\end{gWideParEnv}

\subsection{Das Blut}\index{Blut}\label{topic:blut}

\Layout{0}{\TxBeschreibung}
{%
Das Blut macht ca. \Es{\VonBis{7}{8}\,\% des Körpergewichtes}
aus, das entspricht in etwa \Es{\VonBis{5}{6}\,l}. 
} 
{
\begin{itemize}
  \item 5-6 Liter
  \item festen Bestandteilen (Zellen, Blutkörperchen)
  \item flüssigen Bestandteilen (\index{Plasma}Plasma)
  \item Aufgaben
\end{itemize}
}{}{}{}

\Layout{0}{Bestandteile}{%
Das Blut besteht aus
\Es{festen} Bestandteilen und \Es{flüssigen Bestandteilen} (Blutplasma). 

Zu
den festen Bestandteilen gehören man die \E{roten Blutkörperchen}
(Erythrozyten), die \E{weißen Blutkörperchen} (Leukozyten) und die
\E{Blutplättchen} (Thrombozyten). 

Das Blutplasma besteht zum größten Teil aus
\E{Wasser}. Außerdem enthält es \E{Proteine} und andere \E{gelöste Stoffe}.
}{
\begin{itemize}
  \item feste Bestandteile (Zellen, Blutkörperchen)
  \item flüssige Bestandteile (Plasma)
  \item \prettyref{t:blut-bestandteile}
\end{itemize}
}{}{}{}

\Layout{0}{Aufgaben}{%
Die wichtigste Aufgabe des Blutes ist der \Es{Stoff- und Gastransport}. Es
ist z.\,B. der \E{Sauerstofflieferant} des Körpers. Es holt sich den Sauerstoff
in der Lunge ab, bringt ihn zu den Zellen und nimmt auf dem Rückweg das
abzuatmende \E{Kohlendioxid} wieder mit. Weiters transportiert es die über den
Darm aufgenommenen \E{Nährstoffe} zu den entsprechenden
Weiterverarbeitungsstellen und sorgt für den Abtransport von \E{Abfallprodukten}
hin zu Leber oder Niere.

Eine weitere wichtige Aufgabe ist die \Es{Regulation der
Körpertemperatur} gemeinsam mit den Blutgefäßen. Bei Kälte werden die Blutgefäße
in den Extremitäten enger gestellt - weniger Wärme geht nach ausen verloren,
das Blut und dementsprechend die Wärme wird "`zentralisiert"'. Andersrum
funktioniert es bei großer Hitze. Über eine Weitstellung der Gefäße versucht der
Körper soviel Hitze wie möglich loszuwerden.

Zu guter Letzt sorgt das Blut auch noch für einen ausgeglichenen
\Es{Säure-Basen-Haushalt} im Körper -- der pH des Körpers sollte immer zwischen
7,35 und 7,45 liegen. Bei Abweichungen sorgen \E{Puffersysteme}
(\prettyref{topic:bikarbonatpuffer}) im Blut für eine Wiederherstellung dieses
pH-Wertes. 
}{
\begin{itemize}
  \item Sauerstofftransport zu den Körperzellen
  \index{Sauerstofftransport}
  \item \ce{CO2}-Abtransport \index{CO2-Abtransport@\ce{CO2}-Abtransport}
  \item Transport von Nährstoffen und Abfallprodukten
  \item Wärmeverteilung
  \item pH-Pufferungssystem (Säure/Basen-Gleichgewicht)
\end{itemize}
}{}{}{}

\begin{table}[H]
\begin{gWideParEnv}
\Captionof{table}{Übersicht: Die Bestandteile des
Blutes.\label{t:blut-bestandteile}}

\begin{tabularx}{\linewidth}{XXXXXX}
\hline\addlinespace
\multicolumn{6}{c}{\TabHeading{Blut}}
\\\addlinespace\cmidrule(lr){1-3}\cmidrule(lr){4-6}\addlinespace
\multicolumn{3}{c}{\TabSub{Feste Bestandteile}}
&
\multicolumn{3}{c}{\TabSub{Flüssige Bestandteile}}
\\
\multicolumn{3}{c}{$\sim$\,42\,\%}
&
\multicolumn{3}{c}{$\sim$\,58\,\%}
\\\addlinespace\cmidrule(lr){1-1}\cmidrule(lr){2-2}\cmidrule(lr){3-3}\cmidrule(lr){4-4}\cmidrule(lr){5-5}\cmidrule(lr){6-6}\addlinespace
\TabSub{Erythrozyten}
&
\TabSub{Leukozyten}
&
\TabSub{Thrombozyten}
&
\TabSub{Wasser}
&
\TabSub{Proteine}
&
\TabSub{Gelöste Stoffe}
\\
\noindent\Z{($\sim$\,4\,000\,000\,/$\mu$l)}
&
\noindent\Z{($\sim$\,4\,000\,--\,8\,000\,/$\mu$l)}
&
\noindent\Z{($\sim$\,150\,000\,--\,300\,000\,/$\mu$l)}
&
\noindent$\sim$\,90\,\%
&
\noindent$\sim$\,8\,\%
&
\noindent$\sim$\,2\,\%

{\scriptsize%
Ionen,

Glukose, 

Hormone,

Kreatinin, 

Harnstoff,

{\ldots}%
}
\\\addlinespace\hline
\end{tabularx}

\end{gWideParEnv}

\end{table}

% \noindent\begin{minipage}{\linewidth}
% %  \includegraphics[width=\linewidth]{./images/noauth/AASdev20090228-img32.png} 
% \includegraphics[width=\linewidth]{./icons/under-construction.png}
%  \TODOCRITICAL{GABS}{2010-04-11}{Bild: neu erstellen: Bestandteile des
%  Blutes.}{}
%  
%  
% \end{minipage}





\subsubsection{Feste (zelluläre) Bestandteile}

Die festen Bestandteile des Blutes beinhalten rote Blutkörperchen
(Erythrozyten), weiße Blutkörperchen (Leukozyten) und Blutplättchen
(Thrombozyten).




\Layout{0}{Erythrozyten}
{%
\label{topic:erythrozyten}%
\index{Erythrozyten}%
\index{Rote Blutkörperchen}%
\index{Blutkörperchen!rote}%
\index{Blutkonserve!Erythrozytenkonzentrat}
\index{Erythrozytenkonzentrat}
Die roten Blutkörperchen sind quasi die wichtigsten des ganzen Körpers - ohne
ihren Sauerstofftransport funktioniert im Körper gar nichts. Das Protein, das
für die Bindung des Sauerstoffes an die Erythrozyten verantwortlich ist, heißt
\E{Hämoglobin}. Es ist auch der \E{Blutfarbstoff}.

Die Erythrozyten sind außerdem die häufigsten festen Blutbestandteile
(4\,000\,000\,/$\mu$l) und tragen die \Es{Blutgruppenmerkmale} (A, B bzw. 0 bei
Fehlen der Merkmale) und die Rhesusfaktoren. Es gibt sog.
\T{Erythrozytenkonzentrate} ("`Ery-Konzentrate"'), das sind \Es{Blutkonserven}, welche fast ausschließlich Erythrozyten enthalten.

\Z{Das Kennzeichen der Erythrozyten ist die Abwesenheit des Kernes\ZFootnote{Der
Kern der Erythrozyten geht bei der Reifung verloren.} bzw. jeglicher
Zellorganellen. Nach einer Lebensdauer von ca. 100 Tagen werden die "`Senioren"' der roten
Blutkörperchen in der Milz abgebaut.}
}
{
\begin{itemize}
    \item Blutfarbstoff \T{Hämoglobin}\index{Hämoglobin} (Transportprotein für
    Sauerstoff)
    \item \ce{O2}-Transport
    \item Tragen Blutgruppenmerkmale (Blutgruppe)
    \item Erythrozyten-Konzentrate: Blutkonserven
    \item Häufigsten Blutkörperchen \Z{(4\,000\,000\,/$\mu$l)}
\end{itemize}
}{}{}{}



\Layout{0}{Leukozyten}
{%
\label{topic:leukozyten}% 
\index{Leukozyten}% 
\index{Weiße Blutkörperchen}%
\index{Blutkörperchen!weiße}%
\index{Immunsystem!Leukozyten}
Die weißen Blutkörperchen sind ein wichtiger \Es{Teil des Immunsystems} und
dienen der Infekt- und Fremdkörperabwehr des Körpers. 

\Z{Es befinden sich
zwischen 4\,000 und 8\,000 Leukozyten pro $\mu$l im Blut des Körpers; es gibt
verschiedene Arten von Leukozyten\ZFootnote{\begin{inparaitem}
  \item \E{Granulozyten}: Diese Zellen bilden den Beginn der Immunabwehr. Es
  gibt je nach Erkrankung verschiedene "`Spezialisten"' unter den Granulozyten -
  neutrophile G. (bei Keimen und Fremdkörpern), eosinophile G. (bei Allergien
  und Parasiten) und basophile G.
  \item \E{Lymphozyten}: in diese Gruppe fallen Antikörper-bildende Zellen.
  \item \E{Monozyten}: können sich durch zelleigene Bewegungen im Gewebe bewegen
  und Keime bzw. Fremdkörper verdauen.
\end{inparaitem}}.
}
}
{
\begin{itemize}
    \item Körperabwehr, Entzündungszellen, Immunsystem
    \item \Z{(4\,000\,--\,8\,000\,/$\mu$l)}
\end{itemize}
}{}{}{}


\Layout{0}{Thrombozyten}
{\label{topic:thrombozyten}
\index{Thrombozyten}%
\index{Thrombozytenkonzentrat}

Die Blutplättchen sind kleine kernlose Scheibchen. Sie dienen der
\Es{Blutgerinnung} und dem Wundverschluss. \Z{Ihre Anzahl wird mit 150\,000 bis
300\,000 pro $\mu$l beziffert.}

Es \T{Thrombozytenkonzentrate}, das sind
\Es{Blutkonserven}, welche fast ausschließlich Thrombozyten beinhalten. 
}
{
\begin{itemize}
    \item Wichtige Rolle beim Wundverschluss
    \item Zelluläre Blutstillung
    \item \Z{(150\,000\,--\,300\,000\,/$\mu$l)}
    \item Thrombozytenkonzentrat: Blutkonserve
\end{itemize}
}{}{}{}

\subsubsection{Die flüssigen Bestandteile --
Das Plasma}
\index{Plasma}

\Layout{0}{\TxBeschreibung}
{%
Das Blutplasma besteht zum größten Teil aus Wasser. Der restliche Anteil
enthält Proteine, Enzyme, Hormone, Elektrolyte, Glukose und Fette.

Aus den Inhaltsstoffen ersichtlich besteht die Aufgabe des Plasmas aus dem
Transport von Nährstoffen an den Zielort. Weiters sorgt es für die
Wärmeverteilung im Körper.

An wichtigen Proteinen verteilt das Plasma die Bestandteile für die
plasmatische Blutgerinnung und das Immunsystem (Antikörper, Immunglobuline).
Weiters sorgt es über die Elektrolytverteilung für den ausgeglichenen pH-Wert
im menschlichen Körper.
}
{
\begin{itemize}
    \item 10\% Proteine, Enzyme, Hormone, Elektrolyte, Glukose, Fette
    \item 90\% Wasser
    \item Aufgaben:
    \begin{itemize}
        \item Nährstofftransport
        \item Wärmeverteilung
    	\item Proteine und Enzyme für \Es{Plasmatische Blutgerinnung} (Serum)
    	\item Pufferung des pH-Wertes
    \end{itemize}
\end{itemize}
}{}{}{}

\subsection{Die Blutstillung}
 \index{Blutstillung} 

\Layout{60}{\TxBeschreibung}{%
Die Blutstillung umfasst alle Vorgänge, die zwischen dem Entstehen und dem
Verschluss einer Wunde ablaufen. Sie läuft in zwei Phasen ab: vorläufiger
Wundverschluss durch Bildung eines \E{Thrombozytenpfropfs} 
\index{Thrombozytenpfropf} und endgültiger Wundverschluss
durch Blutgerinnung\index{Blutgerinnung}. 
}
{
\begin{itemize}
  \item vorläufiger Wundverschluss (Thrombozytenpfropf)
  \item endgültiger Wundverschluss (Blutgerinnung)
\end{itemize}
}
{./images/anatomie/Bleeding_finger.jpg}
{\AuthNotesPicLegalOk{}{CC-BY}}
{\SourcePicture{Crystal
(Crystl) from Bloomington, USA
[\url{http://www.flickr.com/people/crystalflickr/}]
Lizenz: CC-BY}}



\paragraph{Bildung eines Thrombozytenpfropfs}

Zunächst
lagern sich an die verletzte und defekte Stelle des Blutgefäßes Thrombozyten an
und verkleben miteinander. Die verklebenden
Blutplättchen bilden einen Thrombozytenpfropf.

\Fazit{Die Bildung des Thrombozytenpfropf wird von den zelluären Bestandteilen
des Blutes erledigt.}



\Layout{0}{Blutgerinnung}
{%
\Z{Die Blutgerinnung ist der wichtigste
Prozess bei der Blutstillung. Sie wird entweder durch die
Gewebeverletzung selber oder durch ein eigenes System von Gerinnungsfaktoren
ausgelöst. Der Vorgang der Blutgerinnung hat zum Ziel,
das im Blutplasma gelöste \Z{Fibrinogen} zu aktivieren und
in einen gallertartigen, klebrigen und später festen Zustand
zu überführen. Um dieses Ziel zu erreichen, steht dem Organismus ein
Gerinnungsystem bestehend aus im Plasma gelösten Gerinnungsfaktoren zur
Verfügung, das kaskadenartig aufgebaut ist. Dieses System funktioniert wie
eine Kette aufgestellter Dominosteine, die nacheinander umfallen, wenn der
erste Stein angestoßen wird. }

\Fazit{Die Blutgerinnung kommt durch im Plasma gelöste Gerinnungsfaktoren
\label{topic:gerinnungsfaktoren} zustande.}

Neben dem System der Blutgerinnung
gibt es auch ein System zur Auflösung von Thromben, die Fibrinolyse. Beide
Systeme müssen sich die Waage halten.

\Fazit{An der Blutstillung sind sowohl die festen Bestandteile des Blutes als
auch das Plasma beteiligt.}
}
{
\begin{itemize}
  \item Thrombozyten
  \item Fibrinogen
  \item Gerinnungsfaktoren
  \item Fibrinolyse
\end{itemize}
}
{}{}{}

\clearpage

\section{Das Nervensystem}
 \index{Nervensystem} 


\Layout{0}{Unterteilung}
{
Man kann das Nervensystem nach seinem Aufbau und nach seiner Funktion
unterteilen.

Vom Blickwinkel des \underline{Aufbaus} wird das Nervensystem grob
unterteilt in ein

\begin{itemize}
\item \T{Zentrales Nervensystem} (\T{ZNS}) und ein 
\item \T{Peripheres Nervensystem}.
\end{itemize}

\noindent Das ZNS ist im Wesentlichen für \E{Steuer-, Regel- und Denkfunktionen}
zuständig, sowie teilweise für die Weiterleitung von Impulsen und
Informationen. Zum ZNS gehören \Es{Gehirn} und \Es{Rückenmark}. 

Das periphere
Nervensystem ist fast ausschließlich für die \E{Weiterleitung von Informationen}
und Impulsen zu den Endorganen (oder umgekehrt, von den Endorganen zum ZNS)
zuständig. Zum peripheren Nervensystem gehören alle Nerven.

Betrachtet man die \underline{Funktion} kann man das
Nervensystem in ein

\begin{itemize}
\item \T{Motorisches} (willkürliche\footnote{\emph{Willkür}: Dem \E{freien
Willen} unterworfen.} Bewegung),
\item \T{Sensorisches} ("`fühlend"', Reizweiterleitung von
den Sinnesorganen) und ein
\item \T{Vegetatives} (autonomes, bzw. unwillkürliches)
Nervensystem
\end{itemize}

unterteilen. 
}
{
Unterteilung nach Aufbau
\begin{itemize}
  \item zentrales Nervensystem
  \item peripheres Nervensystem
\end{itemize}
Unterteilung nach Funktion
\begin{itemize}
  \item motorisches Nervensystem
  \item sensorisches Nervensystem
  \item vegetatives Nervensystem
\end{itemize}
Die Begriffe sind recht unscharf
definiert und geben nur eine grobe, aber wichtige Übersicht
wieder.
}
{}{}{}

\subsection{Das Zentralnervensystem -- ZNS}
% \TagNo{Unterteilung nach Lage und Aufbau}
 \index{Zentralnervensystem}  
 \index{ZNS} 
\index{Nervensystem!Zentral-}
\label{topic:zns}


\Es{Gehirn} und \Es{Rückenmark} sind morphologisch und funktionell eine
Einheit. Gemeinsam verwerten und verarbeiten sie die Signale aus der Peripherie
und wählen diejenigen aus, die entweder unterbewusst verarbeitet oder bewusst
wahrgenommen werden sollen. Das Zentralnervensystem erzeugt schließlich
\Es{Signale}, die -- oft als Antwort auf einen Reiz -- in die
Peripherie abgegeben werden.

\Layout{20}{Unterteilung nach Lage und Aufbau}
{%
Das Zentralnervensystem (ZNS) besteht aus:

\begin{itemize}
    \item \Es{Gehirn}
    \begin{itemize}
        \item Großhirn
        \item \Z{Zwischenhirn}
        \item \Z{Kleinhirn}
        \item Hirnstamm
    \end{itemize}
    \item \Es{Rückenmark}
\end{itemize}

% Der Balken ist die Verbindung zwischen den beiden Großhirnhälften.

\A{STEU: Hirnteile, Balken in grau}
}
{
\begin{itemize}
    \item Gehirn
    \begin{itemize}
        \item Großhirn
        \item \Z{Zwischenhirn}
        \item \Z{Kleinhirn}
        \item Hirnstamm
    \end{itemize}
    \item Rückenmark
\end{itemize}
}
{./images/anatomie/Brain_human_sagittal_section.png}
{Das Hirn im seitlichen Querschnitt.}
{Patrick J. Lynch, medical illustrator,
\url{http://patricklynch.net}; C. Carl Jaffe, MD,
cardiologist. Lizenz: CC-BY 2.5}


\paragraph{Aufgaben allgemein}~

\begin{itemize}
    \item Aufnahme  und Verarbeitung  von äußeren und internen  Reizen
    (Schmerz, Fühlen, Sehen, Hören, {\dots})
    \item Bewusstsein und {\quotedblbase}Bewusstmachung{\textquotedblleft}
    \item Abgabe  von zweckentsprechenden Impulsen (Steuerung, Reaktion
    etc.)
\end{itemize}



\paragraph{Spezielle Aufgaben und Aufbau}~

\index{Großhirn}
\index{Zwischenhirn}
\index{Kleinhirn}
\index{Hirnstamm}
\index{verlängertes Rückenmark}
\index{Rückenmark}

\Layout{0}{Großhirn}
{
Das Großhirn ist verantwortlich für Hören, Sehen, Sprache sowie Riechen.
Außerdem sind dort Zentren für Gefühlsregungen und Gedächnis. Auch bewusste
Bewegungen und die Sensorik werden im Großhirn verarbeitet.

\Z{Es besteht aus 2 Hälften, sie stehen über
        den Balken miteinander in Verbindung. An der Oberfläche sind Furchen
        und Windungen zu erkennen. Das Großhirn ist in Lappen gegliedert.}
}
{
\begin{itemize}
    \item Bewusstes Denken
    \item Sitz des Bewusstseins
    \item Gefühle
    \item Gedächtnis
\end{itemize}
}
{}{}{}

\Layout{0}{\Z{Zwischenhirn}}
{
\Z{Im  Zwischenhirn laufen alle Signale aus dem Körper zusammen, hier befinden
sich wichtige \E{Regulationseinheiten}\ZFootnote{Der wichtigste
Kern, der Thalamus, ist das wichtigste Integrations-, Koordinations- und
Modulationszentrum sowohl für Sensorik als auch für Motorik. Über den
Hypothalamus, ein weiterer Kern des Zwischenhirns, werden alle vegetativen
Funkionen des Körpers geregelt.}. Es ist eine \E{Schaltinstanz} und ein
\E{Filter} zwischen Rückenmark und Großhirn, auch die \E{Regulation der
vegetativen Funktionen} findet hier z.\,T. statt.} 
}
{
\Z{
\begin{itemize}
\item Schaltinstanz und Filter Rückenmark - Großhirn
\item Regulation der vegetativen Funktionen
\end{itemize}}
}{}{}{}

\Layout{0}{\Z{Kleinhirn}}
{\Z{
Das Kleinhirn dient der Koordination- und Regulation der Motorik. Zusammen mit
dem Labyrinth des Innenohres sorgt es für das Körpergleichgewicht. Die
Bewerkstelligung von komplexen Bewegungen hat auch hier ihren Ursprung.

\Z{Das Kleinhirn hat zwei Hälften, sowie feine Windungen und Furchen auf der
        Oberfläche.}
}
}
{\Z{
\begin{itemize}
    \item Bewegungskoordination ({\quotedblbase}Makros{\textquotedblleft})
    \item Gleichgewicht
\end{itemize}
}
}{}{}{}

\Layout{0}{Hirnstamm und verlängertes Rückenmark}
{
Der Hirnstamm ist das Regulationszentrum aller lebenswichtigen Funktionen
(Atmung, Blutdruck, Schlaf usw.). Außerdem ist es die Ursprungsstelle (bzw. in
manchen Fällen die Ankunftstelle) der Hirnnerven.
}
{
\begin{itemize}
    \item Atem-, Kreislaufzentrum
    \item Hirnnerven
\end{itemize}
}{}{}{}

\Layout{0}{Rückenmark}
{
Das Rückenmark bildet die \Es{Verbindung zwischen Gehirn und Peripherie}.

Funktionell steht es unter der Kontrolle übergeordneter Zentren. Es können über
einen Eigenapparat eingehende Signale selbstständig verarbeitet werden
(hierunter versteht man Eigen- und Fremdreflexe). Diese Eigenschaft hat Einfluß
auf das Verhalten wenn übergeordnete Hirnzentren ausfallen, z.\,B. nach einem
Schlaganfall oder einer Querschnittsverletzung: Nach einiger Zeit (Tage
bis Wochen) übernimmt dann das Rückenmark die Steuerung der Muskelspannung
(Entwicklung einer \E{Spastik}); eine durch den Menschen bewusst gesteuerte
Bewegung ist allerdings dadurch nicht möglich.

Das Rückenmark lässt sich durch abgehende Spinalnerven in \E{Segmente} (insg.
31) einteilen. Es gibt 8 Halssegmente, 12 Thorakalsegmente, 5 Lumbalsegmente und 5
Sakralsegmente. Das letzte Segment endet auf Höhe des 1. bis 2. Lumbalwirbel.
}
{
\begin{itemize}
    \item Weiterleitung von Impulsen (Leitungsbahnen)
    \item Vom Rückenmark gehen die peripheren Nerven aus
    \item (Muskeleigen-)Reflexe
    \item 31 Segmente, Ende auf Höhe L1-L2
\end{itemize}
}{}{}{}

\Layout{22}{Exkurs: Der Kopf besteht aus Schichten}
{
%\index{Epiduralraum}
%\index{Dura mater}
\index{Harte Hirnhaut}
%\index{Subduralraum}
%\index{Arachnoidea}
%\index{Subarachnoidalraum}

Unter dem Haupthaar und der damit verbundenen Haut folgt die Kopfschwarte. Die
Besonderheit der Kopfschwarte besteht darin, dass Blutgefäße nicht wie sonst im
Körper bei Verletzung kolabieren. Dadurch kommt es \E{auch bei relativ geringen
Kopfverletzungen} zu vergleichsweise \E{starken Blutungen}.

Hierauf folgt die Schädelkalotte. \Z{Diese ist innen fest mit der harten
Hirnhaut (Dura mater) verbunden - sie ersetzt hier das Periost des Knochens. In die Dura
sind die großen venösen Abflussgefäße des Gehirns eingebaut.

Auf die harte Hirnhaut folgt die Spinngewebshaut (Arachnoidea mater). In sie
sind ebenfalls Gefäße eingebettet.

Darunter folgt die weiche Hirnhaut (Pia mater). Sie liegt den Furchen und
Windungen des Gehirns untrennbar an.}
}
{
\begin{itemize}
    \item Kopfschwarte
    \item Knochen
%    \item Epiduralraum
    \item Harte Hirnhaut \Z{(Dura mater)}
%    \item Subduralraum
    \item Arachnoidea mater (Spinngewebshaut)
%    \item Subarachnoidalraum
%    \item eingebettete Gefäße
    \item Weiche Hirnhaut \Z{(Pia mater)}
\end{itemize}
}
{./images/anatomie/Gray1196-edited.png}
{\AuthNotesPicLegalOk{}{PD-CE}Querschnitt durch
die Schädeldecke. }
{Gray's Anatomy, Copyright expired.}


% Aufbau ZNS
% 
% \begin{itemize}
% \item Großhirn (Telencephalon)
% 
% \end{itemize}
% \begin{itemize}
% \item \index{Graue Substanz}Graue Substanz (Rinde)
% 
% \item \index{Weiße Substanz}Weiße Substanz (Mark)
% 
% \item Balken (Verbindung li/re)
% 
% \end{itemize}
% \begin{itemize}
% \item Zwischenhirn  (Diencephalon)
% 
% \end{itemize}
% \begin{itemize}
% \item Thalamus, Hypothalamus (Hormone, Temperaturregulation etc.)
% 
% \item Hypophyse (Hormone)
% 
% \end{itemize}
% \begin{itemize}
% \item Hirnstamm (Truncus encephali)
% 
% \item Mittelhirn (Mesencephalon)
% 
% \item Brücke (Pons)
% 
% \item verlängertes Mark (Medulla oblongata)
% 
% \item Ventrikelsystem
% 
% \end{itemize}
% \begin{itemize}
% \item Liquorgefüllte Hohlräume (Ventrikel)
% 
% \item haben nichts mit den Vnetrikel im Herz zu tun!
% 
% \end{itemize}
% Liquor = Gehirnwasser, in Ventrikeln gebildet

\noindent Das Hirn und das Rückenmark schwimmt in einer
Flüssigkeit, dem \T{Liquor}. Die
Schädelknochen und der Liquor bieten eine Schutzfunktion. 
Wenn Liquor austritt (z.\,B. durch Nase, Ohr, ...)  ist das
ein Zeichen für eine schwerwiegende Verletzung!

\Fazit{Tritt Liquor in die Umwelt aus ist das ein Zeichen
für eine schwerwiegende Verletzung.} 

\AuthNotes{KOCH: doppelt? GABS: ja. Schlüsselsätze sollen
doppelt (normales Format + Fazit-Box) vorkommen}

% \paragraph{\Z{Blut-Hirn-Schranke}} 
% 
% \Z{Die Blutgefäßwand ist im Hirn besonders dicht. Sie
% filtert daher besonders gut und nur für bestimmte Moleküle (\ce{O2},
% \ce{CO2}, Wasser, \ldots) leicht passierbar. Sie bietet dem Hirn
% Schutz vor schädlichen und giftigen Substanzen. 
% 
% Die Blut-Hirn-Schranke bietet dem Hirn besonderen
% Schutz vor schädlichen Substanzen im Blut.} 

\subsection{Das Periphere Nervensystem -- PNS}
\index{Nervensystem!Peripheres}
\index{PNS}
% \TagNo{Unterteilung nach Lage und Aufbau}

\Definition{Unter dem Begriff \T{Peripheres Nervensystem} fasst man
die Verbindungsnerven zwischen dem ZNS und den Endorganen zusammen.}

\Layout{20}{Unterteilung nach Lage und Aufbau}
{%
Das periphere Nervensystem dient der Kommunikation zwischen zentralem
Nervensystem und den peripheren Organen. Es werden dabei Informtaionen in
\Es{beide Richtungen} übertragen, d.\,h. vom ZNS zu den Organen
(\Speech{"`Hirn an Hand: Leg' dich auf die Herdplatte"'}) und von den Organen
zum ZNS (\Speech{"`Haut von Hand an Hirn: Hier ist es heiß!"'}).

Es werden Hirnnerven und Spinalnerven unterschieden. Hirnnerven (insgesamt 12
Paare) sind Nerven, die ihren Ursprung im Hirnstamm haben. Sie sind
grundsätzlich für die Versorgung vom Kopf zuständig. 

Die \Es{meisten Nerven entspringen segmentweise dem Rückenmark}
und werden Spinalnerven genannt. Sie verbinden das Rückenmark mit den
jeweiligen Organen. \Z{Für die Extremitäten bilden sie jeweils Nervengeflechte,
sogenannte Plexus (Plexus cervicalis, Plexus brachialis, Plexus lumbalis und
Plexus sacralis).}

Die Fortleitung innerhalb eines Nerven funktioniert über \E{elektrische
Impulse}. \Z{Zwischen Nerven und anderen Nerven bzw. den Zielorganen dienen
chemische Signale der Weiterleitung (in sog. Synapsen).}
}
{
\begin{itemize}
    \item Kommunikation ZNS ${\longleftrightarrow}$
    periphere  Organe
    \item Hirnnerven %(12)
    \item Spinalnerven (gehen vom Rückenmark segmental ab)
    \item elektrische Impulse, chemische Signale (Synapsen)
\end{itemize}
}
{./images//hirtler-aass/Rueckenmark.pdf}
{\AuthNotesPicLegalOk{}{}Rückenmark
mit abgehenden peripheren Nerven. Beachte dass die Nerven
geordnet nach Rückenmarkssegmenten abgehen.}
{Hirtler.}



\Layout{0}{Arten von Nerven}
{%
\index{Nerven!motorische}
\index{Nerven!sensorische}
\index{Nerven!vegetative} 
\T{Motorische Nerven} sorgen dafür, dass ein "`Bewegungswunsch"' in den
entsprechenden Muskeln ankommt. Auf das Nervensignal hin kommt es zur
Kontraktion des Muskels.

\T{Sensorische Nerven} liefern dem Gehirn \Es{Sinneseindrücke}. Diese können von
den diversen Sinnesorganen kommen. So gehören Tastempfindungen,
Temperaturempfindungen, Gerüche, Geschmacksempfindungen usw. dazu. Aber auch
weniger offensichtliche Dinge, wie z.B. die Gelenksstellung gehören dazu.

\T{Vegetative Nerven} sind sozusagen die Drahtzieher im Hintergrund. Sie sorgen
für den reibungslosen \Es{Ablauf unbewusster Funktionen}, wie z.\,B. die
Herzfunktion, die Verdauung und die Nierenfunktion. Die vegetativen Nerven können je nach ihrer
Funktion - global stimulierend oder hemmend -  in zwei Systeme unterteilt
werden: In den \Es{Sympathikus} und \Es{Parasympathikus}.
}
{
\begin{itemize}
    \item motorisch  
    \item sensorisch 
    \item vegetativ 
    \begin{itemize}
        \item sympathisch
        \item parasympathisch
    \end{itemize}
\end{itemize}
}{}{}{}

\subsection{Das vegetative (autonome)  Nervensystem}
\index{Nervensystem!Vegetatives} 
\index{Nervensystem!autonomes}
\index{Vegetatives Nervensystem}
\label{topic:vegetatives-nervensystem}
\label{topic:nervensystem-vegetatives}


 
\paragraph{Sympathikus  und Parasympathikus  sind Gegenspieler}~
\index{Sympathikus}  
\index{Parasympathikus}
\label{topic:sympathikus}
\label{topic:parasymapthikus} 
%  \TagNo{Unterteilung nach Funktion}

\vspace{\baselineskip}

\noindent\Parallel{
\T{Sympathikus}

Aktiviert den Aufmerksamkeits- und Kampfmodus.

\emph{"`Fight or flight"'} -- \emph{"`Kampf und Flucht"'}
}{%
\T{Parasympathikus} 

Aktiviert den Ruhemodus.

\emph{"`Rest and digest"'} - \emph{"`Rast und Verdauung"'}
}

z.\,B.: RR + HF

\noindent%
\Parallel{
Sympathikus:  HF ${\uparrow}$, RR ${\uparrow}$
}{
Parasympathikus: HF ${\downarrow}$, RR ${\downarrow}$
}




\begin{figure}[H]
    \begin{center}
    \includegraphics[width=\linewidth]{./images/hirtler-aass/Parasymp-Symp-Cartoon.pdf}
    \Captionof{figure}{"`Fight or flight, rest and digest"'
    \SourcePicture{Hirtler.}}
    \end{center}
\end{figure}

\opt{STATUS-EXPERIMENTAL}{
\begin{figure}[H]
    \begin{center}
   
%     \includegraphics[width=\linewidth]{./images/noauth/AASdev20090228-img37.png} 
\includegraphics[width=\linewidth]{./icons/under-construction.png}

\TODO{GABS}{2010-04-11}{Bild: neu erstellen: Übersicht vegetatives
Nervensystem.}{}

    \Captionof{figure}{\AuthNotesPicLegalNotOk{img37}{}
    Nur zur Illustration: Das vegetative Nervensystem hat
    Einfluss auf viele Organe. Dieses Bild
    \underline{nicht lernen}. Es dient nur der
    Verdeutlichung wie wichtig und einflussreich das vegetative
    Nervensystem im Körper ist. \SourcePicture{Quelle nicht eruierbar.}} 
    \end{center}
\end{figure}
}
 
 
 
% \clearpage
\section{Das Auge}
\index{Auge}
\label{topic:auge} 

\TODO{GABS}{2010-08-05}{Text fehlt}{}

{\TakeNotesLarge}


% \clearpage
\section{Das Ohr}
\index{Ohr}
\label{topic:ohr} 

\TODO{GABS}{2010-08-05}{Text fehlt}{}

{\TakeNotesLarge}


 
\clearpage
\section{Die Haut}
\index{Haut}
\label{topic:haut} 


Die Haut ist die Grenzfläche des Körpers zur
Außenwelt und gleichzeitig das größte Sinnesorgan des menschlichen Körpers.
\Z{Sie macht ca. \nicefrac{1}{6} des Körpergewichts aus und ihre Oberfläche
beträgt beim Erwachsenen ca. 1,6\,m{\texttwosuperior}.}



\Layout{0}{Aufgaben}
{
Die Haut ist \E{die} Barriere gegen Einflüsse aus der Umwelt. Sie verhindert
Eindringen von Fremdkörpern und Keimen. Sie sorgt für einen ausgeglichenen
Wasser- und Elektrolythaushalt, indem sie den Körper vor Austrocknung schützt.

Durch Absonderung von Schweiß bzw. der Regelung der Hautdurchblutung sowie
durch ihre Isolationsfähigkeit sorgt die Haut für eine gleichbleibende
Körperkerntemperatur innerhalb des Körpers bei wechselnden
Witterungsverhältnissen.

Als größtes Sinnesorgan des Körpers dient die Haut der Tastempfindung sowie der
Temperaturempfindung.
}
{
\begin{itemize}
    \item Schutz vor äußeren Einflüssen 
    \item Schutz vor Wasserverlust/Austrocknung
    \item Temperaturregulation 
    \begin{itemize}
        \item Isolation
        \item Wärmeabgabe
\end{itemize}
    \item Sekretionsorgan (Schweiß, Talg)
    \item Sinnesorgan (Tastsinn, Temperatursinn)
\end{itemize}
}{}{}{}

\Layout{0}{Die Haut ist in Schichten aufgebaut}
{\index{Oberhaut}\index{Lederhaut}\index{Unterhaut}%
\begin{minipage}{\linewidth}
\begin{center}
\includegraphics[width=0.48\linewidth]{./images/hirtler-aass/Haut-edited.png}
\end{center}

\CaptionofDiff{figure}{Die Haut im Querschnitt: \TabSub{A:} Oberhaut \Z{Epidermis},
\TabSub{B:} Lederhaut \Z{Dermis},
\TabSub{C:} Unterhaut (Subkutis),
\TabSub{d:} \Z{Meißner'sches Tastkörperchen},
\TabSub{e:} Talgdrüse,
\TabSub{f:} Haar mit Haarfollikel,
\TabSub{g:} Schweißdrüse,
\TabSub{h:} Venole,
\TabSub{i:} Freie Nervenenden (Schmerzempfindung),
\TabSub{j:} Arteriole,
\TabSub{k:} \Z{Pacini-Körperchen}.

\SourcePicture{Hirtler.}}{Die Haut im Querschnitt \SourcePicture{Hirtler.}}
\end{minipage}

Die
oberflächlichste Schicht der Haut ist die \T{Oberhaut} (Epidermis). Sie besteht aus einer mehr oder weniger dicken \E{Hornhaut} (am Dicksten an
Handflächen und Fußsohlen, am Dünnsten an den Augenlidern) und besitzt keine eigene
Gefäßschicht.

Darauf folgt die \T{Lederhaut} (Dermis). In ihr sind alle Hautanhangsgebilde
(Schweißdrüsen, Talgdrüsen, Haare und Nägel) befestigt. Weiters befinden sich in
dieser Hautschicht Tastkörperchen sowie freie Nervenendigungen
(\Es{Schmerzempfindungen}). Außerdem finden sich hier zwei verschiedene
Gefäßsysteme (eines oberflächlicher, eines tiefer). Dieses dient wie bereits erwähnt der
Temperaturregulation.

Unter der Dermis befindet sich die \T{Unterhaut} (\T{Subcutis}). Sie besteht
hauptsächlich aus Fett. Dieses dient der Wärmeisolierung.
}
{
\begin{itemize}
    \item \T{Oberhaut} \Z{(Epidermis)} \index{Oberhaut} 
    \begin{itemize}
        \item Hornschicht, geschichtetes Plattenepithel
    \end{itemize}
    \item \T{Lederhaut} \Z{(Corium, Dermis)}
    \index{Lederhaut}
    \begin{itemize}
        \item Hautanhangsgebilde
        \begin{itemize}
            \item Schweißdrüsen
            \item Talgdrüsen
            \item Haarfollikel, Nagelmatrix
        \end{itemize}
        \item Tastkörperchen, freie Nervenendigungen (Schmerz)
        \item \Z{Kollagenfasern (Stütznetzwerk)}
        \item Arteriolen, Venolen (Versorgung der Haut, Temperaturregulation)
    \end{itemize}
    \item \T{Unterhaut} (\T{Subcutis}) \index{Unterhaut} 
    \index{Subcutis}
    \begin{itemize}
        \item Depotfett, Wärmeisolierung
    \end{itemize}
\end{itemize}\index{Haut!Querschnitt}
\AuthNotes{Haut: Bild Querschnitt: muss vereinfacht werden,
lat. Ausdrücke weg bzw. grau. GABS, 2010-01-23: done. Bild von Hirtler
übernommen.}
}
{./images/hirtler-aass/Haut-edited.png}
{\AuthNotesPicLegalOk{}{Haut
Querschnitt}Die Haut im Querschnitt: \TabSub{A:} Oberhaut \Z{Epidermis},
\TabSub{B:} Lederhaut \Z{Dermis},
\TabSub{C:} Unterhaut (Subkutis),
\TabSub{d:} \Z{Meißner'sches Tastkörperchen},
\TabSub{e:} Talgdrüse,
\TabSub{f:} Haar mit Haarfollikel,
\TabSub{g:} Schweißdrüse,
\TabSub{h:} Venole,
\TabSub{i:} Freie Nervenenden (Schmerzempfindung),
\TabSub{j:} Arteriole,
\TabSub{k:} \Z{Pacini-Körperchen}.
}
{Hirtler.}




\Layout{0}{Die Haut ist wichtig für die
\E{Thermoregulation}}
{
Neben der Lunge dient vor allen Dingen  die Haut als wichtigstes Organ
zur \Es{Thermoregulation}: Die Blutgefäße  der  Haut
sind stark ineinander verwickelt und können sich bei Bedarf stark erweitern. Die
stark durchblutete  Haut strahlt somit Wärme ab. Weiters wird durch
die Verdunstungskälte  des Schweißes  ebenfalls Wärme
abgegeben. 
}
{
\begin{itemize}
  \item Erweiterung - Kontraktion der Blutgefäße
  \item Verdunstungskälte des Schweißes
\end{itemize}
}{}{}{}

\Layout{0}{Wundliegen}
{
\index{Wundliegen}\index{Dekubitus}\label{topic:dekubitus}Wenn der Mensch nicht
fähig ist in ausreichendem Umfang regelmäßig seine Lage zu ändern kann es zum
wundliegen kommen. Die Wunden nennt man dann \T{Dekubitus} oder Druckgeschwüre.
\Z{Besonders gefährdet sind jene Stellen,
an denen ein Knochen knapp unter der Haut liegt
(Druckstelle). Betroffen sind vor allem alte oder
pflegebedürftige bettlägrige Menschen, bewusstlose
Patienten bzw. sedierte Intensivpatienten sowie Patienten
mit Querschnittsläsionen.}
}
{
\begin{itemize}
  \item Unfähigkeit der Lageänderung
  \item besonders bei prädestinierten Druckstellen
  \item Betroffene: Bettlägrige, Bewusstlose, Intensivpatienten,
  Querschnittsgelähmte
\end{itemize}
}{}{}{}

\clearpage
\section{Die Atemwege}
 \index{Atemwege}
 \label{topic:atemwege-lunge}


\AuthNotes{KOCH: mag Nasenrachenraum und Mundrachenraum
streichen, evtl. wie Bild?}

\Layout{37}{}
{}
{}
{./images/anatomie/Respiratory_system_complete_de-edited.pdf}
{\AuthNotesPicLegalOk{}{PD}Die Atemwege, Übersicht.\label{fig:respiraktionstrakt-uebersicht}}
{Mariana Ruiz Villarreal, Public domain}



%\begin{multicols}{2}

Die Atemwege gliedern sich in die oberen und die unteren Atemwege.

\Layout{0}{Obere Atemwege}
{%
Der Weg der Einatemluft beginnt in der Nase bzw. im Mund. Die Nase ist
mit Schleimhaut ausgekleidet, knöcherne Vorsprünge sorgen für
eine Vorwärmung der Luft. In der Schleimhaut be\-finden sich die
Sinneszellen für den Geruchssinn. 

\Z{In den Schädelknochen befinden sich luft\-ge\-füllte
Hohlräume, die so\-genannten \T{Neben\-höhlen} (Stirn-,
Kiefer\-höhlen, Siebbeinzellen). Diese stehen in Verbindung mit dem Nasenrachenraum.
Entzündet sich die Schleimhaut spricht man von der
\E{Neben\-höhlen\-ent\-zündung}, welche oft langwierig
und schmerzhaft sein kann.} 

Die Luft gelangt nun in den \T{Mund-Rachen\-raum} und den
\T{Rachen}. \Es{Hier kreuzt der Weg der Nahrung mit dem
Weg der Atemluft!} Die Nahrung gelangt vom Mundraum in die
hinten gelegene Speiseröhre, die Luft möchte vom Nasen\-rachen\-raum in die vorne gelegene Luftröhre. 

Um ein Verschlucken von Nahrung in die Lunge zu verhindern
verschließt der am Kehlkopf gelegene \T{Kehldeckel}
während des Schluckaktes die Luftröhre. Dieser Reflex ist lebenswichtig! 

Im \E{Kehlkopf}  befinden sich die \T{Stimmbänder},
welche für die Lautbildung beim Sprechen wichtig sind.
}{%
\begin{itemize}
    \item \Es{Nasen}- und \Es{Mundhöhle} \index{Nasenhöhle} 
    \index{Mundhöhle} 
    \item \Es{Rachen} (gemeinsamer Weg von Nahrung und
    Atem!) \index{Rachen} 
    \item \Es{Kehlkopf} (\T{Larynx}) mit
    \Es{Kehldeckel} \index{Kehlkopf} \index{Larynx} 
    \index{Kehldeckel}
\end{itemize}
}{}{}{}

\Layout{0}{Untere Atemwege}
{%
In der \T{Luftröhre} (\T{Trachea}) wird die vom Rachen
und Kehlkopf kommende Luft zur Lunge weitergeleitet. Die Luftröhre teilt sich in einen
\E{rechten} und einen \E{linken} \T{Hauptbronchus},
welche zu den jeweiligen Lungenflügeln ziehen. Jeder Hauptbronchus verzweigt sich immer weiter
bis schlussendlich sehr dünne \T{Bronchioli} in die
\T{Lungenbläschen} (\T{Alveolen}) münden. Hier findet
der Gasaustausch statt.

Die Abschnitte der Atemwege, in denen kein Gasaustausch
stattfindet und der nur dem Transport dient, bezeichnet man
als {\quotedblbase}\index{Totraum}\T{Totraum}{\textquotedblleft}.
}{%
\begin{itemize}
    \item \index{Luftröhre}\Es{Luftröhre}
    (\index{Trachea}\T{Trachea})
    \item \index{Hauptbronchus}\T{Hauptbronchus} (li./re.)
    \item \index{Bronchioli}\T{Bronchioli}
    \item \index{Lungenbläschen}\Es{Lungenbläschen}
    (\index{Alveolen}\T{Alveolen})
\end{itemize}
}{}{}{}



%\end{multicols}


\subsection{Die Lunge ist für einen Teil des
Gasaustausches zuständig}
\index{Gasaustausch!Lunge}  
\index{Lunge}
\index{Lunge!-nflügel}

\Layout{0}{\TxBeschreibung}
{
Die Lunge hat \Es{zwei Flügel}, welche sich jeweils in \Es{3 (rechts)} bzw.
\Es{2 (links) Lappen} gliedern: 

\begin{center}
\begin{tabularx}{\linewidth}[H]{XX}
\toprule
\multicolumn{2}{c}{\TabHeading{Die Lunge}}
\\\toprule
\multicolumn{1}{c}{\TabSub{Rechter Flügel}}
 &
\multicolumn{1}{c}{\TabSub{Linker Flügel}}
\\
\begin{itemize}
    \item \Es{3} Lappen
    \begin{itemize}
        \item Oberlappen \index{Oberlappen} 
        \item Mittellappen \index{Oberlappen} 
        \item Unterlappen \index{Unterlappen} 
    \end{itemize}
\end{itemize}
&
\begin{itemize}
    \item \Es{2} Lappen
    \begin{itemize}
        \item Oberlappen
        \item Unterlappen
    \end{itemize}
\end{itemize}
\\\bottomrule
\end{tabularx}
\end{center}

Die Oberfläche ist mit dem \T{Lungenfell}
\index{Lungenfell} \index{Pleura!lungenseitige} , der
\E{lungenseitigen} \T{Pleura}, überzogen. In einem Lungenflügel verzweigt
sich der Hauptbronchus in immer weitere kleinere Bronchien bis schließlich die
Bronchioli in die \T{Lungenbläschen} (\T{Alveolen}) enden.
}
{
\begin{itemize}
  \item 2 Flügel mit rechts 3 Lappen und links 2
  \item Lungenfell (Pleura)
  \item Lungenbläschen (Alveolen)
\end{itemize}
}{}{}{}

\subsubsection{Die Alveolen sind die Endstücke
des Atemwegs und sind für den Gasaustausch zuständig}

\Layout{0}{\TxBeschreibung}
{
\emph{Siehe Abb. \ref{fig:respiraktionstrakt-uebersicht}
rechts oben.}

In den \T{Lungenbläschen} (\T{Alveolen}) findet der
\index{Gasaustausch!Lunge}\E{Gasaustausch} zwischen der
Einatemluft und dem Blut statt. Sauerstoff gelangt dabei aus den Alveolen in das Blut und
Kohlendioxid wird im Gegenzug vom Blut an die Alveolen
abgegeben.

Dabei besteht zwischen Luft und Blut kein direkter Kontakt, die Gase
müssen Hindernisse wie die Alveolenwand und die Blutgefäßwand
überwinden. \E{Vergrößert sich dieser Weg aufgrund
einer Erkrankung ist die Sauerstoffaufnahme vermindert!} 

\Fazit{
Sauerstoff (\ce{O2}):   Alveolen ${\rightarrow}$ Blut

Kohlendioxid (\ce{CO2}):  Blut ${\rightarrow}$ Alveolen
}
}
{
\begin{itemize}
  \item Gasaustausch: Kohlendioxid gegen Sauerstoff
  \item Hindernisse: Alveolenwand, Blutgefäßwand
\end{itemize}
}{}{}{}

\subsubsection{Die Pleura}\label{topic:pleura}

\Layout{22}{\TxBeschreibung}
{%
Die Pleura hat zwei Blätter. Eine (Rippenfell) liegt an der Rippenseite an und
ist fest mit dem Zwerchfell und der Brustwand verwachsen.
Die Zweite (Lungenfell) überzieht die Lungenoberfläche.

Zwischen beiden Pleurablättern befindet sich ein Flüssigkeitsfilm. Dieser sorgt
dafür, dass bei Zug (durch Unterdruck während der Atmung) an einem der beiden
Blätter sich diese nicht trennen lassen (wie bei zwei Glasplatten mit Wasser
dazwischen) und das Zweite passiv mitbewegt wird.

Beim Gesunden gibt es keine Verwachsungen zwischen diesen beiden
Pleurablättern. Bei Entzündungen kann es jedoch dazu kommen.

Die Pleura ist außerdem die einziger Struktur der Lunge, die schmerzempfindlich
ist.
}
{\begin{itemize}
    \item \E{Thoraxseitige} Pleura (\T{Rippenfell}) \index{Rippenfell} 
    \index{Pleura!thoraxseitige}
    \item \E{Lungenseitige} Pleura (\T{Lungenfell}) \index{Lungenfell} \index{Pleura!lungenseitige}
    \item Atmung aufgrund Unterdruck und Haftung der beiden Blätter mit
    Flüssigkeit dazwischen möglich
    \item Normalerweise keine Verwachsungen der beiden Blätter
\end{itemize}
}
% {./images/noauth/AASdev20090228-img41.png}
{./images/hirtler-aass/Pleura.pdf}
{Pleura, Schema.}
{Hirtler.}




\paragraph{Welche Funktion hat aber nun die Pleura?}

\Layout{0}{\TxBeschreibung}
{
Da die eine Seite mit dem Brustkorb, die andere Seite mit der Lunge
verwachsen ist und der Raum zwischen den beiden Seiten luftdicht
abgeschlossen ist haften die beiden Seiten aneinander, aber sind
gegeneinander verschieblich. (Es gibt keine
Verwachsungen der beiden Seiten und die Pleuraflüssigkeit fungiert sogar ein
bisschen als "`Schmiermittel"'.) 

Wenn sich nun der Brustkorb oder das Zwerchfell bei der
Atmung bewegt, bewegt sich natürlich das Rippenfell mit (es ist ja mit dem Brustkorb
und dem Zwerchfell verwachsen). Da der Pleurazwischenraum
luftdicht ist, muss sich zwangsläufig aber auch das Lungenfell (an dem
wiederum die Lunge angewachsen ist) mitbewegen. Und wenn sich das Lungenfell
mitbewegt wird automatisch auch die Lunge wie eine Ziehharmonika
"`aufgezogen"' und Luft
"`angesaugt"'. 

\subparagraph{Frage:} \emph{Was würde passieren wenn der
Zwischenraum zwischen den Pleurablättern aufgrund einer Verletzung (z.\,B. ein
Loch aufgrund einer Stichwunde im Brustbereich) nicht mehr luftdicht
ist?} Antwort: s. Fußzeile (\footnote{Antwort: Die Lunge
würde in sich zusammenfallen, da sie ja nicht mehr aufgespannt wird. Siehe
"`Pneumothorax"'. \index{Pneumothorax!anatomische
Erklärung}})
}
{
\begin{itemize}
  \item keine Luft zwischen beiden Pleuren
  \item Pleuraflüssigkeit = Schmiermittel
  \item Unterdruck bei Atmung
  \item bei Verletzung bzw. Lufteindringen in den Pleuraspalt Pneumothorax
\end{itemize}
}{}{}{}



\subsection{Atemmechanik}
\label{topic:atemmechanik} \index{Atemmechanik} 

Wenn man sich den Thorax annähernd als Zylinder vorstellt,
so kann eine Volumsvergrößerung (Einatmung) durch eine
Vergrößerung des Querschnittes (Heben der
Zwischenrippenmuskulatur\index{Zwischenrippenmuskulatur})
oder Vergrößerung der Höhe (Anspannung des Zwerchfells) erfolgen.

\begin{equation}
V = A \times h
\end{equation}
\begin{equation}
\mbox{Volumen} = \mbox{Fläche} \times \mbox{Höhe}
\end{equation}

\AuthNotes{KOCH: entweder wird das genauer erklärt wie die
Funktion der Pleura oder es wird ganz einfach mündlich
erklärt.

GABS: Sollte schon schriftlich erklärt werden \ldots wir
wollen ja auch immer dass sie die "`Atemhilfsmuskulatur"'
etc. unterstützen, dann sollten sie wenigstens auch
nachlesen können "`wie man richtig atmet"'

KOCH: soll erklärt werden.}

\begin{figure}[H]
\begin{center}
% \includegraphics[width=\linewidth]{./images/noauth/AASdev20090228-img42.png}
\includegraphics[width=\linewidth]{./images/hirtler-aass/Atmung.pdf}

\Captionof{figure}{Atemmechanik. \SourcePicture{Hirtler.}}
\end{center}
\end{figure}


\Layout{0}{}
{
Für die Atemmechanik ist ein gutes Zusammenspiel von Skelett, Muskeln, Gewebe
und Pleura(spalt) wichtig. 

\Fazit{Ein bis zum Hals Verschütteter erstickt.}

\paragraph{Atemmuskulatur} Der wichtigste Atemmuskel ist
das \Es{Zwerchfell}\index{Zwerchfell}. In Ruhe erledigt es
fast die gesamte Atemarbeit. Bei verstärkter Atmung kommen die
Zwischenrippenmuskeln zum Einsatz \cite{LippertAnatomie7}.

\paragraph{Atemhilfsmuskulatur}%
\label{topic:atemhilfsmuskulatur}%
\index{Atemhilfsmuskulatur}%
Neben den Hauptatemmuskeln (Zwerchfell, Zwischenrippenmuskulatur) gibt es noch
eine Reihe weiterer Atemhilfsmuskeln die die Atmung unterstützen. Diese
verbinden den Thorax und die Oberarme. Fixiert man die Arme, z.\,B. durch
aufstützen, können die Muskeln die Hebung des
Brustkorbes unterstützen. Vgl. dazu \emph{'Asthma', \ref{topic:asthma-bronchiale}} und \emph{'COPD', \ref{topic:copd}}.

\paragraph{Pleura} Die Funktion der Pleura ist unter \ref{topic:pleura} erklärt.

\Niveau{2}{Thoraxwand und Lunge besitzen elastische Eigenschaften: Sie werden gedehnt
und entwickeln elastische Rückstellkräfte. An der Lunge können zwei Mechanismen ein
Zusammenfallen des Lungengewebes bewirken. Erstens sind im Lungenzwischengewebe
reichlich elastische Fasern eingelagert, die die Eigenelastizität der Lunge
bewirken. Durch die Koppelung von Lunge und
Thoraxwand am Pleuraspalt entsteht ein elastisches Gesamtsystem. 

Bei entspannter
Atemmuskulatur und offenem Kontakt mit der Umgebungsluft wird sich ein
Gleichgewicht der elastischen Kräfte einstellen; der Brustkorb befindet sich in
der Atemruhelage. Durch den Zug der Lunge, die sich verkleinern möchte, und den
Gegenzug der Thoraxwand entsteht ein ständiger Unterdruck im Pleuraspalt.}



\Niveau{2}{%
\subparagraph{Atemzyklus} 

Volumenschwankungen des Brustraums bewirken
gleichzeitig auch Volumenschwankungen im Inneren der Lunge. Eine Vergrößerung
des Thoraxinnenraums wird durch Kontraktion der Atemmuskulatur bewirkt. Dieser
Vorgang wird als aktive Phase der Atmung bezeichnet. Dabei kommt es vor allem
zu einer Abflachung und einem Tiefertreten des Zwerchfells, aber auch zu einer
Anhebung der Rippen durch die äußere Zwischenrippenmuskeln. Als Folge entsteht
in der Lunge im Vergleich zur Umgebungsluft ein Unterdruck, der durch
Einströmen von Luft ausgeglichen wird. Thoraxwand und Lunge werden dabei
gedehnt.

Am Ende einer ruhigen Einatmung
erschlafft die Einatmungsmuskulatur. Die elastischen Rückstellkräfte führen das
Thoraxvolumen wieder in die Ausgangslage zurück. Dabei kommt es zu einer
Druckerhöhung in der Lunge, die eine Luftbewegung hin zur Umgebungsluft bewirkt.
Diese erfolgt in aller Regel passiv, d. h. ohne aktive
muskuläre Anstrengung. Dieser Vorgang erfolgt in Ruhe bei einem Erwachsenen
etwa zwölfmal in der Minute und wird als Atemfrequenz bezeichnet.}


\Fazit{Der wichtigste Atemmuskel ist das \Es{Zwerchfell}!}
}
{
\begin{itemize}
  \item Atemmuskulatur: Zwerchfell, Zwischenrippenmuskeln
  \item Atemhilfsmuskulatur
  \item elastische Eigenschaften von Thorax und Lunge
  \item Unterdruck für Atmung wichtig
\end{itemize}
}{}{}{}

\subsection{Atemfrequenz, -tiefe und -volumen}

Die Atemfrequenz, die Atemtiefe und das Atemvolumen
werden im Kapitel Vitalfunktionen unter
\prettyref{topic:atemvolumina} besprochen.

\subsection{Atemgase}

\Definition{Als Atemgas bezeichnet man jenes Gasgemisch,
welches eingeatmet wird.}

\Layout{0}{\TxBeschreibung}
{
Normalerweise ist das ein Gasgemisch aus Stickstoff,
Sauerstoff, Kohlendioxid und Edelgasen ("`Raumluft"').
\index{Stickstoff}\index{Atemgas} 
\index{Sauerstoff!Raumluftgemisch} \index{Kohlendioxid!Raumluftgemisch}
Zwischen der normalen Einatem- und Ausatemluft besteht besonders hinsichtlich
des Sauerstoff- und Kohlendioxid-Anteils ein Unterschied, siehe 
\prettyref{tab:luft-gas}.

\begin{table}[H]
\Captionof{table}{Gasgehalt in der Raum- und
Ausatemluft.\label{tab:luft-gas}}

\begin{tabular}[H]{ccc}
\addlinespace
\centering \TabHeading{Raumluft} & & \TabHeading{Ausatemluft}
\\\addlinespace\midrule\addlinespace
\Es{21\,\%}  & Sauerstoff (\ce{O2}) & \Es{16\,\%}
\\\addlinespace
78\,\% & Stickstoff (\ce{N}) & 78\,\%
\\\addlinespace
\Es{0,03\,\%} & Kohlendioxid (\ce{CO2}) & \Es{4\,\%} 
\\\addlinespace
1\,\% & Edelgase  & 1\,\%
\\\addlinespace\bottomrule
\end{tabular}
\end{table}
}
{
\begin{itemize}
  \item Raumluft
  \item Atemgase: Sauerstoff und Kohlendioxid
\end{itemize}
}{}{}{}

\clearpage
\section{Der Verdaungstrakt -- Von der Schüssel in die
Schüssel}
\index{Magen-Darm-Trakt}
\index{Gastrointestinaltrakt}
\index{GIT}
\label{topic:gastrointestinaltrakt}

\emph{\Lang{Syn.} \T{Magen-Darm-Trakt},
\T{Gastrointestinaltrakt} (\T{GIT})}

\Layout{0}{\TxBeschreibung}
{
Der Verdauungstrakt dient der Verstoffwechselung der mit dem Essen
aufgenommenen Nährstoffe.

Er beginnt im Mund mit der Nahrungsaufnahme und deren Zerkleinerung. Durch
diverse Enzyme erfolgt die Weiterverarbeitung im Magen-Darm-Trakt. Wenn sie
klein genug zerlegt sind, werden sie durch die Darmzellen aufgenommen und in
das Blut überstellt.

Über das Blut werden die Nährstoffe an ihre Zielzellen
tranportiert.

Abfallstoffe landen schlußendlich wieder im Darm (bzw. den Nieren) um wiederum
ausgeschieden zu werden.
}
{
    \begin{itemize}
        \item Nahrungsaufnahme
        \item Weiterverarbeitung im Magen-Darm-Trakt
        \item Transport der Nährstoffe und Energieträger  ins Blut
        \item Zellaufbau und Abbau von Abfallstoffen
    \end{itemize}
}{}{}{}

\paragraph{Übersicht}~

\begin{figure}[H]
    \includegraphics[width=0.8\linewidth]{./images/anatomie/Digestive_system_diagram_de.pdf}
    \Captionof{figure}{\AuthNotesPicLegalOk{}{}Schema des Verdauungstraktes.
    \SourcePicture{Mariana Ruiz Villarreal, Public domain}}
\end{figure}


\begin{gWideParEnv}
\begin{minipage}{\linewidth}
\Captionof{table}{Übersicht über den Verdauungstrakt --
\emph{"`Von der Schüssel in die Schüssel."'}\label{tab:verdauungstrakt-uebersicht}}
\begin{tabulary}{\linewidth}[H]{LLL}
\toprule\addlinespace
\noindent\multirow{3}{8em}{\TabHeading{Mundhöhle}}
 &
 Zunge
 &
\noindent\multirow{2}{15em}{Zerkleinerung der Nahrung und Vermischung
mit Speichel}
\\\addlinespace\hhline{~-~}\addlinespace
 &
Zähne 
&
\\\addlinespace\hhline{~-~}\addlinespace
 &
Speicheldrüsen
&
\\\addlinespace\hhline{---}\addlinespace
Rachen
 &
 &
Kreuzung des Speisebreis mit dem Luftweg

Kehldeckel verschließt beim Schlucken die Luftröhre
\\\addlinespace\hhline{---}\addlinespace
\noindent\Es{Speiseröhre} (\T{Ösophagus})
 &
 &
Speisebrei wird mittels Muskelbewegung zum Magen befördert
\\\addlinespace\hhline{---}\addlinespace
\noindent\TabHeading{Magen}
 &
 &
Vermischung mit Magensäure

Magensäure zerlegt Nahrung

Pförtner gibt Portionsweise Speisebrei an den Zwölffingerdarm ab

\\\addlinespace\hhline{---}\addlinespace
\multirow{4}{8em}{\TabHeading{Dünndarm}}
&
\noindent\TabSub{Zwölffingerdarm} \Z{Duodenum}
 &
Gallenflüssigkeit (${\leftarrow}$Leber) und Bauchspeichel
(${\leftarrow}$Pankreas) wird zugesetzt
\\\addlinespace\hhline{~--}\addlinespace
 &
\noindent\TabSub{Krummdarm} \Z{Jejunum}
 &
Spaltung von Nährstoffen (Eiweiß, Kohlehydrate, Fett) und
Nährstoffaufnahme
\\\addlinespace\hhline{~-~}\addlinespace
 &
\noindent\TabSub{Leerdarm} \Z{Ileum}
 &
\\\addlinespace\hhline{---}\addlinespace
\multirow{6}{8em}{\TabHeading{Dickdarm}\\(\TabHeading{Colon})}
&
\noindent\TabSub{Blinddarm}
&
Sackgasse
\\\addlinespace
&
\noindent\TabSub{Wurmfortsatz} (\TabSub{Appendix}) 
 &
Anhängsel der Sackgasse, kann sich entzünden.
\\\addlinespace\hhline{~--}\addlinespace
 &
\noindent\TabSub{Aufsteigender} Teil
 &
\multirow{4}{\columnwidth}{Wasserentzug, \\Eindickung, \\
Bakterienbesiedlung, \\Produktion und Speicherung des
Stuhls} 
\\\addlinespace\hhline{~-~}\addlinespace
 &
\noindent\TabSub{Querender} Teil
 &
\\\addlinespace\hhline{~-~}\addlinespace
 &
\noindent\TabSub{Absteigender} Teil
 &
\\\addlinespace\hhline{~-~}\addlinespace
 &
\noindent\TabSub{S-förmig} gebogener Teil (\Z{Sigma})
 &
\\\addlinespace\midrule\addlinespace
\noindent\TabHeading{Mastdarm} (\TabHeading{Rektum})
 &
 &
{\quotedblbase}Speicherung{\textquotedblleft} des ausscheidungsbereiten
Stuhls
\\\addlinespace\hhline{---}\addlinespace
\noindent\TabHeading{Anus}
 &
 &
Absetzen des Stuhls

\\\addlinespace\bottomrule
\end{tabulary}

\end{minipage}
\end{gWideParEnv}


\subsection{Mundhöhle}\index{Mundhöhle}

\Layout{0}{\TxBeschreibung}
{
In der Mundhöhle kommt es durch die Zähne zur Zerkleinerung der Nahrung. Durch
die Kaubewegungen der Zähne und durch die Bewegungen der Zunge wird der Brei
mit Speichel aus den diversen Speicheldrüsen vermischt.

Der Speichel dient schon im Mund der Zerkleinerung von Stärke und Zucker in der
Mahlzeit.

Die Zunge schiebt den Bissen Schritt für Schritt in den Rachen. Beim Schlucken
sorgt sie dafür, dass der Kehldeckel verschloßen wird, da es sonst zu einem
Übertritt der Nahrung in die Luftröhre und damit zu einer Aspiration kommen
kann.
}
{
\begin{itemize}
    \item Zähne: Nahrungszerkleinerung
    \item Speicheldrüsen
    \item Zunge: Nahrungsfortbewegung, Kehldeckelverschluss
\end{itemize}
}{}{}{}

\subsection{Die Speiseröhre (Ösophagus)}
\index{Speiseröhre}
\index{Oesophagus@Ösophagus}
\label{topic:speiseroehre}


\emph{\Lang{Term.} \Terminus{Ösophagus}}


\Layout{22}{\TxBeschreibung}
{
Die Speiseröhre (Ösophagus) ist ein ungefähr 25\,cm langer Muskelschlauch. Er
dient dem Transport über den Mund aufgenommener Nahrung vom Rachen in den Magen.
Diese Fortbewegung funktioniert über eine wellenförmige Kontraktur der
Wandmuskeln der Speiseröhre. Im Ösophagus werden keine verdauenden Enzyme
ausgeschieden -- er dient rein zur Fortbewegung des Speisebreies.

\Z{Morphologisch weist die Speiseröhre 3 Engen auf, die die Nahrung überwinden
muss. Die erste ist hinter dem Kehlkopf (wenn hier ein Speisebrocken Stecken
bleibt kann er außerdem die Luftröhre verschließen. Landläufig nennt man das
den "`Wiener-Würschtel-Tod"'.). Die zweite Enge liegt dort, wo die Aorta die
Speiseröhre kreuzt. Die Dritte Enge entspricht dem Durchtritt des Ösophagus
durch das Zwerchfell.}

In der Wand des Ösophagus finden sich zahlreiche Venen. Bei manchen
Lebererkrankungen kann es zu einer Stauung kommen; es
entwickeln sich sog. \E{Ösophagusvarizen}.\ZFootnote{Z.\,B. fortgeschrittene
Leberzirrhose: Umbau der Leber in Bindegewebe durch z.B. chronische Entzündung und damit resultierendem Verschluss der Pfortader. Dabei
fungieren diese Venen als Umgehungskreislauf in die Vena cava.} Wenn eine solche
Ösophagusvarize platzt, kann es zu einem großen Blutverlust kommen
(\prettyref{topic:oesophagusvarizenblutung}). }
{
\begin{itemize}
    \item $\sim$ 25\,cm langer Muskelschlauch
    \item Nahrung: Rachen $\rightarrow$ Magen
    \item keine Verdauung
    \item Ösophagusvarizen
\end{itemize}
}
{./images/wikipedia-pd/esophagus.pdf}
{Ösophagus, Schema.}
{???}





\subsection{Der Magen}
\index{Magen}

\Z{\emph{\Lang{Term.} Gaster, Ventriculus.}}

%\AuthNotes{KOCH: Fließtext!}
\Layout{0}{\TxBeschreibung}
{\index{Salzsäure}\index{Magen!-pförtner}
Der Magen ist ein muskuläres Hohlorgan und dient der Nahrungsspeicherung,
Nahrungsdurchmischung sowie der kontrollierten Weiterleitung des
durchmischten Speisebreis. 

Verantwortlich für den Beginn der Verdauung der Nahrung ist der durch die
Magenschleimhaut gebildete saure, aus \Es{Salzsäure} bestehende, \Es{Magensaft}.
Die Magenschleimhaut produziert ausserdem eine \Es{schützende Schleimschicht},
die die Oberfläche vor der starken Säure schützt. 

Nachdem die Wandmuskulatur den Speisebrei mit dem Magensaft \Es{durchmischt} hat,
gibt der \Es{Magenpförtner} \Z{(Pylorus)} den Speisebrei portionsweise an den
Dünndarm -- genauer: den Zwölffingerdarm -- weiter.
}{
\begin{itemize}
    \item Magensaft mit Salzsäure
    \item Schleim schützt Magenschleimhaut
    \item Wandmuskulatur durchmischt Speisebrei
    \item Portionsweise Abgabe von Speisebrei durch den
    Magenpförtner in
    den Zwölffingerdarm
\end{itemize}
}


% \Parallel{
% \TODO{GABS}{2010-04-11}{Bild: neu erstellen: Lage des Magens.}{}
% \centering
% % \includegraphics[width=0.7\linewidth]
% % {./images/noauth/AASdev20090228-img46.png} 
% \includegraphics[width=0.7\linewidth]{./icons/under-construction.png}
% 
% \Captionof{figure}{\AuthNotesPicLegalNotOk{img46
% img45}{}Lage des Magens. \SourcePicture{Quelle nicht
% eruierbar.}}
% }{
% \TODO{GABS}{2010-04-11}{Bild: neu erstellen: Magen, Pförtner, Duodenum.}{}
% 
% %  \includegraphics[width=\linewidth]{./images/noauth/AASdev20090228-img47.png} 
% \includegraphics[width=\linewidth]{./icons/under-construction.png}
% 
%  \Captionof{figure}{\AuthNotesPicLegalNotOk{img47}{}Der
%  Magen, der Pförtner und das erste Stück
%  des Zwölffingerdarms im Durchschnitt.
%  \SourcePicture{Quelle nicht eruierbar.}}
%  
% }




\subsection{Zwölffingerdarm (Duodenum)}
% TODO GIT: Fließtext mangelhaft
% TODO GIT: SLIDE fehlt

\begin{itemize}
\item Beginn des Dünndarmes, ca 25\,cm lang
\item Hier münden der Gallengang und der Pankreasgang
ein. Es wird aus dem Gallengang die Gallenflüssigkeit,
welche für die Fettverdauung wichtig ist (Emulsion der
Fette), und aus dem Pankreas der Bauchspeichel mit
Verdauungsenzymen abgegeben. 
\end{itemize}

\opt{STATUS-EXPERIMENTAL}{
\Parallel{

\TODO{GABS}{2010-04-11}{Bild: neu erstellen: Lage des Duodenums.}{}

% \includegraphics[width=\linewidth]{./images/noauth/AASdev20090228-img48.png}
\includegraphics[width=\linewidth]{./icons/under-construction.png}

\Captionof{figure}{\AuthNotesPicLegalNotOk{}{}Die Lage des
Zwölffingerdarms im Körper.} 

}{

\TODO{GABS}{2010-04-11}{Bild: neu erstellen: Duodenum und Pankreas.}{}

\includegraphics[width=\linewidth]{./icons/under-construction.png}
% \includegraphics[width=\linewidth]{./images/noauth/AASdev20090228-img49.png}
\Captionof{figure}{\AuthNotesPicLegalNotOk{}{}Der
Zwölffingerdarm und die Bauchspeicheldrüse.} 

}
}




\subsection{Die Leber}\index{Leber}

% \AuthNotes{KOCH: fehlt: was ist sie, eine Drüse, ein
% Hohlmuskel, \ldots } GABS: Done
% TODO GIT: Fließtext mangelhaft
% TODO GIT: SLIDE fehlt

\Z{\Lang{Term.} Hepar.}

\noindent Die Leber ist ein solides Organ und befindet sich im \Es{rechten
Oberbauch}. Sie hat einen linken und rechten Lappen.

\Layout{0}{Aufgaben}{
\begin{itemize}
\item "`Chemiefabrik"' des Körpers: Produktion von vom
    Körper benötigten Eiweißen
    \begin{itemize}
        \item Produktion von vom
    Körper benötigten Eiweißen
    \item Entgiftung: Abbau von Giftstoffen (die Leber nimmt
    durch die Giftstoffe mit der Zeit selbst auch Schaden
    (Leberentzündung, Leberzirrhose)
    \item erhält über \index{Pfortader}\T{Pfortader}
    (Vena portae) nährstoffreiches Blut vom Darm zwecks Entgiftung   
\end{itemize}
    \item {\quotedblbase}Verdauungsdrüse{\textquotedblleft}
    \begin{itemize}
        \item produziert \index{Gallenflüssigkeit}Gallenflüssigkeit
        ({\quotedblbase}Galle{\textquotedblleft}), welche in der Gallenblase
        gespeichert wird: Fettverdauung
        \item Von der Leber gehen die Gallengänge, an den Gallengängen
        hängt die
    \end{itemize}
    \item Speicherung von Fett \& Zucker
    
\end{itemize}

\TODO{GABS}{2010-04-11}{Bild: neu erstellen: Leber.}{}

% \includegraphics[width=\linewidth]{./images/noauth/AASdev20090228-img50.png}
}{

}
{}
{}
{}    


\subsubsection{Gallenblase}
\index{Gallenblase}
\label{topic:gallenblase}

Am Gallengang angehängt findet sich die \T{Gallenblase}.
Sie dient nur als Zwischenspeicher für die Gallenflüssigkeit.

\AuthNotes{GABS: Erklärung/Bild Pfortader einfügen? fragt
STEU}




\Dias{Bilderserie: Leber und Gallenblase.}
{
% \includegraphics[width=\linewidth]{./images/noauth/AASdev20090228-img51.png}
\includegraphics[width=\linewidth]{./images/wikipedia-pd/Liver_1.png}

\Captionof{figure}{Eine Leber.
\SourcePicture{WmCo/Mikael Häggström, PD}}
}
{
% \includegraphics[width=\linewidth]{./images/noauth/AASdev20090228-img52.png}
\includegraphics[width=\linewidth]{./images/wikipedia-pd/Liver_and_pangreas_-_transparent.png}

\Captionof{figure}{\AuthNotesPicLegalNotOk{img52}{}Leber, Gallenblase und
Pankreas.
\SourcePicture{Quelle nicht eruierbar.}}
}
{
\includegraphics[width=\linewidth]{./images/wikipedia-pd/Gray1095-gall_bladder.png}

\Captionof{figure}{Gallenblase.
\SourcePicture{Gray's Anatomy, PD, Copyright expired.}}
}


 

\subsection{Bauchspeicheldrüse}

\index{Bauchspeicheldrüse}%
\index{Pankreas}%
\Lang{Syn.} \T{Pankreas}

\Layout{22}{\TxBeschreibung}
{
Das Pankreas ist ca. 15\,--\,22\,cm lang und liegt hinter dem
Magen. Es hat einen \Es{Ausführungsgang}, welcher
zusammen mit dem Gallengang \Es{in den Zwölffingerdarm}
(Duodenum) mündet.

\paragraph{Aufgaben} 

Das Pankreas hat \Es{zwei} wesentliche Funktionen:

\begin{description}
    \item[Enzymproduktion] für Verdauung
    ${\rightarrow}$ Zwölffingerdarm
    \item[Hormonproduktion] Ausschüttung ins Blut
    \begin{itemize}
        \item \index{Insulin}Das Hormon Insulin spielt eine wichtige Rolle im 
        Zuckerstoffwechsel (\prettyref{topic:insulin}) ${\rightarrow}$
        Blutzucker ${\downarrow}$
        \item \Z{Glukagon ${\rightarrow}$ Blutzucker ${\uparrow}$}
    \end{itemize}
\end{description}
}{}
{./images/wikipedia-pd/Gray1100.png}
{Das Pankreas und der Zwöffingerdarm.}
{}



%\TempPicInWork{Gallensystem}

\subsection{Krumm- und Leerdarm -- Jejunum und Ileum}
\index{Krummdarm}\index{Leerdarm}\index{Jejunum}\index{Ileum}

\T{Jejunum}, \T{Ileum}
% TODO GIT: Fließtext mangelhaft

\Layout{60}{\TxBeschreibung}
{
\begin{itemize}
    \item ca. 3\,m lang
    \item Oberflächenvergrößerung durch Darm\-zotten
    \item \Es{Nährstoffaufnahme}\index{Nährstoffaufnahme}
    \item \Z{beweglich, am Darmgekröse (Mesenterium)
    aufgehängt}
\end{itemize}
}{}
{./images/wikipedia-pd/Small_intestine.pdf}
% {./images/noauth/AASdev20090228-img56.png}
{Dünndarm.}
{WmCo/Mikael Häggström, PD.}



\subsection{Dickdarm (Colon)}
\index{Dickdarm}\index{Colon}

Der Dickdarm (Colon) ist ca. 1,5\,m lang und wird in
verschiedene Abschnitte unterteilt: 

% TODO Dickadarm: Tab: Enddarm, etc. fehlt
\begin{center}
\begin{tabulary}{\linewidth}[H]{LLL}
\toprule
\multirow{6}{8em}{\Es{Dickdarm}\\(\T{Colon})}
&
\Es{Blinddarm}\index{Blinddarm}
&
Sackgasse
\\
&
\Es{Wurmfortsatz} (\T{Appendix})
\index{Wurmfortsatz}\index{Appendix} &
Anhängsel der Sackgasse, kann sich entzünden.
\\\hhline{~--}
 &
\Es{Aufsteigender} Teil \index{Dickdarm!aufsteigender
Teil} &
\multirow{4}{\columnwidth}{Wasserentzug, Eindickung, \\
Bakterienbesiedlung, \\Speicherung} 
\\\hhline{~-~}
 &
\Es{Querender} Teil \index{Dickdarm!querender 
Teil}
 &
\\\hhline{~-~}
 &
\Es{Absteigender} Teil \index{Dickdarm!absteigender
Teil}
 &
\\\hhline{~-~}
 &
\Es{S-förmig} gebogener Teil (\T{Sigma})
\index{Dickdarm!S-förmig gebogener Teil} &
\\\hhline{~-~}
\Es{Enddarm} (\T{Rectum})\index{Enddarm}&
\\\bottomrule
\end{tabulary}
\end{center}

\Layout{60}{\TxBeschreibung}{
Der Leerdarm mündet in den Dickdarm. Diese Mündung ist wie eine
\E{T-Kreuzung} gebaut: \E{Nach oben} geht es in den
\Es{aufsteigenden Teil} des Dickdarmes, \E{nach unten} in
den \Es{Blinddarm}. Wie der Name des Blinddarmes schon sagt
endet er blind, er ist daher eine Sackgasse. Am Blinddarm
angehängt findet man den \Es{Wurmfortsatz}
(\T{Appendix}) welcher sich häufig entzündet
(\T{Appendiz\underline{itis}\footnote{Das Anhängsel
\emph{"`-itis"' zeigt immer eine Entzündung an, vgl.
Kapitel Hygiene}}}). Der Wurmfortsatz erfüllt beim
Menschen keine nennenswerte Funktion, außer der
Arbeitsplatzsicherung von Anästhesisten und Chirurgen
(wegen der vielen Blinddarm-Operationen).
}{

}
{./images/wikipedia-pd/Colon_anatomy.pdf}
{Der Dickdarm.}
{???}


%%%%%%%%%%%%%%%%%%%%%%%%%%%%%%%%%%%%%%%%%%%%%%%%%%%%%%%%%%%
% TODO : Vervollständigen Mastdarm/Rektum fehlt
\TODO{GABS}{2010-06-02}{Vervollständigung notwendig: Mastdarm/Rektum}{}
%%%%%%%%%%%%%%%%%%%%%%%%%%%%%%%%%%%%%%%%%%%%%%%%%%%%%%%%%%%



% \clearpage
% \vspace{2cm}

\section{Sonstige Organe und Strukturen im Bauchraum}
\label{topic:sonstige-strukturen}
\label{topic:sonstige-bauchorgane}

\subsection{Bauchfell}
\index{Bauchfell}
% REVIEW Bauchfell: Überarbeitung
\A{Review: Bauchfell: Überarbeitung}
\A{Review: Kommentare eingefügt.}


\Layout{20}
{}
{%
Das Bauchfell (\T{Peritoneum}) ist eine dünne Haut, welche \Es{die gesamte
Bauchhöhle auskleidet}. Dieser Überzug bedeckt die meisten
Bauchorgane\ZFootnote{außer z.\,B. die Nieren} \Z{und das große und kleine 
Netz\ZFootnote{Omentum majus und minus}}.
\ACh{GABS}{2010-03-09}{Unter großem und
kleinen Netz wird sich kaum wer was vorstellen können \ldots} 

Es sorgt durch seine feuchte und
glatte Oberfläche dafür, dass die Organe, wie z.\,B. Darm, ihre Bewegungen zur Verdauungen besser durchführen können.
Das Bauchfell ist auch verantwortlich für die \Es{Schmerzen}, die z.\,B. bei
einer Entzündung der Verdauungsorgane (z.\,B. Blinddarmentzündung
(Appendizitis)) entstehen. Es kann sich auch selbst entzünden
(\E{Bauchfellentzündung} (Peritonitis, \prettyref{topic:bauchfellentzuendung})),
z.\,B. durch eine Infektion oder wenn es bei einem Durchbruch
(\E{Perforation}) des Darmes zu einem Eintritt von Darminhalt in die Bauchhöhle
kommt. 
}{
\begin{itemize}
    \item Auskleidung der gesamten Bauchhöhle
    \item dünner Überzug
    \item Bewegung der Organe
    \item schmerzempfindlich
    \item Bauchfellentzündung (Peritonitis)
\end{itemize}
}
{./images/gabriel-sebastian-aass/Torso_GABS_IMG_1251_00800_frei.jpg}
{Der Brust- und Bauchraum, geöffnet am Torso in
Aufsicht. Zu sehen sind die Lunge, dahinter das Herz, darunter die
Leber (braun) mit einem kleinen Stück der Gallenblase (grün), der Magen mit
einem Teil der gelben Fettschürze \Z{"`großes Netz"'}, sowie der Dickdarm und
der Dünndarm. } 
{GABS.}


\subsection{Die Milz} 
\index{Milz}
% TODO Milz: Fließtext mangelhaft
\TODO{GABS}{yyyy-mm-dd}{Milz: Fließtext mangelhaft}{}
% TODO Milz: Beschreibung mangelhaft
\TODO{GABS}{yyyy-mm-dd}{Milz: Beschreibung mangelhaft}{}
% TODO Milz: SLIDE fehlt
\TODO{GABS}{yyyy-mm-dd}{Milz: SLIDE fehlt}{}

\Z{\Lang{Syn.} Lien, Splen.}

\Layout{2}{}
{
\begin{itemize}
\item hat nichts  mit der Verdauung  zu tun 
\item 12\,cm, oval, von \Es{Kapsel}\index{Kapsel!Milz}
umgeben
\item im linken Oberbauch unter den Rippen gelegen
\item Teil des Immunsystems
\item entfernt alte Erythrozyten, beim Embryo auch Blut\-bildung
\item sehr gut durchblutet, Blutspeicher (Gefahr:
\index{Milzriss}Milzriss)
\end{itemize}

\vspace{1em}

\Cave{{\quotedblbase}\index{Zweizeitiger
Milzriss}\index{Milz!-riss, zweizeitiger}\T{Zweizeitiger
Milzriss}{\textquotedblleft}:

Zuerst reißt die Milz ein, die Kapsel bleib erhalten, es blutet zuerst in die Kapsel und der Druck steigt.
Nach einer gewissen Zeit reißt auch die Kapsel und es blutet nun
ungehindert aus der inneren Wunde. \Es{Der Patient wirkt
zuerst {\textquotedblleft}ganz normal{\textquotedblleft}, unerwartet und
plötzlich verfällt er.}}

\TODO{GABS}{2010-04-11}{Bild: neu erstellen: Milz, Lage.}{}


}{}
{./icons/under-construction.png}
% {./images/noauth/AASdev20090228-img60.png}
{\AuthNotesPicLegalNotOk{img60}{}Milz:
Lage.}
{Quelle nicht eruierbar.}



 



\subsection{Die Niere(n) (Ren)}
\index{Niere}
% TODO Niere: Fließtext mangelhaft
\TODO{GABS}{yyyy-mm-dd}{Niere: Fließtext mangelhaft}{}
% TODO Niere: SLIDE fehlt
\TODO{GABS}{yyyy-mm-dd}{Niere: SLIDE fehlt}{}

\Layout{23}{}
{
\begin{itemize}
    \item paarig angelegtes Organ
    \item stark durchblutet, empfindlich auf RR- Abfall ${\rightarrow}$
    Funktion ${\downarrow}$
    \item reguliert
    \begin{itemize}
        \item \Es{Wasser-} und
        \Es{Elektrolythaushalt}
        \index{Wasserhaushalt!Niere}
        \index{Elektrolythaushalt!Niere}
        \item Volumshaushalt\index{Volumshaushalt!Niere}
        \item pH
        (Säure/Basen-Haushalt)
        \index{Säure-Basen-Haushalt!Niere}
    \end{itemize}
    \item \Z{Hormonproduktion}
    \item Geht im Alter, besonders bei Diabetikern, kaputt ${\rightarrow}$
    {\quotedblbase}\T{Chronische
    Niereninsuffizienz}{\textquotedblleft} ${\rightarrow}$
    {\quotedblbase}\T{Blutwäsche}{\textquotedblleft}
    (\index{Dialyse}\T{Dialyse})
\end{itemize}

\TODO{GABS}{2010-04-11}{Bild: neu erstellen: Niere, Lage.}{}

}{}
{./images/wikipedia-pd/Kidneys.png}
% {./images/noauth/AASdev20090228-img61.png}
{Die Nieren.}
{WmCo/Mikael Häggström and Madhero88, PD.}




\paragraph{Aufbau}~

\begin{itemize}
\item Rinde, Mark
\item Hilus (Stiel)
\item Gefäße
\item Nierenbecken mit Anschluss an den
Harntrakt\index{Nierenbecken}
\end{itemize}

\AuthNotes{Niere: Anatomie und Funktion muss noch
ausgearbeitet werden.}


\AuthNotes{KOCH: grau:}
\paragraph{\Z{Mikroskopischer Aufbau}}

\Z{
\begin{itemize}
    \item Nephron ist die funktionelle Einheit (ca. 1
    Mio/Niere)
    \begin{itemize}
        \item Glomerulus: Hier wird das Blut gefiltert und Plasmawasser
        abgepresst (ca. 150\,l\,/\,Tag)
        \item Schleifenapparat (Tubulus): Das abgepresste Plasma\-wasser wandert
        den Schleifen\-apparat entlang, hier werden die meisten wichtigen
        Stoffe und der Groß\-teil des Wassers wiederum zurückgeholt,
        sodass am Ende nur mehr 1,5 -- 2~l Harn pro Tag in der Harnblasen
        landen.
    \end{itemize}
\end{itemize}
}

\Z{
\paragraph{Funktionsweise}

\begin{itemize}
    \item 1500\,l Blut pro Tag werden
    {\quotedblbase}bearbeitet{\textquotedblleft}
    \item Glomerula: Filtration ${\rightarrow}$ 150\,l
    \index{Primärharn}Primärharn (enthält Elektrolyte, etc.) wird
    gewonnen und fließt durch den Schleifenapparat
    \item Schleifenapparat: {\quotedblbase}Recycling{\textquotedblleft},
    Gutes kommt zurück ins Blut, der Rest fließt weiter Richtung
    Harnblase ${\rightarrow}$ 1,5-2\,l (End-)Harn/Tag
\end{itemize}
}
\marginline{\cite{LoefflerBiochemie7}}

\subsection{Die Nebenniere(n) (Adren)}

% \TempPicMissing{Nebennieren}

\Z{Paariges, jeweils am oberen Nierenpol gelegenes Organ.

Hat nichts  mit dem Harntrakt oder der eigentlichen Nierenfunktion zu
tun.

\Es{Hormonproduktion}
\Z{
\begin{itemize}
    \item Cortison
    \item Sexualhormone (Androgene)
    \item Adrenalin , Noradrenalin
\end{itemize}}

Die Nebennieren sitzen wie {\quotedblbase}Mützen{\textquotedblleft}
auf den Nieren, haben mit deren Funktion allerdings
\E{nichts} zu tun.
}

\subsection{Harnableitendes System -- Harntrakt}
\index{Harn!-ableitendes System}\index{Harn!-trakt}

\Layout{20}{\TxBeschreibung{}}
{
Der Harntrakt besteht aus 
\begin{inparaitem}
    \item Nieren inkl. Nierenbecken
    \item \index{Harn!-leiter}je 1 Harnleiter 
    \Z{(Ureter)}, 30\,cm lang, mündet  in die 
    \item \index{Blase}Blase \index{Harn!-blase}
    \item Harnröhre  (\Z{Urethra,} Mann: 25\,cm, Frau
    4\,cm)\index{Harn!-röhre}
\end{inparaitem}

Die Harnröhre verläuft beim Mann von der Harnblase durch die
Prostata \index{Harn!-röhre!Verlauf} und den Penis und
mündet an der Eichel in die Außenwelt. Bei einer Vergrößerung der Prostata
kann der Harnabfluss bis zur Unpassierbarkeit erschwert sein.  

Bei der Frau ist die Harnröhre kürzer als beim Mann und
mündet oberhalb des Scheideneinganges. 
}{
\begin{itemize}
    \item Nieren inkl. Nierenbecken
    \item \index{Harn!-leiter}je 1 Harnleiter 
    \Z{(Ureter)}, 30\,cm lang, mündet  in die 
    \item \index{Blase}Blase \index{Harn!-blase}
    \item Harnröhre  (\Z{Urethra,} Mann: 25\,cm, Frau
    4\,cm)\index{Harn!-röhre}
\end{itemize}
}
{./images/anatomie/anatomie-harntrakt-edited.png}
{\AuthNotesPicLegalOk{}{Lena}Harntrakt
mit Nieren (rechte Niere im Querschnitt mit Nierenbecken),
Harnleiter, Blase und Harnröhre. \label{fig:harntrakt}}
{Lena Hirtler}







\subsection{Was gibt es noch?}

Außer den "`klassischen Bauchorganen"' liegen auch andere
Organsysteme im Bauchraum auf die nicht vergessen werden
sollte. Besonders betrifft das die weiblichen
Fortpflanzungsorgane wie Eierstöcke, Eileiter und Uterus.
Diese werden allerdings im Kapitel "`Geschlechtsorgane"'
gesondert besprochen. 

\clearpage
\section{Geschlechtsorgane}
\subsection{Männliche Geschlechtsorgane}

% TODO Männliche Geschlechtsorgane: Fließtext mangelhaft
% TODO Männliche Geschlechtsorgane: SLIDE fehlt

\begin{gWideParEnv}
\Parallel{
\begin{itemize}
    \item innen:
    \begin{itemize}
        \item Hoden \index{Hoden}
        \begin{itemize}
            \item Spermienproduktion
            \item Geschlechtshormon-Produktion
        \end{itemize}
        \item Nebenhoden: Spermienreifung und Lagerung
        \item Samenstrang, besteht aus
        \begin{itemize}
            \item Samenleiter \index{Samenleiter}
            \item Nerven
            \item Gefäße
            \item durchzieht Leistenkanal
        \end{itemize}
        \item \index{Prostata}\T{Prostata}
        (Vorsteherdrüse)
        \begin{itemize}
            \item produziert Samenflüssigkeit (Beweglichkeit der Spermien)
            \item \Es{Vergrößert sich im Alter} und
            kann die \Es{Harnröhre abdrücken}, es kommt
            zum Harnstau ${\rightarrow}$ häufiger Notfall!
        \end{itemize}
    \end{itemize}
    \item außen:
    \begin{itemize}
        \item Glied (Penis): Schwellkörper, Harnsamenröhre, Begattungsorgan
        \item Hodensack \Z{(Skrotum)}: umhüllt Hoden und
        Nebenhoden
        \begin{itemize}
            \item Hodenstiel kann sich verdrehen
            \index{Hodenstiel} ${\rightarrow}$ \Ca{Notfall!}
        \end{itemize}
    \end{itemize}
\end{itemize}
}{
\includegraphics[width=\linewidth]{./images/anatomie/anatomie-genitale-maennlich-edited.png}
\Captionof{figure}{\AuthNotesPicLegalOk{}{}Die männlichen Geschlechtsorgane im Querschnitt. \SourcePicture{Lena Hirtler}} 
}
\end{gWideParEnv}


 

\subsection{Weibliche Geschlechtsorgane}
% TODO Weibliche Geschlechtsorgane: Fließtext mangelhaft
% TODO Weibliche Geschlechtsorgane: SLIDE fehlt

\begin{gWideParEnv}
\Parallel{
\begin{itemize}
    \item innen:
    \begin{itemize}
        \item \Es{Eierstöcke} \Z{(Ovar)}: Reifung der
        Eizellen \index{Eierstöcke}
        \item \Es{Eileiter} \Z{(Tuben)}: Transport des
        reifen Eis in Richtung Gebärmutter \index{Eileiter}
        \item \Es{Gebärmutter} \index{Gebärmutter}
        (\index{Uterus}\T{Uterus}): dehnbarer Hohlmuskel,
        mit Schleimhaut ausgekleidet,
        \begin{itemize}
            \item Schleimhaut wird regelmäßig aufgebaut und wieder
            abgestoßen (Menstruation, Regelblutung)
            \item Hier nistet sich die befruchtete Eizelle
            ein.
            \item Der Gebärmutterhals
            \index{Gebärmutterhals}
            (\T{Cervix}\index{Cervix}) führt in Richtung der Scheide.
            \item Die Mündung der Gebärmutter in die
            Scheide ist der \T{Muttermund}.
            \index{Muttermund!Anatomie}
        \end{itemize}
        \item \Es{Scheide}
        (\T{Vagina})\index{Scheide}\index{Vagina}
    \end{itemize}
    \item außen:
    \begin{itemize}
        \item Schamlippen \Z{(Labien)}:
        goße/äußere, kleine/innere
        \item \T{Klitoris}\index{Klitoris}: entspricht der
        männlichen Eichel
        \item Öffnung der Harnröhre
    \end{itemize}
\end{itemize}
}{
\includegraphics[width=\linewidth]{./images/anatomie/anatomie-genitale-weiblich-edited.png} \Captionof{figure}{\AuthNotesPicLegalOk{}{Lena }Weibliche Geschlechtsorgane, Querschnitt.
    \SourcePicture{Lena Hirtler}}
}
\end{gWideParEnv}




\subsubsection{Der weibliche Zyklus}
\label{topic:weiblicher-zyklus}\index{Zyklus, weiblicher}

\begin{itemize}
    \item Der weibliche Zyklus dauert 28 Tage
    \item gerechnet von Menstruation zu Menstruation  (Monatsblutung)
    \item Uterusschleimhaut wird aufgebaut und wieder abgestoßen
    \item hormongesteuert
    \begin{itemize}
        \item Östrogen, Progesteron
    \end{itemize}
    \item ca. am 14. Tag: Eisprung
    \item Befruchtung nein:
    \begin{itemize}
        \item Zyklus läuft normal weiter, Menstruation
    \end{itemize}
    \item Befruchtung ja:
    \begin{itemize}
        \item Eizelle nistet sich im Uterus ein, Zyklus gestoppt
    \end{itemize}
\end{itemize}

Was kann dabei schief gehen? Wenn sich die Eizelle schon im Ei\-leiter
einnistet kommt es zu einer
\index{Eileiter\-schwanger\-schaft}\Es{Eileiter\-schwanger\-schaft}
die lebens\-bedrohlich ist (Eileiter reißt und die Frau verblutet).

\noindent\begin{minipage}{\linewidth}

\TODOCRITICAL{GABS}{2010-04-11}{Bild: neu erstellen: Zyklus.}{}

\includegraphics[width=\linewidth]{./images/noauth/AASdev20090228-img65.png}

\Captionof{figure}{\AuthNotesPicLegalNotOk{img65}{}Der weibliche Zyklus. Oben:
Follikelreifung. Unten: Auf- und Abbau der Uterusschleimhaut. Am Tag 0 des
Zyklus wird die Uterusschleimhaut abgestoßen, es kommt zur Regelblutung,
danach wird die Schleimhaut wieder aufgebaut. Am Tag 14 erfolgt der Eisprung.
Kommt es zu keiner Befruchtung wird die Schleimhaut erneut am Tag 28
abgestoßen. Nach einer Befruchtung wandert die Eizelle weiter in den Uterus
und nistet sich dort ein.
\SourcePicture{Quelle nicht eruierbar.}}
\end{minipage}



 

% \chapter[{[NEW]} Newsflash -- Aktueller Dienst]{[NEW]\\Newsflash -- Aktueller
Dienst}
% \addtocounter{chapter}{1}


% 


% \chapter{Fallbeispiele und Arbeitsaufgaben}
% 
% \addtocounter{chapter}{0}
% \part{ILS: Notfallrettung Grundlagen}
% {\mdseries
% Entspricht SAN 3-2-1-Einschulung}
% 
% \chapter{Anamnese und Untersuchung}
% 
% \chapter{MPG}
% 
% \chapter{Megacodetraining}
% 
% \addtocounter{chapter}{0}
% \part{Intermediate: NFS-Aufbaumodul}
% \addtocounter{chapter}{0}
% \part{Advanced: Arzneimittel}
% \addtocounter{chapter}{0}
% \part{Advanced: Minimalinvasive Tätigkeiten und Behandlung}

\appendix

%\setcounter{page}{0}

% \pagenumbering[Appendix]{bychapter}


% \chapter{Verzeichnisse}


\chapter{Curricula}
\label{chp:INT}



\section{Österreichische Sanitäter-Ausbildungsverordnung}

\HeadingSub{SanAV}

\subsection{Ausbildungsziele des Modul 1}


\begin{WidePar}
% \scriptsize

\begin{multicols}{2}

Ausbildungsziele für österreichische Rettungssanitäter gem. der \emph{Verordnung
der Bundesministerin für Gesundheit und Frauen über die Ausbildung zum
Sanitäter - Sanitäter-Ausbildungsverordnung - San-AV StF: BGBl. II Nr. 420/2003
-- Anlage 1, Modul I}.

\begin{enumerate}
  
  \item \TabHeading{Erste Hilfe und erweiterte Erste Hilfe}
  
  (Bestandteil des Sachgebiets 'Sanitätshilfe')
  
  \begin{inparadesc}
    \item[Stundenanzahl:] 14
    \item[Lehrkraft:] Notarzt, Arzt für Allgemeinmedizin, approbierter Arzt,
    Facharzt, Lehrsanitäter, fachkompetente Person
  \end{inparadesc}
  \begin{enumerate}
    \item
    Notwendigkeit, Verpflichtung der Hilfeleistung
    \\
    \noindent\FullRef{topic:behandlungspflichtJUS},
    \FullRef{topic:behandlungspflicht-ehi}
    \item
    Unfallverhütung
    \\
    \FullRef{topic:verhalten-notfallort},
    \FullRef{topic:gefahrenzone}
    \item
    Rettungskette
    \\
    \noindent\FullRef{topic:rettungskette}
    \item
    Gefahrenzonen
    \\
    \noindent\FullRef{topic:gefahrenzone}
    \item
    Bergen (Wegziehen, Rautekgriff aus dem  Auto, Sturzhelm)
    \\
    \FullRef{topic:wegziehen},
    \FullRef{topic:rautekgriff}, 
    \FullRef{topic:sturzhelmabnahme}
    \item
    Notfallpatient, Notfalldiagnose, Notfallhilfe
    \\
    \FullRef{topic:notfalldiagnose-bewusstlosigkeit},
    \FullRef{topic:notfalldiagnose-atem-kreislaufstillstand},
    \FullRef{topic:starke-blutung},
    \FullRef{topic:schock-ehi}
    \item    Kontrolle der Lebensfunktionen
    \\
    \noindent\FullRef{topic:lebensfunktionen-kontrolle}
    \item    Notfalldiagnose Bewusstlosigkeit
    \\
    \FullRef{topic:notfalldiagnose-bewusstlosigkeit}
    \item    Notfalldiagnose  Atem-Kreislaufstillstand
    \\
    \noindent\FullRef{topic:notfalldiagnose-atem-kreislaufstillstand}
    \item    Starke Blutung (Blutstillung, Fingerdruck,  Druckverband, Abbindung)
    \\
    \noindent\FullRef{topic:starke-blutung}
    \item    Schock (Schockbekämpfung,  Lagerungen)
    \\
    \noindent\FullRef{topic:schock-ehi}
    \item    Wunden und Wundverbände (Tierbisse, Insektenstiche, Verätzungen,
    Verbrennungen, Erfrierungen, Unterkühlung) 
    \\
    Wunden: \FullRef{topic:wunden-ehi} \\
    Wundverbände: \FullRef{topic:wundverbaende}
    \item    Quetschungen
    \\
    \FullRef{topic:wunden-einteilung}
    \item    Gelenksverletzungen
    \\
    Verstauchung: \FullRef{topic:verstauchung-eh},
    \\
    Verrenkung: \FullRef{topic:verrenkung-eh}
    \item
    Knochenbrüche, allgemein
    \\
    \noindent\FullRef{topic:knochenbruch-eh}
    \item
    Brustkorbverletzungen
    \\
    \FullRef{topic:thoraxtrauma}
    \item    Plötzlich auftretende Erkrankungen
    \\
    \FullRef{chp:MED}
    \item    Vergiftungen
    \\
    \noindent\FullRef{topic:vergiftungen-ehi}
    \item    Transport
    \\
    \FullRef{topic:rettungskette}
  \end{enumerate}
  \item \TabHeading{Hygiene}
  
  (Bestandteil des Sachgebiets 'Sanitätshilfe')
  
  \begin{inparadesc}
    \item[Stundenanzahl:] 2
    \item[Lehrkraft:] Arzt für Allgemeinmedizin, approbierter Arzt,
    Facharzt, Angehörige des gehobenen Dienstes für Gesundheits- und Krankenpflege,
    Hygienefachkraft,  Lehrsanitäter
  \end{inparadesc}
  
  \begin{enumerate}
    \item    Persönliche Hygiene
    \\
    \FullRef{topic:persoenliche-hygiene}
    \item    Grundbegriffe der Desinfektion
    \\
    \FullRef{topic:desinfektion-grundzuege}
    \item    Grundbegriffe der  Sterilisation
    \\
    \FullRef{topic:sterilisation-grundzuege}
    \item    Entsorgung von infektiösem Abfall
    \\
    \FullRef{topic:entsorgung-infektioeser-abfall}
    \item    Allgemeine Infektionslehre
    \\
    Erreger: \FullRef{topic:infektion-erreger},
    
    Übertragungswege: \FullRef{topic:infektion-uebertragungswege},
    
    Kreuzinfektion: \FullRef{topic:kreuzinfektion}.
    \item
    Vorgehen bei Verletzungen des Personals
    \\
    Nadelstichverletzungen: \FullRef{topic:nadelstichverletzung}
    \item    Hygienemaßnahmenplan
    \\
    \FullRef{topic:hygieneplan}
    \item    Infektionstransport
    \\
    \Abk{MRSA} und andere resistente Keime: \prettyref{m:mrsa-transport}
  \end{enumerate}
  \item \TabHeading{Berufsspezifische rechtliche Grundlagen}  ~
  
  
  
  \begin{inparadesc}
    \item[Stundenanzahl:] 3
    \item[Lehrkraft:]  Jurist, fachkompetente Person
  \end{inparadesc}
  
  \begin{enumerate}\item
    
    Aufgaben und Kompetenzen des  Rettungssanitäters
    \\
    \noindent\FullRef{topic:kompetenzen-sanitaeter},
    \FullRef{topic:fortbildungspflicht}
    \item    Dokumentation im Rettungswesen (Einsatzprotokoll, Leitstellendokumentation, Transportnachweis)
    \\
    \noindent\FullRef{topic:dokumentation-rechtlich}
    \item    Hilfs- und Rettungswesen
    \\
    \FullRef{topic:vetrag-rettungswesen}
    \item    Straßenverkehrsordnung
    \\
    \noindent\FullRef{topic:strassenverkehrsordnung}
    \item    Patientenrechte
    \\
    integrativ in \FullRef{topic:langer-dienst}
    %\noindent\FullRef{topic:patientenrechte}
    \item    Grundlagen des Haftungsrechtes
    \\
    \FullRef{topic:vetrag-rettungswesen},
    \FullRef{topic:strafbarkeitspruefung},
    \FullRef{topic:schadenersatz}
    %\noindent\FullRef{topic:haftungsrecht}
    \item    Unterbringungsgesetz
    \\
    \noindent\FullRef{topic:unterbringung-jus},
    \FullRef{topic:unterbringungsgesetz},
    \FullRef{topic:unterbringung}
    \item    Reversfähigkeiten und Effekten
    \\
    \FullRef{topic:jus-fall-aufklaerung}
    %\FullRef{topic:effekte}
    \item    Mitnahme von Begleitpersonen
    \\
    wird extern abgehandelt
    %\noindent\FullRef{topic:begleitpersonen-mitnahme}
  \end{enumerate}
  
  \item \TabHeading{Anatomie und Physiologie}
  
  (Bestandteil des Sachgebiets 'Sanitätshilfe')
  
  \begin{inparadesc}
    \item[Stundenanzahl:] 4
    \item[Lehrkraft:]  Notarzt, Arzt für Allgemeinmedizin, approbierter Arzt,
    Facharzt, Turnusarzt,   Angehörige  des  gehobenen  Dienstes für Gesundheits-
    und Krankenpflege, Lehrsanitäter, fachkompetente Person
  \end{inparadesc}
  
  \begin{enumerate}
    \item    Blutkreislauf - Grundzüge
    \\
    \FullRef{topic:blutkreislauf}
    \item    Gliedmaßen - Grundzüge
    \\
    Obere Extremität:
    \FullRef{topic:schulterguertel},
    \FullRef{topic:obere-extremitaet},
    \\
    untere Extremität:
    \FullRef{topic:beckenguertel},
    \FullRef{topic:untere-extremitaet}
    \item    Haut -- Grundzüge
    \\
    \FullRef{topic:haut}
    \item    Schädel und Rumpf - Grundzüge Skelett
    \\
    Schädel: \FullRef{topic:schaedel}
    \item    Brustkorb - Grundzüge
    \\
    Thorax: \FullRef{topic:brustkorb-skelett}, \\
    Herz: \FullRef{topic:herz} \\
    Atemwege/Lunge: \FullRef{topic:atemwege-lunge}\\
    Speiseröhre: \FullRef{topic:speiseroehre}
    \item    Bauchraum - Grundzüge
    \\
    Gastro-intestinaltrakt:
    \FullRef{topic:gastrointestinaltrakt}, \\ 
    andere Bauchorgane:
    \FullRef{topic:sonstige-bauchorgane}.
  \end{enumerate}
  
  \item \TabHeading{Störungen der  Vitalfunktionen  und Regelkreise und zu setzende Maßnahmen}
  
  (Bestandteil des  Sachgebiets  'Sanitätshilfe')
  
  \begin{inparadesc}
    \item[Stundenanzahl:]  8
    \item[Lehrkraft:]  Notarzt, Arzt für Allgemein- medizin, approbierter Arzt,
    Facharzt,  Lehr- sanitäter,  fachkompe- tente Person
  \end{inparadesc}
  
  \begin{enumerate}\item
    
    Definition Vitalfunktion
    \\
    \noindent\FullRef{topic:vitalfunktionen-definition},
    \FullRef{topic:vitalwerte}
    \item    NACA-Schema
    \\
    \FullRef{topic:naca}
    \item    Bewusstsein
    \\
    (Bewusstseinsstörungen, Bewusstlosigkeit)
    \\
    \noindent Bewusstsein
    allgemein: \FullRef{topic:bewusstsein};
    Bewusstseinsstörungen:
    \FullRef{topic:bewusstseinsstoerung},
    \\
    Bewusstlosigkeit: \FullRef{topic:bewusstseinsstoerung}
    \item
    Atmung
    \begin{enumerate}
      \item      (Atemstörungen,
      \\
      \noindent\FullRef{topic:atemstörungen},
      \FullRef{topic:respiratorische-notfaelle}
      \item      Atemstillstand)
      \\
      \noindent\FullRef{topic:atemstillstand},
      \FullRef{topic:atemstillstand-aed}
    \end{enumerate}
    \item    Kreislauf
    \begin{enumerate}
      \item      (Kreislaufstörungen,
      \\
      Blutdruck: \FullRef{topic:blutdruck},\\
      Herz-Kreislaufstörungen:
      \FullRef{topic:herz-kreislauf-stoerungen},\\ 
      Schock:
      \FullRef{topic:schock-ehi}, \FullRef{topic:schock}
    \end{enumerate}
    \item    Kreislaufstillstand)
    \\
    \FullRef{topic:kreislaufstillstand},
    \FullRef{topic:kreislaufstillstand-aed}
    \item    Regelkreise
    \begin{enumerate}
      \item      Wärmehaushalt
      \\
      \FullRef{topic:temperaturregulation}
      \item      Wasser- und Elektrolythaushalt
      \\
      \FullRef{topic:wasser-elektrolythaushalt}
      \item      \Säure-Basen-Haushalt
      \\
      \FullRef{topic:saeure-basen-haushalt}
      \item      Stoffwechsel
      \\
      \FullRef{topic:stoffwechsel}
    \end{enumerate}
    \item
    Starke Blutung
    \\
    \FullRef{topic:starke-blutung},
    \FullRef{topic:druckverband},
    \FullRef{topic:schock-hypovolaemischer}
    \item
    Schock (Ursachen,
    \FullRef{topic:schock}
    \begin{enumerate}
      \item      Wirkung, Schockformen,
      \\
      \FullRef{topic:schock}, \FullRef{topic:schock-arten}
      \item      Verlauf, Schockzeichen)
      \\
      \FullRef{topic:schock-allgemeine-symptomatik}
    \end{enumerate}
    \item Akute Störung der Atmung
    \begin{enumerate}
      \item      Atembehinderung
      \\
      \noindent\FullRef{topic:atemstörungen}
      \item      Verlegung der Atemwege  durch Fremdkörper
      \\
      \noindent\FullRef{topic:atemstörungen},
      \FullRef{topic:mechanische-atemwegsverlegung}
      \item      Verlegung durch Schwellung
      \\
      \noindent\FullRef{topic:atemstörungen},
      \FullRef{topic:asthma-bronchiale}
      \item      Akute Atemnot
      \\
      \noindent\FullRef{topic:atemstörungen},
      \FullRef{topic:respiratorische-notfaelle}
      \item      Absaugung
      \\
      \FullRef{topic:absaugung}
      \item      Sauerstoff (Inhalationsrichtlinien, Dosierungsvorschriften)
      \\
      Sauerstoff: \FullRef{topic:sauerstoff},\\
      Beatmungshilfen: \FullRef{topic:beatmung-hilfsmittel}
      \item      Assistierte Beatmung (Indikationen, Durchführung)
      \\
      \FullRef{topic:beatmung-assistiert}
    \end{enumerate}
    \item    Feststellung des Todes
    \\
    \FullRef{topic:tod}
  \end{enumerate}\item \TabHeading{Notfälle bei verschiedenen Krankheitsbildern  und zu setzende Maßnahmen}
  
  (Bestandteil des Sachgebiets 'Sanitätshilfe')
  
  \begin{inparadesc}
    \item[Stundenanzahl:]  6
    \item[Lehrkraft:]  Notarzt, Arzt für Allgemeinmedizin, approbierter Arzt,
    Facharzt, Turnusarzt, Lehrsanitäter, fachkompetente Person
  \end{inparadesc}
  \begin{enumerate}\item
    Koma aus vorerst unbekannter Ursache
    \begin{enumerate}
      \item      (Schlaganfall,
      \\
      \FullRef{topic:insult}
      \item      Meningitis,
      \\
      \FullRef{topic:meningitis}
      \item      Diabetes,
      \\
      \FullRef{topic:diabetes-mellitus}
      \item      Vergiftung)
      \\
      \FullRef{topic:vergiftung}
    \end{enumerate}
    \item
    Krampfanfall \\
    \noindent\FullRef{topic:zerebrale-krampfanfaelle},
    \FullRef{topic:krampfanfaelle-kinder},
    \FullRef{topic:Wundstarrkrampf} \\
    \begin{enumerate}
      \item      (Epilepsie, Tetanie)
      \\
      \noindent\FullRef{topic:zerebrale-krampfanfaelle}
    \end{enumerate}
    \item Akuter Gefäßverschluss an den Gliedmaßen
    \begin{enumerate}
      \item      Venenthrombose,
      \\
      \FullRef{topic:venenthrombose}
      \item      Arterielle Embolie)
      \\
      \FullRef{topic:arterieller-gefaessverschluss}
    \end{enumerate}
    \item Pulmonale Notfälle
    \begin{enumerate}
      \item      Asthma bronchiale
      \\
      \FullRef{topic:asthma-bronchiale}
      \item      Lungenödem
      \\
      \FullRef{topic:lungenoedem}
      \item      Lungenembolie
      \\
      \FullRef{topic:lungenembolie}
      \item      Lungenentzündung
      \\
      \FullRef{topic:lungenentzuendung}
    \end{enumerate}
    \item Cardiale Notfälle
    \begin{enumerate}
      \item
      (Herz-Rhythmus-Störungen,
      \\
      \FullRef{topic:herzrhythmusstoerungen}
      \item
      Herzversagen,
      Linksherzschwäche,
      Rechtsherzschwäche,
      \\
      \FullRef{topic:herzinsuffizienz},
      \FullRef{topic:linksherzinsuffizienz},
      \FullRef{topic:rechtsherzinsuffizienz}
      \item
      Angina pectoris,
      \\
      \FullRef{topic:angina-pectoris}
      \item
      Herzinfarkt,
      \\
      \FullRef{topic:herzinfarkt}
      \item
      Hochdruckkrise)
      \\
      \FullRef{topic:hypertensive-krise}, \FullRef{topic:hypertensiver-notfall}
    \end{enumerate}
    \item Allgemeinchirurgische Notfälle
    \begin{enumerate}
      \item
      Akutes Abdomen
      \\
      \FullRef{topic:akutes-abdomen}
      \item
      Gastritis, Ulcus
      \\
      \FullRef{topic:gastritis}
      \item  Gastroenteritis,
      \\
      \FullRef{topic:gastroenteritis}
      \item  Pankreatitis,
      \\
      \FullRef{topic:pankreatitis}
      \item   Appendizitis,
      \\
      \FullRef{topic:appendizitis}
      \item   Gallenkolik,
      \\
      \FullRef{topic:gallenkolik}
      \item   Ileus,
      \\
      \FullRef{topic:ileus}
      \item    Mesenterialinfarkt,
      \\
      \FullRef{topic:mesenterialinfarkt}
      \item  Lebensmittelvergiftung,
      \\
      \FullRef{topic:lebensmittelvergiftung}
      \item gastrointestinale Blutung)
      \\
      \FullRef{topic:gastrointestinale-blutung},\\
      \FullRef{topic:gastrointestinale-blutung},
      \FullRef{topic:oesophagusvarizenblutung}
    \end{enumerate}
    \item Gynäkologische und urologische Erkrankungen
    \begin{enumerate}
      \item Nierenkolik,
      \\
      \FullRef{topic:nierenkolik}
      \item Nierenbeckenentzündung,
      \\
      \FullRef{topic:nierenbeckenentzuendung}
      \item Harnwegsinfekt,
      \\
      \FullRef{topic:harnwegsinfekt}
      \item akutes Harnverhalten,
      \\
      \FullRef{topic:harnverhalten}
      \item chronische Niereninsuffienz und Hämodialyse,
      \\
      Niereninsuffizienz:
      \FullRef{topic:niereninsuffizienz};
      
      Elektrolytstörungen durch
      Dialyse: \FullRef{topic:dialyse-elektrolyte};
      
      Dialyse als Therapie bei chronischer
      Niereninsuffizenz: \FullRef{topic:cni-therapie-dialyse};
      
      Dialysepatienten, besondere
      Gefahren: \FullRef{topic:dialysepatienten-gefahren}.
      \item extrauterine Gravidität,
      \\
      \FullRef{topic:extrauterine-graviditaet}
      \item Abortus,
      \\
      \FullRef{topic:abort}
      \item  Unterleibsblutungen,
      \\
      \FullRef{topic:unterleibsblutungen}
      \item Ovarialtumore,
      \\
      fraglich
      \item Vergewaltigung (typische Verletzungen, psychische Betreuung)
      \\
      wird extern abgehandelt
      %\FullRef{topic:vergewaltigung}
    \end{enumerate}
  \end{enumerate}
  
  \item \TabHeading{Spezielle Notfälle und zu setzende Maßnahmen  }
  
  (Bestandteil des Sachgebiets  'Sanitätshilfe')
  
  \begin{inparadesc}
    \item[Stundenanzahl:]  15
    \item[Lehrkraft:]   Notarzt, Arzt für  Allgemein-  medizin,  approbierter
    Arzt, Facharzt,   Turnusarzt,  Lehr-  sanitäter,  fachkompetente Person
  \end{inparadesc}
  
  \begin{enumerate}
    \item Traumatologische Notfälle
    \begin{enumerate}
      \item      Schädel-Hirn-Verletzungen,
      \\
      \FullRef{topic:sht}, \FullRef{topic:sht-diagnostik}
      \item      Hirnblutung
      \\
      \FullRef{topic:hirnblutung},
      \FullRef{topic:sht-diagnostik}
      \item      Hirndruck
      \\
      Hirndruckzeichen: \FullRef{topic:hirnblutung},
      \FullRef{topic:sht-diagnostik};
      
      Hirnstammeinklemmung bei
      Hirndrucksteigerung: \FullRef{topic:hirnblutung},
      \FullRef{topic:sht-diagnostik};
      
      Ursachen:
      \FullRef{topic:hirnblutung}, \FullRef{topic:insult};
      
      Lagerung: \FullRef{topic:hirnblutung}.
      \item
      Halswirbelsäulen- und  Wirbelsäulentrauma
      \\
      
      \begin{enumerate}
      	\item      (Sturzhelmabnahme,
      	\\
      	\FullRef{topic:sturzhelmabnahme}	
      	\item      HWS-Scheine Beschreibung (MP): \\
     	 Material \FullRef{topic:hws-schiene-beschreibung}, \\
      	Anwendung \FullRef{topic:hws-schiene-anlegen}
      
      	\item     
      	Motorik-Durchblutung-Sensibilitätskontrolle,
      	\\
      	\FullRef{topic:sehen-hoeren-fuehlen},
      	\FullRef{topic:ersteinschaetzung}
      	\item      Body-Check, Umgang mit  Schaufeltrage und
      	Vakuummatratze, Sandwich-Technik)\\
      	Body-Check: \FullRef{topic:traumacheck}, \\
      	Material:
      	\FullRef{topic:schaufeltrage-beschreibung},
      	\FullRef{topic:vakuummatratze-beschreibung}, 
      	\\ Anwendung:
      	\FullRef{topic:schaufeltrage-anwendung}, 
      	\FullRef{topic:vakuummatratze-anwendung}
      \end{enumerate}
      
      \item Extremitätentrauma
      \FullRef{topic:extrimitaetentrauma}
      \begin{enumerate}
        
        \item (Verletzungsarten,
        \FullRef{topic:extrimitaetentrauma}
        
        \item Motorik-Durchblutung-Sensibilitätskontrolle,\\
        \FullRef{topic:sehen-hoeren-fuehlen},
        \FullRef{topic:ersteinschaetzung}
        
        \item Stiefelgriff,
        Prinzip der Schienung,
        Ruhigstellung,\\
        \FullRef{topic:vacuumbeinschiene-beschreibung},
        \FullRef{topic:vakuummatratze-anwendung},
        \FullRef{topic:frakturzeichen},
        \FullRef{topic:knochenbruch-eh},
        \FullRef{topic:verstauchung-tra},
        \FullRef{topic:verstauchung-eh},
        \FullRef{topic:verrenkung-eh}
        \item        Schienung des Armes
        \\
        \FullRef{topic:schienung-arm}
        \item        Schienung des Beines
        \\
        Material:
        \FullRef{topic:vacuumbeinschiene-beschreibung},\\
        Anwendung: 
        \FullRef{topic:vacuumbeinschiene-anwendung}
        \item        Pneumatische Schiene
        \\
        wird extern abgehandelt
        \item        Vakuumschiene
        \\
        Vakuumbeinschiene,\\
        Material:
        \FullRef{topic:vacuumbeinschiene-beschreibung},\\
        Anwendung: 
        \FullRef{topic:vacuumbeinschiene-anwendung}
        \item        Extensionsschiene
        \\
        wird extern abgehandelt
      \end{enumerate}
      \item Thoraxtrauma
      \begin{enumerate}
        \item
        offene, geschlossene Brustkorbverletzung
        \\
        \FullRef{topic:thoraxtrauma}
        \item        Serienrippenbruch
        \\
        \FullRef{topic:serienrippenfraktur}
        \item        Pneumothorax
        \\
        \FullRef{topic:pneumothorax}
        \item        Spannungspneumothorax
        \\
        \FullRef{topic:pneumothorax}
      \end{enumerate}
      \item      Bauchtrauma
      \begin{enumerate}
        \item        offene, geschlossene Bauchverletzung
        \\
        \FullRef{topic:bauchtrauma-offen},
        \FullRef{topic:bauchtrauma-geschlossen}
        \item        Verletzung der Harn- und
        Geschlechtsorgane
        \\
        wird extern abgehandelt
        \item        Beckentrauma
        \\
        \FullRef{topic:beckentrauma}
      \end{enumerate}
      \item      Polytrauma (Definition, Prioritäten, Management)
      \\
      \FullRef{topic:polytrauma}
      \item      Wunden (mechanische, chemische, thermische)
      \\
      \FullRef{topic:wunden}
      \item      Dekubitus Prophylaxe und Lagerung bei Dekubitus
      \\
      \FullRef{topic:dekubitusprophylaxe}
      \item      Verbandlehre
      \\
      \FullRef{topic:druckverband},
      \FullRef{topic:starke-blutung}
      \item      Akut auftretende Blutungen
      \\
      \FullRef{topic:starke-blutung},
      \FullRef{topic:druckverband},
      \FullRef{topic:schock-hypovolaemischer}
      \item      Vergiftung (Ursachen und Verdacht, Aufnahmearten)
      \\
      \FullRef{topic:vergiftungen-ehi},
      \FullRef{topic:vergiftung}
    \end{enumerate}
    \item Psychiatrische Notfälle
    \begin{enumerate}
      \item      (Suizid,
      \\
      \FullRef{topic:suizid1}, \FullRef{topic:suizid2}
      \item      Psychose
      \\
      \FullRef{topic:psychose}
      \item      Suchterkrankungen und Entzugssyndrom
      \\
      \FullRef{topic:alkohol-entzug}
      %\FullRef{topic:sucht},
      %\FullRef{topic:entzugssysndrom}
      \item      Depression
      \\
      \FullRef{topic:depression}
      \item      Manie
      \\
      \FullRef{topic:manie}
    \end{enumerate}
    \item
    Schwangerschaft und Geburt
    \begin{enumerate}
      \item      Notfälle in der Schwangerschaft
      \\
      \FullRef{topic:notfall-fruehschwangerschaft}, 
      \FullRef{topic:notfall-spaetschwangerschaft}
      \item      Geburt
      \\
      \FullRef{topic:geburt}
      \item      Geburtskomplikationen
      \\
      \FullRef{topic:notfall-geburt}
      \item      Versorgung des Neugeborenen\\
      \FullRef{topic:geburt}
    \end{enumerate}
    \item Notfälle im Säuglings- und Kleinkindalter
    \begin{enumerate}
      \item      anatomische und physiologische Besonderheiten
      \\
      \FullRef{topic:kind-besonderheiten}
      \item      Kontrolle der Lebensfunktionen und lebensrettende Sofortmaßnahmen
      \\
      \FullRef{topic:pediatric-life-support}
      \item      Krampfanfälle
      \\
      Allgemein: \FullRef{topic:zerebrale-krampfanfaelle},\\
      Fieberkrampf: \FullRef{topic:krampfanfaelle-kinder},
      \\
      Wundstarrkrampf: \FullRef{topic:Wundstarrkrampf}
      \item      Pseudokrupp
      \\
      \FullRef{topic:laryngitis}
      \item      Epiglottitis
      \\
      \FullRef{topic:epiglottitis}
      \item      Keuchhusten
      \\
      \FullRef{topic:}
      \item      plötzlicher Kindstod
      \\
      \FullRef{topic:sids}
    \end{enumerate}
  \end{enumerate}
  \item \TabHeading{Defibrillation mit halbautomatischen Geräten}
  
  (Bestandteil des Sachgebiets  'Sanitätshilfe')
  
  \begin{inparadesc}
    \item[Stundenanzahl:]   8
    \item[Lehrkraft:]   Notarzt,  Arzt für Allgemeinmedizin,  approbierter Arzt,
    Facharzt, Lehrsanitäter, fachkompetente Person
  \end{inparadesc}
  
  \begin{enumerate}
    \item Der halbautomatische Defibrillator\\
    \FullRef{chp:AED}
    \item Handhabung eines halbautomatischen Defibrillators\\
    \FullRef{chp:AED}
    \item Gerätemanagement während der Reanimation\\
    \FullRef{chp:AED}
    \item Erfolgskontrolle\\
    \FullRef{chp:AED}
  \end{enumerate}
  \item \TabHeading{Gerätelehre und Sanitätstechnik }
  
  \begin{inparadesc}
    \item[Stundenanzahl:]  12
    \item[Lehrkraft:]   Lehrsanitäter, fachkompetente Person
  \end{inparadesc}
  
  \begin{enumerate}\item
    
    Medizinproduktegesetz -- Information
    \\
    \FullRef{topic:mpg-info}
    \item    Bergungs- und Lagerungstechniken
    \begin{enumerate}
      \item      Bergetuch
      \\
      \FullRef{topic:bergetuch}
      \item      Einheitskrankentrage
      \\
      \FullRef{topic:fahrtrage}
      \item      Tragsessel
      \\
      \FullRef{topic:tragsessel}
      \item      Fahrtrage
      \\
      \FullRef{topic:fahrtrage},
      \FullRef{topic:lagerungsarten}
      \item      Rollstuhl
      \\
      \FullRef{topic:rollstuhl}
    \end{enumerate}
    \item    Einsatzfahrzeug
    \\
    \FullRef{topic:einsatzmittel}
    \item    Beatmungsbeutel
    \\
    \FullRef{topic:beatmungsbeutel}
    \item    Absauggeräte
    \\
    \FullRef{topic:absaugung}
    \item    Sauerstoff
    \\
    \FullRef{topic:sauerstoff}
    \item    Infusionen und
    \\
    \FullRef{topic:infusion}, \FullRef{topic:venflon}
    \item    Infusionsgeräte
    \\
    \FullRef{topic:infusion}
    \item    Blutdruckmessung
    \\
    \FullRef{topic:rr-messung}
    \item    Stabilisierungs- und Schienungstechniken\\
    Material: \FullRef{topic:schienungsmaterial}, \\
    Anwendungen: \FullRef{topic:schienung}
    \begin{enumerate}
      \item      Stabilisierung der Halswirbelsäule
      \\
      Material \FullRef{topic:hws-schiene-beschreibung}, \\
      Anwendung \FullRef{topic:hws-schiene-anlegen}
      \item      Schaufeltrage
      \\
      Material: \FullRef{topic:schaufeltrage-beschreibung},
      \\ Anwendung: \FullRef{topic:schaufeltrage-anwendung}
      \item      Vakuummatratze
      \\
      Material: \FullRef{topic:vakuummatratze-beschreibung},
      \\ Anwendung:
      \FullRef{topic:vakuummatratze-anwendung}
    \end{enumerate}
    \item    Geburtenkoffer
    \\
    wird extern abgehandelt
    \item    Transportinkubator
    \\
    wird extern abgehandelt
  \end{enumerate}
  
  \item \TabHeading{Rettungswesen}
  
  \begin{inparadesc}
    \item[Stundenanzahl:]  4
    \item[Lehrkraft:]   Fachkompetente Person
  \end{inparadesc}
  
  \begin{enumerate}\item
    
    Rechtliche Grundlagen
    \\
    \FullRef{topic:rettungswesen-recht}
    \item    Zusammenarbeit mit  anderen Organisationen
    \\
    \FullRef{topic:rettungswesen-wien}
    \item    Einsatzarten
    \\
    \FullRef{topic:einsatzarten}
    \item    Transport- und Fahrzeugarten (Land,  Wasser, Luft)
    \\
    \FullRef{topic:einsatzmittel}
    \item    Normen, persönliche Schutzausrüstung
    \\
    wird extern abgehandelt
    \item    Fahrzeugausstattung (Dienstvorschriften)
    \\
    wird extern abgehandelt (TD/Fahrzeugeinschulung)
    \item    Rettungskette
    \\
    \FullRef{topic:rettungskette}
    \item    Hilfsfrist
    \\
    wird extern abgehandelt
    \item    Dienststellennetz
    \\
    \FullRef{topic:rettungswesen-wien}
    \item    Personal im Rettungsdienst
    \\
    \FullRef{chp:TAE}
    \item    Notarztsysteme
    \\
    \FullRef{topic:rettungswesen-wien},
    \FullRef{topic:einsatzmittel}
    \item    Leitstelle
    \\
    \FullRef{topic:leitstelle}
    \item    Kommunikationsarten
    \\
    \FullRef{topic:kommunikationswege}
    \item    Gefahren an der Einsatzstelle
    \\
    \FullRef{topic:gefahrenzonen}
    \item    Gefahrguteinsätze
    \\
    \FullRef{topic:gefahrengutunfall}
    \item    Sondertransporte
    \\
    wird extern abgehandelt
    
  \end{enumerate}
  \item \TabHeading{Katastrophen, Großschadensereignisse, Gefahrgutunfälle}
  
  \begin{inparadesc}
    \item[Stundenanzahl:]  4
    \item[Lehrkraft:]  Fachkompetente Person
  \end{inparadesc}
  
  \begin{enumerate}\item
    
    Katastrophen
    \begin{enumerate}
      \item
      Rechtliche Grundlagen,  Geltungsbereiche
      \\
      \FullRef{topic:katastrophe},
      \FullRef{topic:katastrophenhilfsgesetze}
      \item      Arten der Katastrophen
      \\
      \FullRef{topic:katastrophe},
      \FullRef{topic:katastrophenschutz}
      \item      Phasen der Katastrophenbewältigung
      \\
      \FullRef{topic:katastrophe-phasen}
      \item      Katastrophenhilfseinheiten
      \\
      \FullRef{topic:katastrophenschutz}
      \item      Führungsorganisation
      \\
      \FullRef{topic:katastrophe-fuehrung},
      \FullRef{topic:sanhist}
      \item      personelle, materielle und finanzielle Vorsorge
      \\
      \FullRef{topic:katastrophe-material}
      \item      Einsatzgrundsätze
      \\
      \FullRef{topic:einsatzgrundsaetze}
      \item      generelle Einsatzrichtlinien
      \\
      \FullRef{topic:einsatzrichtlinie}
      \item      Grundzüge der Triage
      \\
      \FullRef{topic:triage}
    \end{enumerate}
    \item
    Großschadensereignisse
    \begin{enumerate}
      \item
      Rechtliche Grundlagen
      \\
      \FullRef{topic:grossschaden}
      \item      Einstufung
      \\
      \FullRef{topic:grossschaden}
      \item      Alarmierung
      \\
      \FullRef{topic:alarmierung}
      \item      Schadensraum
      \\
      \FullRef{topic:sanhist}
      \item      Schadensplatz, Sicherheitseinrichtungen
      \\
      \FullRef{topic:sanhist}
      \item      Organisation beim Großunfall
      \\
      \FullRef{topic:sanhist}
      \item      Grundzüge der Triage
      \\
      \FullRef{topic:triage}
      \item      österreichisches Patientenleitsystem
      \\
      \FullRef{topic:pls}
      \item      Material und Ausrüstung
      \\
      \FullRef{topic:katastrophe-material}
      \item      Kommunikation
      \\
      \FullRef{topic:khd-organisation}
    \end{enumerate}
    \item
    Gefahrgutunfälle
    \begin{enumerate}
      \item      Arten von Gefahrgutunfällen
      \\
      \FullRef{topic:gefahrenzonen}
      \item      Gefahrzettel
      \\
      \FullRef{topic:gefahrzettel}
      \item      Gefahrsymbole
      \\
      \FullRef{topic:gefahrsymbole}
      \item      Warntafel
      \\
      \FullRef{topic:warntafel},
      \FullRef{topic:gefahrnummer}
      \item
      Verhalten am Unfallort
      \\
      \FullRef{topic:einsatzrichtlinie}
      \item      Koordination mit anderen Einsatzorganisationen
      \\
      wird extern abgehandelt
      \item      Absperrmaßnahmen
      \\
      \FullRef{topic:absperrmassnahmen}
      \item      Sofortmaßnahmen
      \\
      \FullRef{topic:einsatzrichtlinie}
    \end{enumerate}
  \end{enumerate}\item \TabHeading{Angewandte Psychologie und Stressbewältigung}
  
  
  
  \begin{inparadesc}
    \item[Stundenanzahl:]  4
    \item[Lehrkraft:]  Fachkompetente Person
  \end{inparadesc}
  
  \begin{enumerate}\item
    
    Belastung, Anforderung, Beanspruchung, workflow
    \\
    \FullRef{topic:belastung}
    \item    Überforderung, Unterforderung
    \\
    wird extern abgehandelt (PEER-Vortrag)
    \item    Beanspruchungsfolgen
    \\
    wird extern abgehandelt (PEER-Vortrag)
    \item    Stressursachen, -entstehung und -faktoren
    \\
    \FullRef{topic:stress}
    \item    Stressauswirkung
    \\
    \FullRef{topic:stress}
    \item    Früherkennung
    \\
    \FullRef{topic:stress}
    \item    Grundsätze der Stressvermeidung
    \\
    \FullRef{topic:stress}
    \item    Maßnahmen zur Verhütung und Verminderung von Beanspruchungsfolgen
    \\
    wird extern abgehandelt (PEER-Vortrag)
    \item    Psychische Betreuung von Kranken/Verletzten
    \\
    \FullRef{topic:kommunikationsregeln},
    \FullRef{topic:umgang-psy}
    \item    Gesprächsführung
    \\
    \FullRef{topic:kommunikationsregeln},
    \FullRef{topic:umgang-psy}
    \item    Vertrauensaufbau und Patienteninformation
    \\
    \FullRef{topic:kommunikationsregeln},
    \FullRef{topic:umgang-psy}
    \item    psychische Belastungssyndrome
    \\
    \FullRef{topic:ptss}
    \item    verwirrte Patienten
    \\
    wird integrativ in allen Kapitel abgehandelt
    \item    Begleitung und Betreuung Sterbender
    \\
    wird extern abgehandelt
    \item    Supervision
    \\
    wird extern abgehandelt (PEER-Vortrag)
  \end{enumerate}
  \item \TabHeading{Praktische Übungen ohne Patientenkontakt}
  
  \begin{inparadesc}
    \item[Stundenanzahl:] 16
    \item[Lehrkraft:] Lehrsanitäter, fachkompetente Person
  \end{inparadesc}
  
  \begin{enumerate}
    \item Regloser Notfallpatient
    \\
    \TakeNotes
    \item    Kontrolle der  Lebensfunktionen  (erwachsener Notfallpatient)
    \\
    \TakeNotes
    \item    Notfalldiagnose Bewusstlosigkeit (stabile Seitenlage)
    \\
    \TakeNotes
    \item    Notfalldiagnose Atemstillstand (Beatmung)
    \\
    \TakeNotes
    \item    Notfalldiagnose Kreislaufstillstand (Beatmung und Herzmassage)
    \\
    \TakeNotes
    \item    Blutstillung (Fingerdruck, Abdrückstellen, Druckverband, Abbindung, Amputatsversorgung)
    \\
    \TakeNotes
    \item    Blutdruckmessung
    \\
    \TakeNotes
    \item    Schockbekämpfung (Lagerungsarten)
    \\
    \TakeNotes
    \item
    Halswirbelsäulen- und Wirbelsäulentrauma
    \\
    \begin{enumerate}
      \item      (Sturzhelmabnahme,
      \\
      \TakeNotes
      \item      Halswirbelsäulenschienung,
      \\
      \TakeNotes
      \item      Motorik-Durchblutung-Sensibilitätskontrolle,
      \\
      \TakeNotes
      \item      Body-Check,
      \\
      \TakeNotes
      \item      Umgang mit Schaufeltrage,
      \\
      \TakeNotes
      \item      Umgang mit Vakuummatratze,
      \\
      \TakeNotes
      \item      Sandwich-Technik)
      \\
      \TakeNotes
    \end{enumerate}
    \item
    Extremitätentrauma
    \begin{enumerate}
      \item
      (Motorik-Durchblutung-Sensibilitätskontrolle,
      \\
      \TakeNotes
      \item      Stiefelgriff,
      \\
      \TakeNotes
      \item      Ruhigstellung,
      \\
      \TakeNotes
      \item      Schienung des Armes,
      \\
      \TakeNotes
      \item      Schienung des Beines,
      \\
      \TakeNotes
      \item      pneumatische Schiene,
      \\
      \TakeNotes
      \item      Vakuumschiene,
      \\
      \TakeNotes
      \item      Extensionsschiene)
      \\
      \TakeNotes
    \end{enumerate}
    \item
    Verbandlehre
    \\
    \TakeNotes
    \item Notfälle im Säuglings- und Kleinkindalter
    \begin{enumerate}
      \item
      (Kontrolle der Lebensfunktionen und lebensrettende Sofortmaßnahmen)
      \\
      \TakeNotes
    \end{enumerate}
    \item An- und Auskleiden, Körperpflege und Hygiene, Nahrungs- und Flüssigkeitsaufnahme, Harn- und Stuhlentleerung,
    Erbrechen
    \\
    \TakeNotes
    \item Ergonomische und schonende Arbeitsweise (Richtiges Heben und Tragen)
    \\
    \TakeNotes
    \item Handhabung der in Einsatzfahrzeugen zu verwendenden Geräte  (insbesondere
    Krankentrage, Tragsessel,  Sauerstoffgeräte, Kommunikationseinrichtungen) sowie
    die  Handhabung von Rollstühlen und Gehhilfen
    
    
  \end{enumerate}
\end{enumerate}

\end{multicols}

\end{WidePar}

\section{Berücksichtigte Lehrmeinungen}

\subsection{Arbeiter-Samariter-Bund Österreichs}

\subsubsection{Bundesverband}

\begin{WidePar}
\small
\begin{tabulary}{\linewidth}{LcLcL}
\toprule\addlinespace
\noindent\TabHeading{Kurztitel}
&
\noindent\TabHeading{Nr.}
&
\noindent\TabHeading{Referenz}
&
\noindent\TabHeading{seit}
&
\noindent\TabHeading{Querverweise}
\\\addlinespace\midrule\addlinespace
EH Reanimation%KT
&
01/2011%Nr
&
\bibentry{AsboeLehrmeinungEh2011-01} \cite{AsboeLehrmeinungEh2011-01}
%bibentry,cite
&
0.18.0%seit
&
\FullPrettyRef{topic:reanimation-eh}%FullPrettyRef
\\\addlinespace
EH Helmabnahme%KT
&
02/2011%Nr
&
\bibentry{AsboeLehrmeinungEh2011-02} \cite{AsboeLehrmeinungEh2011-02}%bibentry,cite
&
%seit
&
%FullPrettyRef
\\\addlinespace
RD Reanimation%KT
&
02/2011%Nr
&
\bibentry{AsboeLehrmeinungRd2011-02} \cite{AsboeLehrmeinungRd2011-02}%bibentry,cite
&
0.18.0%seit
&
\prettyref{chp:BLS}, \prettyref{chp:ALS}%FullPrettyRef
\\\addlinespace
EH Verbrennung%KT
&
03/2011%Nr
&
\bibentry{AsboeLehrmeinungEh2011-03} \cite{AsboeLehrmeinungEh2011-03}%bibentry,cite
&
%seit
&
%FullPrettyRef
\\\addlinespace
RD Nitrogabe%KT
&
03/2011%Nr
&
\bibentry{AsboeLehrmeinungRd2011-03} \cite{AsboeLehrmeinungRd2011-03}%bibentry,cite
&
%seit
&
%FullPrettyRef
\\\addlinespace
RD \ce{O2} bei MCI und ROSC%KT
&
04/2011%Nr
&
\bibentry{AsboeLehrmeinungRd2011-04} \cite{AsboeLehrmeinungRd2011-04}%bibentry,cite
&
%seit
&
%FullPrettyRef
\\\addlinespace
RD Not-/ schnelle Rettung%KT
&
05/2011%Nr
&
\bibentry{AsboeLehrmeinungRd2011-05} \cite{AsboeLehrmeinungRd2011-05}%bibentry,cite
&
%seit
&
%FullPrettyRef
\\\addlinespace
RD Rautekgriff%KT
&
06/2011%Nr
&
\bibentry{AsboeLehrmeinungRd2011-06} \cite{AsboeLehrmeinungRd2011-06}%bibentry,cite
&
0.18.0%seit
&
\FullPrettyRef{topic:rautekgriff}%FullPrettyRef
\\\addlinespace
Sturzhelmabnahme
&
07/2011
&
\bibentry{AsboeLehrmeinungRd2011-07} \cite{AsboeLehrmeinungRd2011-07,ASBOEAkademieUpdate2010Sturzhelm}
&
0.16.0
&

\\\addlinespace
RD \ce{O2} allgem.%KT
&
09/2011%Nr
&
\bibentry{AsboeLehrmeinungRd2011-09} \cite{AsboeLehrmeinungRd2011-09}%bibentry,cite
&
0.18.0%seit
&
\FullPrettyRef{m:sauerstoffberieselung}%FullPrettyRef
\\\addlinespace
Spineboard%KT
&
10/2011
&
\bibentry{AsboeLehrmeinungRd2011-10} \cite{AsboeLehrmeinungRd2011-10}%bibitem
&
 0.18.0
&
%FullPrettyRef
\\\addlinespace
RD \ce{O2}-Gabe Asthma%KT
&
11/2011%Nr
&
\bibentry{AsboeLehrmeinungRd2011-11} \cite{AsboeLehrmeinungRd2011-11}%bibentry,cite
&
0.18.0%seit
&
\FullPrettyRef{m:asthmaanfall}%FullPrettyRef
\\\addlinespace
Boa-Rettung%KT
&
14/2011
&
\bibentry{AsboeLehrmeinungRd2011-14} \cite{AsboeLehrmeinungRd2011-14}%bibitem
&
ausständig
&
%FullPrettyRef
\\\addlinespace
RD Guedel-Tubus%KT
&
15/2011%Nr
&
\bibentry{AsboeLehrmeinungRd2011-15} \cite{AsboeLehrmeinungRd2011-15}%bibentry,cite
&
%seit
&
%FullPrettyRef
\\\addlinespace
RD Larynxtubus%KT
&
16/2011%Nr
&
\bibentry{AsboeLehrmeinungRd2011-16} \cite{AsboeLehrmeinungRd2011-16}%bibentry,cite
&
%seit
&
%FullPrettyRef
\\\addlinespace
Absaugen Neugeborener%KT
&
18/2011
&
\bibentry{AsboeLehrmeinungRd2011-18} \cite{AsboeLehrmeinungRd2011-18}%bibitem
&
0.18.0%seit
&
\FullPrettyRef{m:neugeborenes}%FullPrettyRef
\\\addlinespace
Abnabelung%KT
&
19/2011%Nr
&
\bibentry{AsboeLehrmeinungRd2011-19} \cite{AsboeLehrmeinungRd2011-19}%bibitem,cite
&
%seit
&
%FullPrettyRef
\\\addlinespace
Geburt, Dammschutz%KT
&
20/2011%Nr
&
\bibentry{AsboeLehrmeinungRd2011-20} \cite{AsboeLehrmeinungRd2011-20}%bibentry,cite
&
0.18.0%seit
&
\ref{topic:dammschutz}, \ref{m:geburt}%FullPrettyRef
\\\addlinespace
Verbrennung
&
21/2011
&
\bibentry{AsboeLehrmeinungRd2011-21} \cite{AsboeLehrmeinungRd2011-21}

\bibentry{ASBOEAkademieUpdate2010Verbrennung} \cite{ASBOEAkademieUpdate2010Verbrennung}
&
0.16.0
&
\FullPrettyRef{m:verbrennung-eh}, \FullPrettyRef{m:verbrennung-leicht}, \FullPrettyRef{m:verbrennung-schwer}
\\\addlinespace
Notfallkompetenzen%KT
&
23/2011%Nr
&
\bibentry{AsboeLehrmeinungRd2011-23} \cite{AsboeLehrmeinungRd2011-23}%bibitem,cite 
&
in Arbeit%seit
&
\ref{m:hypertensiver-notfall}%FullPrettyRef
\\\addlinespace
%KT
&
%Nr
&
%bibentry,cite
&
%seit
&
%FullPrettyRef
\\\addlinespace
%KT
&
%Nr
&
%bibentry,cite
&
%seit
&
%FullPrettyRef
\\\addlinespace\bottomrule
\end{tabulary}


\end{WidePar}


\subsubsection{Landesverband Wien}

\subsubsection{Gruppe Floridsdorf-Donaustadt}


%%%%%%%%%%%%%%%%%%%%%%%%%%%%%%%%%%%%%%%%%%%%%%%%%%%%%%%%%%%
%%% Kommentare
\printnotes

%%%%%%%%%%%%%%%%%%%%%%%%%%%%%%%%%%%%%%%%%%%%%%%%%%%%%%%%%%%
%%% Literaturverzeichnis
% \opt{STATUS-WITH-BIBLIO}{{\loadgeometry{FullPage}\twocolumn\small\setlength{\parskip}{0pt}\cleardoublepage\bibliography{Literatur-AASS}\restoregeometry}\onecolumn}
% \opt{STATUS-WITH-BIBLIO}{{\setlength{\parskip}{0pt}\bibliography{Literatur-AASS}}}

%%%%%%%%%%%%%%%%%%%%%%%%%%%%%%%%%%%%%%%%%%%%%%%%%%%%%%%%%%%
%%% Index
% \clearpage
{\small\printindex}




\chapter{Übersichten}

% \SectionUnderConstruction


\Layout{57}{}{}{}{
\begin{tabularx}{\linewidth}{llXXXXXXX}%
\noindent
&
&
\noindent\TabHeading{NG}
&
\noindent\TabHeading{Sg}
&
\noindent\TabHeading{KK}
&
\noindent\TabHeading{SK}
&
\noindent\TabHeading{Ju}
&
\noindent\TabHeading{Erw \Male}
&
\noindent\TabHeading{Erw \Female}
\\\addlinespace\midrule\addlinespace
\TabSub{RR}
&
[\MetricUnit{mmHg}]
&
$\hookrightarrow$
&
$\hookrightarrow$
&
$\hookrightarrow$
&
$\hookrightarrow$
&
$\hookrightarrow$
&
\VonBis{100}{140}
&
$\hookleftarrow$
\\\addlinespace
\TabSub{HF} 
&
[\PerMin]
&
\VonBis{140}{180}
&
\VonBis{110}{160}\nocite{Lissauer2007}
&
\VonBis{95}{140}\nocite{Lissauer2007}
&
\VonBis{80}{120}\nocite{Lissauer2007}
&
\VonBis{60}{100}\nocite{Lissauer2007}
&
\VonBis{60}{100}
&
$\hookleftarrow$
\\\addlinespace\midrule\addlinespace
\TabSub{AF} 
&
[\PerMin]
&
\VonBis{30}{50}\nocite{Lissauer2007}
&
\VonBis{20}{30}\nocite{Lissauer2007}
&
\VonBis{20}{30}\nocite{Lissauer2007}
&
\VonBis{15}{20}\nocite{Lissauer2007}
&
\VonBis{14}{20}
&
\VonBis{12}{16}
&
$\hookleftarrow$
% \\\addlinespace
% \TabSub{AZV} 
% &
% [\nicefrac{\Ml}{\Kg}]
% &
% 
% &
% 
% &
% 
% &
% 
% &
% 
% &
% 
% &
% 
\\\addlinespace
\TabSub{AZV}\footnote{Bezogen auf Normalgewicht, \E{nicht}
auf Ist-Gewicht!} 
&
[\Ml]
&
\VonBis{20}{30}\nocite{LutomskyLeitRett3}
&
40\,--\,55\,--\,80\nocite{LutomskyLeitRett3}
&
\VonBis{80}{180}\nocite{LutomskyLeitRett3}
&
\VonBis{240}{350}\nocite{LutomskyLeitRett3}
&
500\nocite{LutomskyLeitRett3}
&
800\nocite{LutomskyLeitRett3}
&
700\nocite{LutomskyLeitRett3}
% \\\addlinespace
% \TabSub{AMV} 
% &
% [\nicefrac{\Liter}{\Min}]
% &
% \VonBis{600}{1500}
% &
% 
% &
% 
% &
% 
% &
% 
% &
% 
% &
% 
\\\addlinespace\midrule\addlinespace
\TabSub{BZ}\footnote{nüchtern}
&
[\MgDl]
&
$\hookrightarrow$
&
$\hookrightarrow$
&
$\hookrightarrow$
&
$\hookrightarrow$
&
$\hookrightarrow$
&
\VonBis{80}{120}
&
$\hookleftarrow$
\\\addlinespace
\TabSub{Sp\ce{O2}}
&
[\PC]
&
$\hookrightarrow$
&
$\hookrightarrow$
&
$\hookrightarrow$
&
$\hookrightarrow$
&
$\hookrightarrow$
&
\VonBis{95}{100}
&
$\hookleftarrow$
\\\addlinespace
\TabSub{Temp}\footnote{Körperkerntemperatur}
&
[\textcelsius]
&
$\hookrightarrow$
&
$\hookrightarrow$
&
$\hookrightarrow$
&
$\hookrightarrow$
&
$\hookrightarrow$
&
37
&
$\hookleftarrow$
\\\addlinespace\bottomrule
\end{tabularx}
}{Übersicht Normalwerte}{}






% \begin{table}\FontTableBase
% \caption{SI Basiseinheiten.}
% \label{tab:unit:base}
% \centering
% \begin{tabular}{llll}
% \toprule
% \TabHeading{Maß} & \Z{\TabHeading{Symbol}} & \TabHeading{Einheit} & \TabHeading{Symbol} \\
% \midrule
% Stromfluß 	& $I$, $i$ &  Ampere 		& \MetricUnit{A} \\
% 			& $I_v$ & \Z{Candela} 	& \Z{\MetricUnit{cd}} \\
% Temperatur 	& $T$ & Kelvin 		& \MetricUnit{K} \\
% Masse 		& $m$ & Kilogramm 	& \MetricUnit{kg} \\
% Weg 		& & Meter 		& \MetricUnit{m} \\
% Stoffmenge 	& $n$ & Mol 			& \MetricUnit{mol} \\
% Zeit 		& $t$ & Sekunde 		& \MetricUnit{s} \\
% \bottomrule
% \end{tabular}
% \end{table}

\Layout{57}{}{}{}{
\begin{tabulary}{\linewidth}{llllllL}
\toprule\addlinespace
\TabHeading{Maß}		& \TabHeading{Symbol}& \TabHeading{Einheit}	& \TabHeading{Zeichen}& \TabHeading{Typ}& \TabHeading{Anmerkungen, Anwendungen} & \TabHeading{}\\\addlinespace\midrule\addlinespace
\EE{Stromfluß} 			& $I$, $i$ 	&  \EE{Ampere} 		& \MetricUnit{A} 		& SI Basis \\
\EE{}					& $I_v$ 	& \Z{Candela} 	& \Z{\MetricUnit{cd}}  	& SI Basis \\
\EE{Temperatur} 		& $T$ 		& \EE{Kelvin} 		& \MetricUnit{K}  		& SI Basis \\
\EE{Masse} 				& $m$ 		& \EE{Kilogramm} 	& \MetricUnit{kg} 		& SI Basis  \\
\EE{Weg} 				& $s$		& \EE{Meter} 		& \MetricUnit{m}  		& SI Basis \\
\EE{Stoffmenge} 		& $n$ 		& \EE{Mol} 			& \MetricUnit{mol}  	& SI Basis \\
\EE{Zeit} 				& $t$ 		& \EE{Sekunde} 		& \MetricUnit{s}  		& SI Basis \\\addlinespace\midrule\addlinespace
\Z{Absorptionsdosis}	&  			& \Z{Gray} 		&  \Z{\MetricUnit{Gy}} 			& \\
% \A{???}					& $I_v$ 	& \Z{Candela} 	& \Z{\MetricUnit{cd}} 			& SI 	& \\
Dosisäquivalent			&  			& Sievert 		& \MetricUnit{Sv} 			& 	& Strahlenbelastung \\\addlinespace
Druck 					& $p$ 		& \T{Bar} 		& \MetricUnit{bar} \\
Druck 					& $p$		& Pascal 		& \MetricUnit{Pa} \\
Druck 					& $p$		&\T{Millimeter Quecksilbersäule} & \MetricUnit{mmHg} & & Blutdruck\\
Druck 					& $p$		&\T{Zentimeter Wassersäule} & \MetricUnit{cm\ce{H2O}} \\\addlinespace
% Elektr. Leitfähigkeit	&  			& Siemens 		& \MetricUnit{S} \\
Elektr. Ladung			& 			& \Z{Coulomb}	&  \Z{\MetricUnit{C}}   		& \\
Elektr. Spannung 		& $U$		& Volt 			& \MetricUnit{V} \\
Elektr. Widerstand 		& $R$		& Ohm 			& \MetricUnit{$\Omega$} 		& \\\addlinespace
Energie, Arbeit			& 			& \T{Joule}		&  \MetricUnit{J}     			& \\
Frequenz				&  			& Hertz 		&  \MetricUnit{Hz}     			& \\
% \Z{Kapazität} 			&  			& \Z{Farad} 	&  \Z{\MetricUnit{F}}    		& \\
% \A{???}	 				&  			& \Z{Katal} 	& \Z{\MetricUnit{kat}}  		& \\
% \Z{Indukktivität} 		&  			& \Z{Henry} 	&  \Z{\MetricUnit{H}}    		& \\
\Z{Geschwindigtkeit}	&			& \Z{Lichtgeschwindigkeit}	& $c_0$ 				& \\\addlinespace
Leistung 				& 			& Watt 			& \MetricUnit{W} \\\addlinespace
% \A{???}					&  			& \Z{Lumen} 	& \Z{\MetricUnit{lm}}   		& \\
% \A{???}					&  			& \Z{Lux} 		& \Z{\MetricUnit{lx}}   		& \\
Magn. Feldstärke 		& 			&  Tesla 		& \MetricUnit{T} & & MRT\\\addlinespace
% Magn. Fluß 				&  			& \Z{Weber} 	& \Z{\MetricUnit{}} \\
Masse					& 			& Dalton 		& \MetricUnit{Da}       		& & Angabe der Molekulargröße bei kolloidalen Infusionslösungen \\
Masse					& $m$		& Tonne			& \MetricUnit{t}\\\addlinespace
Schalldruck 			&  			& Dezibel		& \MetricUnit{dB}\\\addlinespace
Stoffmenge 				& $n$ 		& Mol 			& \MetricUnit{mol} 				& 	& Stoffmenge bei Infusionszubereitung\\\addlinespace
Temperatur 				& $T$		& Grad Celsius 	& \MetricUnit{\textcelsius} 	& 	& \\
Temperatur 				& $T$ 		& Kelvin 		& \MetricUnit{K} 				& 	& \\\addlinespace
Volumen 				&			& Liter			& \MetricUnit{l}, {\Liter}\\\addlinespace
Winkel 					&			& Grad 			& {\textdegree}\\
Winkel 					&			& Winkelminute 	& \\
Winkel 					&			& Winkelsekunde & \\\addlinespace
Zeit					& $t$		& Tag 			& {\Day}\\
Zeit					& $t$		& Stunde 			& {\Hour}\\
Zeit					& $t$		& Minute (Zeit) 	& {\Min}\\\addlinespace



% {astronomical unit}{astronomicalunit} &  & \MetricUnit{} \\
% {atomic mass unit}{atomicmassunit} &  & \MetricUnit{} \\
% \A{???}					&			& \Z{Bohr} 			& \MetricUnit{} \\


% {electron mass}{electronmass} &  & \MetricUnit{}  \\
%  						& 			& {Elektronenvolt} & \MetricUnit{eV} \\\addlinespace
% {elementary charge}{elementarycharge} &  & \MetricUnit{}  \\
% {hartree} &  & \MetricUnit{}  \\
% {reduced Planck constant}{planckbar} &  & \MetricUnit{}  \\


% Distanz					& 			&{{\aa}ngstr{\"o}m} 	& \MetricUnit{°A} &  & \\

% {barn} &  & \\
% && {Bel}   \\

% {knot} &  & \\

% {nautical mile}{nauticalmile} \\
% {neper} &  & \\
% {{\aa}ngstr{\"o}m}{angstrom} &  & \\
\\\addlinespace\bottomrule
\end{tabulary}
}{Übersicht Gebräuchliche Einheiten}{}


% \begin{table}\FontTableBase
% \caption{Coherent derived units in the SI with special names and
% symbols.}
% \label{tab:unit:derived}
% \centering
% \begin{tabular}{lll>{\qquad}lll}
% \toprule
% \TabHeading{Unit} & \TabHeading{Macro} & \TabHeading{Symbol} & \TabHeading{Unit} & \TabHeading{Macro} & \TabHeading{Symbol}
% \\
% \midrule
% &&\MetricUnit{} 		& &&\MetricUnit{N}    \\
% Grad Celsius & Temperatur 			&\MetricUnit{\textcelsius} & Ohm 	& Elektr. Widerstand 		&\MetricUnit{$\Omega$}\\
% \Z{Coulomb}	& Elektr. Ladung		& \Z{\MetricUnit{C}}   	& Pascal 	& Druck 					&\MetricUnit{Pa} \\
% \Z{Farad} 	& \Z{Kapazität} 		& \Z{\MetricUnit{F}}    & 			&  							&\MetricUnit{} \\
% \Z{Gray} 	& \Z{Absorptionsdosis} 	& \Z{\MetricUnit{Gy}} 	& Siemens 	& Elektr. Konduktivität		&\MetricUnit{S} \\
% Hertz 		& Frequenz				& \MetricUnit{Hz}     	& Sievert 	& Dosisäquivalent			&\MetricUnit{Sv} \\
% \Z{Henry} 	& \Z{Indukktivität} 	& \Z{\MetricUnit{H}}    & 			&  							&\MetricUnit{} \\
% Joule 		& Energie, Arbeit		& \MetricUnit{J}     	& Tesla 	& magnetsiche Feldstärke 	&\MetricUnit{T} \\
% \Z{Katal} 	&  						& \Z{\MetricUnit{kat}}  & Volt 		& Spannung 					&\MetricUnit{V} \\
% \Z{Lumen} 	&  						& \Z{\MetricUnit{lm}}   & Watt 		& Leistung 					&\MetricUnit{W} \\
% \Z{Lux} 	&  						& \Z{\MetricUnit{lx}}   & \Z{Weber} & Magn. Fluß 				&\Z{\MetricUnit{}} \\
% \bottomrule
% \end{tabular}
% \end{table}


% \begin{table}\FontTableBase
% \caption{Non-SI units accepted for use with the
%    International System of Units.}
% \label{tab:unit:accepted}
% \centering
% \begin{tabular}{ll}
% \toprule
% Unit &  Symbol
% \\\midrule
% % Tag 			& {\day}\\
% Grad 			& {\textdegree}\\
% 				& {} \\
% Stunde 			& {\Hour}\\
% Liter 			& \MetricUnit{l}, {\Liter}\\
% Winkelminute 	& \\
% Minute (Zeit) 	& {\Min}\\
% Winkelsekunde 	& \\
% 				& \MetricUnit{t}
% \\\bottomrule
% \end{tabular}
% \end{table}
% 
% \begin{table}\FontTableBase
% \caption{Non-SI units whose values in SI units must be obtained
% experimentally.}
% \label{tab:unit:physical}
% \centering
% \begin{tabular}{lll}
% \toprule
% \TabHeading{Unit} & \TabHeading{Macro} & \TabHeading{Symbol} \\
% \midrule
% % {astronomical unit}{astronomicalunit} &  & \MetricUnit{} \\
% % {atomic mass unit}{atomicmassunit} &  & \MetricUnit{} \\
% {bohr} &  & \MetricUnit{} \\
% % {speed of light}{clight} &  & \MetricUnit{} \\
% Dalton &  & \MetricUnit{Da}       \\
% % {electron mass}{electronmass} &  & \MetricUnit{}  \\
% {electronvolt} &  & \MetricUnit{eV} \\
% % {elementary charge}{elementarycharge} &  & \MetricUnit{}  \\
% % {hartree} &  & \MetricUnit{}  \\
% % {reduced Planck constant}{planckbar} &  & \MetricUnit{}  \\
% \bottomrule
% \end{tabular}
% \end{table}




\Layout{57}{}{}{}{
\begin{tabular}{lcrllcrr}
\toprule
\TabHeading{Präfix} & \TabHeading{Symbol} & \TabHeading{{$10^x$}} & & \TabHeading{Präfix} & \TabHeading{Symbol} & \TabHeading{{$10^x$}} & \\
\midrule
{deci}  & d &  -1 		& 0,1 		& {deca}  & da &   1 & 10 \\
{centi} & c &  -2 		& 0,01 		& {hecto} & h &   2 & 100\\
{milli} & m &  -3 		& 0,001 	& {kilo}  & k &   3 & 1\,000\\
{micro} & \textmu & -6 	& 0,000\,1 	& {mega}  & M &   6 & 1\,000\,000\\
{nano}  & n &  -9 		& \ldots 	& {giga}  & G &   9 & \ldots\\
{pico}  & p & -12 		& & {tera}  & T &  12 & \\
{femto} & f & -15 		& & {peta}  & P &  15 & \\
{atto}  & a & -18 		& & {exa}   & E &  18 & \\
{zepto} & z & -21 		& & {zetta} & Z &  21 & \\
{yocto} & y & -24 		& & {yotta} & Y &  24 & \\
\bottomrule
\end{tabular}
}{Präfixe für Maßeinheiten\label{tab:unit:prefix}}{}








\nocite{Erc2005Sec1}
\nocite{Erc2005Sec2}
\nocite{Erc2005Sec3}
\nocite{Erc2005Sec4}
\nocite{Erc2005Sec5}
\nocite{Erc2005Sec6}
\nocite{Erc2005Sec7}
\nocite{Erc2005Sec8}
\nocite{Erc2005Sec9}

\nocite{Erc2005DeSec1}
\nocite{Erc2005DeSec2}
\nocite{Erc2005DeSec3}
\nocite{Erc2005DeSec4}
\nocite{Erc2005DeSec5}
\nocite{Erc2005DeSec6}
\nocite{Erc2005DeSec7}
\nocite{Erc2005DeSec8}
\nocite{Erc2005DeSec9}






% \begin{figure}[H]
% \begin{WidePar}
% \includegraphics[width=\linewidth]{./images/nasa-pd/114125main_image_feature_323_ys_full-00800.jpg}
% 
% \Captionof{figure}{Orbital ›Thumbs Up‹.
% The Expedition 10 and 11 crews give the >>thumbs up<< after the change of
% command ceremony on the International Space Station on April 22, 2005. From left:
% Expedition 11 Flight Engineer John L. Phillips, Expedition 11 Commander Sergei K.
% Krikalev, Expedition 10 Commander Leroy Chiao, European Space Agency (ESA)
% astronaut Roberto Vittori of Italy; and Expedition 10 Flight Engineer Salizhan S.
% Sharipov. \SourcePicture{NASA/JSC}}
% \end{WidePar}
% \end{figure}
SO LONG..........

AND THANKS FOR ALL THE FISH
\MultiColListEnd{lomassnahmen}
\MultiColListEnd{lofazit}
\MultiColListEnd{toc}
% \begingroup
% \pagestyle{empty}
% \cleardoublepage
% \endgroup
% \emph{Notizen}
% 
% \clearpage
% \thispagestyle{empty}
% \sffamily
% 
% {\marginline{\flushright{\scalefont{5}\AASS}}\sffamily\Large{\color{ubuntured}%
% \noindent Arbeits- und Ausbildungsstandards}\\ 
% {\color{darkBlue}für den
% Sanitätsdienst}}
% 
% \noindent\vspace{\baselineskip}
% 
% \DocVersion\\
% \DocVersionInfo
% 
% \parpic[r][r]{\includegraphics[width=2cm]{./images/gabriel-sebastian-aass/avatar-defense.jpg}}
% Das vorliegende Werk ist das (Zwischen-)Ergebnis einer seit mehr als zwei Jahren
% andauernden Arbeit. Oft werden wir gefragt \Speech{"`Wann wird Euer Buch fertig
% sein?"'} Nun, das wird es wohl nie sein. Die {\AASS} sind darauf ausgelegt,
% ständig weiterentwickelt, gewartet und verbessert zu werden. Das medizinische
% und organisatorische Umfeld wandelt sich ständig, sodass man es sich gar nicht
% leisten könnte, still zu stehen.
% 
% Ziel dieses Werkes ist es, das Wissen und die Welt des Sanitätsdienstes abseits
% bereits breitgetretener Pfade darzustellen, und sowohl eine Unterlage für die
% Ausbildung, als auch ein Hilfsmittel für die Praxis zu schaffen.  Besonderen
% Wert haben wir dabei nicht nur auf die Praxis, sondern auch -- heutzutage nicht
% mehr ganz modern -- auf das Hintergrundwissen und relevante Randbereiche gelegt.
% Denn in einer Zeit, in der das praktische Wissen innerhalb weniger Jahre bereits
% wieder veraltet ist, bildet dieses Fundament die Chance für eine
% langjährige Tätigkeit und erlaubt selbstständiges, regel\E{geleitetes} Arbeiten
% (anstatt strikten, regel\E{gebundenen} Handelns). 
% 
% In diesem Sinne entstanden in Zusammenarbeit mit dem Ausbildungszentrum des ASB
% Florisdorf-Donaustadt ausgehend von der Ausbildungsverordnung 
% österreichischer Rettungssanitäter und den Erfordernissen der entsprechenden
% Kurse die {\AASS}. Sie bilden die Grundlage der Curricula und konnten sich
% schon, etappenweise und unter kontrollierten Bedingungen eingeführt, seit über
% einem Jahr als Lern- und Ausbildungsunterlage bewähren.
% 
% \begin{labeling}[:]{mmmmm}\ttfamily
%   \item[APID] ---
%   \item[ISBN] ---
%   \item[DOI] ---
% \end{labeling}
\end{document}